% Orthogonalité dans l'espace — Mathématiques 1ère D — Cours signé OnBuch+
% Programme officiel MINESEC. Compiler avec : tectonic N18-orthogonalite-espace.tex
\newcommand{\DOCMATIERE}{Mathématiques}
\newcommand{\DOCNIVEAU}{1ère D}
\newcommand{\DOCMODULE}{Géométrie dans l'espace}
\newcommand{\DOCLECON}{N18}
\newcommand{\DOCTITRE}{Orthogonalité dans l'espace}
% OnBuch+ — préambule « cours signés » ENRICHI (compilé avec tectonic / XeLaTeX)
% Système visuel « Pop » d'OnBuch+ : fond crème (jamais blanc), bordures encre
% épaisses, ombres « dures » décalées, titres Archivo Black en capitales.
% Version MATHÉMATIQUES : graphes (pgfplots/tikz), boîtes pédagogiques
% (définition, propriété, méthode, attention, le savais-tu, exercice résolu,
% exercice, TP, bilan), page de garde et signature OnBuch+ ancrée en bas.
%
% Macros attendues dans main.tex AVANT \input{preamble} :
%   \DOCMATIERE  \DOCNIVEAU  \DOCMODULE  \DOCLECON (numéro)  \DOCTITRE
\documentclass[11pt]{article}

\usepackage{fontspec}
\usepackage[french]{babel}
\usepackage[a4paper,margin=2cm,top=3.4cm,bottom=2.8cm,headheight=1.4cm,footskip=1.3cm]{geometry}
\usepackage{amsmath,amssymb,amsfonts}
\usepackage{mathtools} % \xmapsto, \coloneqq... souvent utilisés par les modèles
\usepackage{cancel} % simplifications barrées
% \abs{z} (module, valeur absolue) et \norm{u} (norme), utilisés par les modèles
\providecommand{\abs}{}\let\abs\relax\DeclarePairedDelimiter{\abs}{\lvert}{\rvert}
\providecommand{\norm}{}\let\norm\relax\DeclarePairedDelimiter{\norm}{\lVert}{\rVert}
\usepackage[table]{xcolor} % [table] : fournit aussi \rowcolors (alternance de lignes)
\usepackage{pagecolor}
\usepackage{tikz}
\usetikzlibrary{arrows.meta,positioning,calc,shapes.geometric,shapes.misc,shapes.arrows,decorations.pathmorphing,decorations.pathreplacing,decorations.markings,patterns,shadows,mindmap,trees,fit,backgrounds,matrix,intersections,angles,quotes}
\usepackage{pgfplots}
\usepackage{pgf-pie} % diagrammes circulaires \pie (statistiques)
\pgfplotsset{compat=1.18}
\usepgfplotslibrary{fillbetween,groupplots} % aires entre courbes (calcul intégral)
\usepackage{enumitem}
\usepackage{fancyhdr}
% [most] charge toutes les bibliothèques tcolorbox (listings, minted, raster,
% documentation...) et gonfle inutilement la mémoire TeX sur un document long
% (~300 boîtes) : on ne charge que ce qui est réellement utilisé.
\usepackage[skins,breakable]{tcolorbox}
\usepackage{graphicx}
\usepackage{colortbl}
\usepackage{booktabs}
\usepackage{tabularx}
\usepackage{diagbox} % case d'en-tête diagonale des tableaux à double entrée
% Colonnes extensibles centrée / gauche / droite (C, L, R), souvent utilisées par les modèles
\newcolumntype{C}{>{\centering\arraybackslash}X}
\newcolumntype{L}{>{\raggedright\arraybackslash}X}
\newcolumntype{R}{>{\raggedleft\arraybackslash}X}
\usepackage{array}
\usepackage{multicol}
\usepackage{siunitx}
% parse-numbers=false : accepte aussi \SI{2,5 \times 10^{-3}}{...}
\sisetup{locale=FR,output-decimal-marker={,},group-minimum-digits=4,parse-numbers=false}
\DeclareSIUnit\ohm{\ensuremath{\symup{\Omega}}} % U+2126 absent de la police
\usepackage{qrcode}
% \degree : commande fréquemment utilisée par les modèles pour un angle ou une
% température en dehors de \SI{}{} (non fournie par les paquets chargés).
\providecommand{\degree}{^{\circ}}
% \qed (fin de démonstration), utilisé par les modèles sans amsthm
\providecommand{\qed}{\hfill$\square$}
% \barre : double barre des valeurs interdites dans un tableau de variations (tkz-tab)
\providecommand{\barre}{\ensuremath{\|}}
% environnement proof (amsthm non chargé) utilisé par les modèles
\ifdefined\proof\else\newenvironment{proof}[1][Démonstration]{\par\noindent\textit{#1.}\ }{\qed\par}\fi
\usepackage{lastpage}
\usepackage{hyperref}
\hypersetup{hidelinks}

% ============================================================
% Palette « Pop » OnBuch+ — mêmes hex que la classe P (plus_ui.dart)
% ============================================================
\definecolor{poporange}{HTML}{F59321}
\definecolor{popdark}{HTML}{E07A0C}
\definecolor{popcream}{HTML}{FFF6E8}
\definecolor{popink}{HTML}{1C1714}
\definecolor{poppurple}{HTML}{7A5AE0}
\definecolor{popgreen}{HTML}{2E9E6A}
\definecolor{poppink}{HTML}{E0457B}
\definecolor{popblue}{HTML}{2F7DE1}
\definecolor{popgold}{HTML}{C98A12}
\definecolor{popmuted}{HTML}{8A7F76}
\definecolor{popline}{HTML}{EADFD2}
\definecolor{popgray}{HTML}{8A7F76}
\colorlet{popgrey}{popgray}
% Alias tolérants (le modèle nomme parfois les couleurs différemment)
\colorlet{popwhite}{white}
\colorlet{popblack}{popink}
\colorlet{popred}{poppink}
\colorlet{popyellow}{popgold}
% Teintes claires pour fonds de tableaux / zones
\colorlet{popblueL}{popblue!10!white}
\colorlet{poporangeL}{poporange!14!white}
\colorlet{popgreenL}{popgreen!12!white}
\colorlet{poppurpleL}{poppurple!12!white}
\colorlet{poppinkL}{poppink!12!white}
\colorlet{popgoldL}{popgold!14!white}

\pagecolor{popcream}

% ============================================================
% Typographie
% ============================================================
\defaultfontfeatures{Ligatures=TeX}
\setmainfont{PlusJakartaSans-Medium.ttf}[Path=../../../fonts/,BoldFont=PlusJakartaSans-Bold.ttf]
% Maths en sans-serif (Fira Math) avec chiffres et lettres droites de Plus
% Jakarta, pour que les formules se fondent dans le texte.
\usepackage[math-style=ISO,bold-style=ISO]{unicode-math}
\setmathfont{FiraMath-Regular.otf}[Path=../../../fonts/]
\setmathfont{PlusJakartaSans-Medium.ttf}[Path=../../../fonts/,range={up/num,up/{Latin,latin},\mathup}]
\usepackage{icomma} % virgule décimale sans espace en mode math
% Caractères mathématiques Unicode absents des polices de texte (Plus Jakarta,
% Archivo Black) : rendus par leur équivalent mathématique, en texte comme en
% formule — sinon ils s'affichent en carré vide (ex. « Arithmétique dans ℤ »).
\usepackage{newunicodechar}
\newunicodechar{ℝ}{\ensuremath{\mathbb{R}}}
\newunicodechar{ℤ}{\ensuremath{\mathbb{Z}}}
\newunicodechar{ℂ}{\ensuremath{\mathbb{C}}}
\newunicodechar{ℕ}{\ensuremath{\mathbb{N}}}
\newunicodechar{ℚ}{\ensuremath{\mathbb{Q}}}
\newunicodechar{𝔻}{\ensuremath{\mathbb{D}}}
\newunicodechar{α}{\ensuremath{\alpha}}
\newunicodechar{β}{\ensuremath{\beta}}
\newunicodechar{γ}{\ensuremath{\gamma}}
\newunicodechar{δ}{\ensuremath{\delta}}
\newunicodechar{ε}{\ensuremath{\varepsilon}}
\newunicodechar{θ}{\ensuremath{\theta}}
\newunicodechar{λ}{\ensuremath{\lambda}}
\newunicodechar{μ}{\ensuremath{\mu}}
\newunicodechar{σ}{\ensuremath{\sigma}}
\newunicodechar{τ}{\ensuremath{\tau}}
\newunicodechar{φ}{\ensuremath{\varphi}}
\newunicodechar{ψ}{\ensuremath{\psi}}
\newunicodechar{ω}{\ensuremath{\omega}}
\newunicodechar{Σ}{\ensuremath{\Sigma}}
% majuscules produites par \MakeUppercase dans les titres (α-aminés -> Α)
\newunicodechar{Α}{\ensuremath{\alpha}}
\newunicodechar{Β}{\ensuremath{\beta}}
\newunicodechar{Γ}{\ensuremath{\Gamma}}
\newunicodechar{Δ}{\ensuremath{\Delta}}
\newunicodechar{Ε}{\ensuremath{\varepsilon}}
\newunicodechar{Θ}{\ensuremath{\Theta}}
\newunicodechar{Λ}{\ensuremath{\Lambda}}
\newunicodechar{Μ}{\ensuremath{\mu}}
\newunicodechar{Τ}{\ensuremath{\tau}}
\newunicodechar{Φ}{\ensuremath{\Phi}}
\newunicodechar{Ψ}{\ensuremath{\Psi}}
\newunicodechar{Ω}{\ensuremath{\Omega}}
\newunicodechar{↦}{\ensuremath{\mapsto}}
\newunicodechar{∈}{\ensuremath{\in}}
\newunicodechar{∉}{\ensuremath{\notin}}
\newunicodechar{∧}{\ensuremath{\wedge}}
\newunicodechar{⇔}{\ensuremath{\Leftrightarrow}}
\newunicodechar{⟺}{\ensuremath{\Longleftrightarrow}}
\newunicodechar{⇒}{\ensuremath{\Rightarrow}}
\newunicodechar{⟹}{\ensuremath{\Longrightarrow}}
\newunicodechar{∘}{\ensuremath{\circ}}
\newunicodechar{∪}{\ensuremath{\cup}}
\newunicodechar{∩}{\ensuremath{\cap}}
\newunicodechar{≡}{\ensuremath{\equiv}}
\newunicodechar{⊂}{\ensuremath{\subset}}
\newunicodechar{⊥}{\ensuremath{\perp}}
\newunicodechar{∃}{\ensuremath{\exists}}
\newunicodechar{∀}{\ensuremath{\forall}}
\newunicodechar{∛}{\ensuremath{\sqrt[3]{\,}}}
\newunicodechar{∜}{\ensuremath{\sqrt[4]{\,}}}
\newunicodechar{⃗}{\ensuremath{{}^{\to}}}
\newunicodechar{ⁿ}{\ifmmode^{n}\else\textsuperscript{\ensuremath{n}}\fi}
\newunicodechar{ˣ}{\ifmmode^{x}\else\textsuperscript{\ensuremath{x}}\fi}
\newunicodechar{ᵇ}{\ifmmode^{b}\else\textsuperscript{\ensuremath{b}}\fi}
\newunicodechar{ʳ}{\ifmmode^{r}\else\textsuperscript{\ensuremath{r}}\fi}
\newunicodechar{ᵘ}{\ifmmode^{u}\else\textsuperscript{\ensuremath{u}}\fi}
\newunicodechar{ʸ}{\ifmmode^{y}\else\textsuperscript{\ensuremath{y}}\fi}
\newunicodechar{ᵃ}{\ifmmode^{a}\else\textsuperscript{\ensuremath{a}}\fi}
\newunicodechar{ᵉ}{\ifmmode^{e}\else\textsuperscript{\ensuremath{e}}\fi}
\newunicodechar{ᵏ}{\ifmmode^{k}\else\textsuperscript{\ensuremath{k}}\fi}
\newunicodechar{ᵖ}{\ifmmode^{p}\else\textsuperscript{\ensuremath{p}}\fi}
\newunicodechar{ᴺ}{\ifmmode^{N}\else\textsuperscript{\ensuremath{N}}\fi}
\newunicodechar{ᶜ}{\ifmmode^{c}\else\textsuperscript{\ensuremath{c}}\fi}
\newunicodechar{ʰ}{\ifmmode^{h}\else\textsuperscript{\ensuremath{h}}\fi}
\newunicodechar{ᵠ}{\ifmmode^{q}\else\textsuperscript{\ensuremath{q}}\fi}
\newunicodechar{ₙ}{\ifmmode_{n}\else\textsubscript{\ensuremath{n}}\fi}
\newunicodechar{ₐ}{\ifmmode_{a}\else\textsubscript{\ensuremath{a}}\fi}
\newunicodechar{ᵢ}{\ifmmode_{i}\else\textsubscript{\ensuremath{i}}\fi}
\newunicodechar{ᵣ}{\ifmmode_{r}\else\textsubscript{\ensuremath{r}}\fi}
\newunicodechar{ₖ}{\ifmmode_{k}\else\textsubscript{\ensuremath{k}}\fi}
\newunicodechar{ᵦ}{\ifmmode_{\beta}\else\textsubscript{\ensuremath{\beta}}\fi}
\newunicodechar{ₓ}{\ifmmode_{x}\else\textsubscript{\ensuremath{x}}\fi}
\newfontfamily\popdisplay{ArchivoBlack-Regular.ttf}[Path=../../../fonts/]
\newfontfamily\poplogo{SpaceGrotesk-Bold.ttf}[Path=../../../fonts/]
\newfontfamily\popsemi{PlusJakartaSans-SemiBold.ttf}[Path=../../../fonts/]
\newfontfamily\popxbold{PlusJakartaSans-ExtraBold.ttf}[Path=../../../fonts/]
\newfontfamily\popxboldls{PlusJakartaSans-ExtraBold.ttf}[Path=../../../fonts/,LetterSpace=3.0]

\setlist[enumerate]{leftmargin=1.7em,itemsep=5pt,topsep=4pt}
\setlist[itemize]{leftmargin=1.7em,itemsep=4pt,topsep=4pt,label={\color{poporange}\textbullet}}
\setlength{\parindent}{0pt}
\setlength{\parskip}{5pt}
\renewcommand{\arraystretch}{1.25}

% pgfplots : style Pop par défaut
\pgfplotsset{
  every axis/.append style={axis line style={popink,line width=1pt},tick style={popink},
    grid style={popline,line width=0.5pt},label style={font=\small},tick label style={font=\footnotesize}},
  popcurve/.style={poporange,line width=1.8pt,no markers,smooth},
}
% Même style disponible dans un \draw TikZ simple (hors pgfplots)
\tikzset{popcurve/.style={poporange,line width=1.8pt}}

% ============================================================
% Pill / squiggle
% ============================================================
\newcommand{\poppill}[3]{%
  \tikz[baseline=-0.55ex]\node[draw=popink,line width=1.2pt,rounded corners=2.6pt,%
    fill=#2,inner xsep=6.5pt,inner ysep=3pt]{%
    \popxboldls\fontsize{7}{7}\selectfont\color{#3}\MakeUppercase{#1}};%
}
\newcommand{\popsquiggle}[1]{%
  \tikz[baseline=-0.4ex]{
    \draw[#1,line width=2.6pt,line cap=round]
      plot [smooth] coordinates {(0,0.09) (0.34,0.01) (0.68,0.21) (1.02,0.01) (1.36,0.10)};
  }%
}
% Mot-clé surligné (marqueur orange sous le texte)
\newcommand{\cle}[1]{{\popxbold\color{popdark}#1}}

% ============================================================
% En-tête / pied
% ============================================================
\pagestyle{fancy}
\fancyhf{}
\renewcommand{\headrulewidth}{0pt}
\renewcommand{\footrulewidth}{0pt}
\lhead{\raisebox{-0.15ex}{\poplogo\fontsize{13}{13}\selectfont\color{popink}OnBuch\color{poporange}+}}
\rhead{\poppill{\DOCMATIERE\ \textbullet\ \DOCNIVEAU\ \textbullet\ Leçon \DOCLECON}{white}{popink}}
\lfoot{\popxbold\fontsize{7.6}{7.6}\selectfont\color{popmuted}\MakeUppercase{Cours signé OnBuch+}\ \textbullet\ {\popsemi\DOCTITRE}}
\rfoot{\tikz[baseline=-0.6ex]\node[rounded corners=8.5pt,fill=popink,minimum height=17pt,minimum width=17pt,inner xsep=5pt,inner sep=0]{\popxbold\fontsize{8}{8}\selectfont\color{popcream}\thepage\,/\,\pageref*{LastPage}};}

% ============================================================
% Titres
% ============================================================
\newcommand{\chaptitre}[2]{%
  \begin{tcolorbox}[enhanced,colback=poporange,colframe=popink,boxrule=2.6pt,arc=11pt,
      drop shadow=popink,left=15pt,right=15pt,top=12pt,bottom=11pt,before skip=3pt,after skip=18pt]
    \poppill{#2}{popink}{popcream}\par\vspace{9pt}
    {\raggedright\hyphenpenalty=10000\exhyphenpenalty=10000\popdisplay\fontsize{22}{24}\selectfont\color{popcream}\MakeUppercase{#1}\par}
  \end{tcolorbox}
}
% Section numérotée : pastille orange avec le numéro + titre Archivo
\newcounter{popsec}
\newcommand{\coursec}[1]{%
  \stepcounter{popsec}\setcounter{popsub}{0}%
  \par\vspace{16pt}\noindent
  \begin{minipage}{\linewidth}
  \tikz[baseline=-0.7ex]\node[draw=popink,line width=1.6pt,fill=poporange,rounded corners=4pt,minimum size=22pt,inner sep=0]{\popdisplay\fontsize{12}{12}\selectfont\color{popcream}\thepopsec};%
  \hspace{8pt}{\popdisplay\fontsize{15}{17}\selectfont\color{popink}\MakeUppercase{#1}}\par
  \vspace{4pt}\noindent{\color{poporange}\rule{36pt}{3.6pt}}
  \end{minipage}\par\nopagebreak\vspace{8pt}
}
\newcounter{popsub}
\newcommand{\courssub}[1]{%
  \stepcounter{popsub}%
  \par\vspace{9pt}\noindent{\popxbold\fontsize{11.5}{13}\selectfont\color{popink}{\color{poporange}\thepopsec.\thepopsub}\hspace{6pt}#1}\par\nopagebreak\vspace{4pt}
}

% ============================================================
% Boîtes pédagogiques — même recette : fond blanc, bordure épaisse colorée,
% ombre dure encre, badge plein. Titre optionnel [#1].
% ============================================================
\tcbset{popbase/.style={breakable,enhanced,colback=white,boxrule=2.3pt,arc=9pt,
  shadow={3pt}{-3pt}{0pt}{popink},left=14pt,right=14pt,top=11pt,bottom=11pt,before skip=11pt,after skip=15pt,
  fontupper=\popsemi\fontsize{10.6}{14.5}\selectfont\color{popink},
  fontlower=\popsemi\fontsize{10.6}{14.5}\selectfont\color{popink}}}
% Coche dessinée (le glyphe ✓ n'existe pas dans les polices du document)
\newcommand{\popcheck}[1]{\tikz[baseline=-0.5ex]\draw[#1,line width=1.6pt,line cap=round,line join=round](-0.13,0.0)--(-0.04,-0.09)--(0.13,0.1);}
\newcommand{\popboxhead}[3]{\poppill{#1}{#2}{white}\ifx\relax#3\relax\else\hspace{6pt}{\popxbold\fontsize{10.6}{12}\selectfont\color{popink}#3}\fi\par\vspace{7pt}}

\newtcolorbox{aretenir}[1][]{popbase,colframe=popgreen,before upper={\popboxhead{À retenir}{popgreen}{#1}}}
\newtcolorbox{exemplebox}[1][]{popbase,colframe=poporange,before upper={\popboxhead{Exemple}{poporange}{#1}}}
\newtcolorbox{definition}[1][]{popbase,colframe=popblue,before upper={\popboxhead{Définition}{popblue}{#1}}}
\newtcolorbox{propriete}[1][]{popbase,colframe=poppurple,before upper={\popboxhead{Propriété}{poppurple}{#1}}}
\newtcolorbox{methode}[1][]{popbase,colframe=popgold,before upper={\popboxhead{Méthode}{popgold}{#1}}}
\newtcolorbox{attention}[1][]{popbase,colframe=poppink,before upper={\popboxhead{Attention}{poppink}{#1}}}
\newtcolorbox{experience}[1][]{popbase,colframe=popblue,colback=popblueL,before upper={\popboxhead{Expérience}{popblue}{#1}}}
% « Le savais-tu ? » : carte violette pleine (contexte, culture, Cameroun)
\newtcolorbox{savaistu}[1][]{popbase,colback=poppurple,colframe=popink,
  fontupper=\popsemi\fontsize{10.4}{14.2}\selectfont\color{white},
  before upper={\poppill{Le savais-tu\,?}{white}{poppurple}\ifx\relax#1\relax\else\hspace{6pt}{\popxbold\color{white}#1}\fi\par\vspace{7pt}}}
% Exercice résolu : énoncé en haut, \tcblower puis la solution sur fond crème
\newcounter{popexo}\newcounter{popexr}
\newtcolorbox{exoresolu}[1][]{popbase,colframe=poporange,colbacklower=popcream,
  segmentation style={popink,line width=1.2pt,dashed},
  before upper={\stepcounter{popexr}\popboxhead{Exercice résolu \thepopexr}{poporange}{#1}},
  before lower={\poppill{Solution}{popink}{popcream}\par\vspace{6pt}}}
% Exercice (non résolu en place) — la correction est dans la partie Corrigés
\newtcolorbox{exercice}[1][]{popbase,colframe=popink,
  before upper={\stepcounter{popexo}\popboxhead{Exercice \thepopexo}{popink}{#1}}}
% Corrigé
\newtcolorbox{corrige}[1][]{popbase,colframe=popgreen,colback=popgreenL,
  before upper={\popboxhead{Corrigé}{popgreen}{#1}}}
% Objectifs / prérequis / bilan
\newtcolorbox{objectifs}{popbase,colframe=popink,colback=poporangeL,
  before upper={\popboxhead{Objectifs de la leçon}{popink}{}}}
\newtcolorbox{prerequis}{popbase,colframe=popink,colback=popblueL,
  before upper={\popboxhead{Prérequis}{popblue}{}}}
\newtcolorbox{bilan}{popbase,colframe=popink,colback=white,boxrule=2.8pt,
  before upper={\popboxhead{Fiche bilan}{poporange}{L'essentiel à connaître pour l'examen}}}
% Figure encadrée avec légende
\newtcolorbox{popfigure}[1][]{enhanced,breakable=false,colback=white,colframe=popline,boxrule=1.4pt,arc=8pt,
  left=10pt,right=10pt,top=10pt,bottom=8pt,before skip=10pt,after skip=12pt,halign=center,
  lower separated=false,halign lower=center,
  fontlower=\popsemi\fontsize{9}{11.5}\selectfont\color{popmuted},
  #1}
\newcommand{\legende}[1]{\par\vspace{6pt}{\popsemi\fontsize{9}{11.5}\selectfont\color{popmuted}#1}}

% ============================================================
% Signature OnBuch+
% ============================================================
\newcommand{\poplogoinline}[1]{{\poplogo\fontsize{#1}{#1}\selectfont\color{popink}OnBuch\color{poporange}+}}
\newcommand{\signatureonbuch}{%
  \begin{tcolorbox}[enhanced,colback=popink,colframe=popink,boxrule=2.2pt,arc=10pt,
      drop shadow=poporange,left=14pt,right=14pt,top=10pt,bottom=10pt,before skip=0pt,after skip=0pt]
    \tikz[baseline=-0.7ex]\node[circle,fill=popgreen,draw=popcream,line width=1.3pt,minimum size=22pt,inner sep=0]{\popcheck{white}};%
    \hspace{9pt}\begin{minipage}[c]{0.62\linewidth}
      {\popxbold\fontsize{12}{13}\selectfont\color{popcream}Signé\ \textbullet\ {\poplogo OnBuch\color{poporange}+}}\par\vspace{2pt}
      {\raggedright\popsemi\fontsize{8}{10.5}\selectfont\color{popline}\MakeUppercase{Équipe pédagogique OnBuch+}\\\MakeUppercase{Conforme au programme officiel MINESEC}\par}
    \end{minipage}\hfill
    {\poplogo\fontsize{22}{22}\selectfont\color{popcream}OnBuch\color{poporange}+}
  \end{tcolorbox}%
}

% ============================================================
% Page de garde
% #1 titre · #2 matière · #3 niveau · #4 module · #5 n° leçon · #6 durée ·
% #7 description / objectifs (texte ou itemize)
% ============================================================
\newcommand{\pagegarde}[7]{%
  \thispagestyle{empty}
  \newgeometry{a4paper,margin=1.7cm,top=1.6cm,bottom=1.4cm}%
  \begin{tikzpicture}[remember picture,overlay]
    \fill[poporange!18] ([xshift=-3.2cm,yshift=-3.6cm]current page.north east) circle (4.6cm);
    \fill[poppurple!14] ([xshift=2.4cm,yshift=7.6cm]current page.south west) circle (3.8cm);
  \end{tikzpicture}%
  {\poplogo\fontsize{30}{30}\selectfont\color{popink}OnBuch\color{poporange}+}\hfill
  \raisebox{0.6ex}{\poppill{Cours signé \textbullet\ Premium}{popink}{popcream}}\par
  \vspace{1pt}\popsquiggle{poppurple}\par
  \vspace{8pt}
  \begin{tcolorbox}[enhanced,colback=poporange,colframe=popink,boxrule=2.8pt,arc=13pt,
      drop shadow=popink,left=18pt,right=18pt,top=12pt,bottom=13pt,after skip=10pt]
    \poppill{#2 \textbullet\ #3}{popink}{popcream}\hspace{5pt}\poppill{#4}{white}{popink}\par\vspace{8pt}
    {\popdisplay\fontsize{14}{14}\selectfont\color{popink}LEÇON #5}\par\vspace{4pt}
    {\raggedright\hyphenpenalty=10000\exhyphenpenalty=10000\popdisplay\fontsize{25}{27}\selectfont\color{popcream}\MakeUppercase{#1}\par}
    \vspace{8pt}
    {\popxbold\fontsize{9.5}{9.5}\selectfont\color{popink}\MakeUppercase{Durée conseillée : #6}}
  \end{tcolorbox}
  \begin{tcolorbox}[enhanced,colback=white,colframe=popink,boxrule=2.2pt,arc=10pt,
      drop shadow=popink,left=16pt,right=16pt,top=10pt,bottom=10pt,after skip=8pt]
    \poppill{Dans cette leçon}{poppurple}{white}\par\vspace{5pt}
    \popsemi\fontsize{9.3}{12.6}\selectfont\color{popink}#7
  \end{tcolorbox}
  \noindent\begin{minipage}[t]{0.487\linewidth}
    \begin{tcolorbox}[enhanced,colback=popgreen,colframe=popink,boxrule=2.2pt,arc=10pt,
        drop shadow=popink,left=11pt,right=11pt,top=10pt,bottom=10pt,height=3.1cm,valign=center]
      \tcbox[colback=white,colframe=popink,boxrule=1.3pt,arc=5pt,boxsep=0pt,nobeforeafter,
        left=3pt,right=3pt,top=3pt,bottom=3pt,valign=center]{\qrcode[height=1.9cm]{https://play.google.com/store/apps/details?id=com.luvvix.onbuch&hl=fr}}%
      \hspace{8pt}\begin{minipage}[c]{0.5\linewidth}
        {\popdisplay\fontsize{11}{12.5}\selectfont\color{white}\MakeUppercase{Télécharge OnBuch}}\par\vspace{4pt}
        {\raggedright\popsemi\fontsize{8.4}{11}\selectfont\color{white}Gratuit \textbullet\ Google Play\\Tuteur IA, annales, concours, résultats\par}
      \end{minipage}
    \end{tcolorbox}
  \end{minipage}\hfill
  \begin{minipage}[t]{0.487\linewidth}
    \begin{tcolorbox}[enhanced,colback=poppurple,colframe=popink,boxrule=2.2pt,arc=10pt,
        drop shadow=popink,left=11pt,right=11pt,top=10pt,bottom=10pt,height=3.1cm,valign=center]
      \tikz[baseline=-0.7ex]\node[circle,fill=white,draw=popink,line width=1.3pt,minimum size=1.6cm,inner sep=0]{\popdisplay\fontsize{24}{24}\selectfont\color{poppurple}+};%
      \hspace{8pt}\begin{minipage}[c]{0.6\linewidth}
        {\popdisplay\fontsize{11}{12.5}\selectfont\color{white}\MakeUppercase{Passe à OnBuch+}}\par\vspace{4pt}
        {\raggedright\popsemi\fontsize{8.4}{11}\selectfont\color{white}Tous les cours signés, Tuteur IA illimité, fiches PDF — onglet OnBuch+ de l'app\par}
      \end{minipage}
    \end{tcolorbox}
  \end{minipage}
  \vspace{8pt}
  \popcontenu
  \vfill
  \signatureonbuch
  \restoregeometry
  \newpage
}

% Rangée « Ce cours contient » (page de garde)
\newcommand{\poptile}[3]{% #1 couleur #2 gros texte #3 légende
  \begin{tcolorbox}[enhanced,colback=#1,colframe=popink,boxrule=2pt,arc=9pt,shadow={2.5pt}{-2.5pt}{0pt}{popink},
      left=6pt,right=6pt,top=9pt,bottom=9pt,halign=center,valign=center,height=1.75cm,width=\linewidth,nobeforeafter]
    {\popdisplay\fontsize{12.5}{13}\selectfont\color{white}\MakeUppercase{#2}}\par\vspace{3pt}
    {\centering\popsemi\fontsize{7.8}{9.5}\selectfont\color{white}#3\par}
  \end{tcolorbox}}
\newcommand{\popcontenu}{%
  \par\noindent{\popxboldls\fontsize{7.5}{7.5}\selectfont\color{popmuted}CE COURS SIGNÉ CONTIENT}\par\vspace{7pt}
  \noindent\begin{minipage}[t]{0.235\linewidth}\poptile{popblue}{Cours}{complet, illustré, pas à pas}\end{minipage}\hfill
  \begin{minipage}[t]{0.235\linewidth}\poptile{poporange}{Méthodes}{et exercices résolus}\end{minipage}\hfill
  \begin{minipage}[t]{0.235\linewidth}\poptile{poppink}{Exercices}{type examen + corrigés}\end{minipage}\hfill
  \begin{minipage}[t]{0.235\linewidth}\poptile{popgreen}{Fiche bilan}{l'essentiel à retenir}\end{minipage}\par}

% Sommaire
\newtcolorbox{sommaire}{enhanced,colback=white,colframe=popink,boxrule=2.2pt,arc=10pt,
  drop shadow=popink,left=18pt,right=18pt,top=13pt,bottom=13pt,before skip=6pt,after skip=20pt,
  before upper={\poppill{Sommaire}{poppurple}{white}\par\vspace{9pt}\popsemi\fontsize{10.6}{14.5}\selectfont\color{popink}}}

% Signature de fin de cours — poussée tout en bas de la dernière page
\newcommand{\findecours}{%
  \par\vfill\nopagebreak
  \begin{center}\popsquiggle{poporange}\end{center}\vspace{4pt}
  \signatureonbuch
}
\AtBeginDocument{\def\llbracket{\mathopen{[\mkern-3.5mu[}}\def\rrbracket{\mathclose{]\mkern-3.5mu]}}} % intervalles d'entiers (glyphe ⟦ absent de la police)
% Glyphes absents de Fira Math : redéfinis à partir de glyphes disponibles
\AtBeginDocument{%
  \def\setminus{\mathbin{\backslash}}\let\smallsetminus\setminus
  \def\vdots{\mathord{\vbox{\baselineskip3.2pt\lineskiplimit0pt\kern4pt\hbox{.}\hbox{.}\hbox{.}}}}%
  \def\ddots{\mathinner{\mkern1mu\raise6.5pt\hbox{.}\mkern2mu\raise3.6pt\hbox{.}\mkern2mu\raise0.7pt\hbox{.}\mkern1mu}}%
  \def\triangle{\mathord{\scalebox{0.9}{$\symup{\Delta}$}}}%
  \def\nabla{\mathord{\raisebox{\depth}{\scalebox{1}[-1]{$\symup{\Delta}$}}}}%
  \def\checkmark{\popcheck{popdark}}%
  \def\star{\mathbin{\ast}}\def\bigstar{\mathord{\ast}}%
  \def\diamond{\mathbin{\scalebox{0.6}{$\Diamond$}}}%
  \def\bigcup{\mathop{\mathchoice{\raisebox{-0.3em}{\scalebox{1.7}{$\cup$}}}{\scalebox{1.25}{$\cup$}}{\cup}{\cup}}\displaylimits}%
  \def\bigcap{\mathop{\mathchoice{\raisebox{-0.3em}{\scalebox{1.7}{$\cap$}}}{\scalebox{1.25}{$\cap$}}{\cap}{\cap}}\displaylimits}%
  \def\lesseqqgtr{\mathrel{\lessgtr}}\def\bigoplus{\mathop{\oplus}\displaylimits}%
}
\providecommand{\ssi}{si et seulement si} % abréviation fréquente
\AtBeginDocument{\providecommand{\wideparen}{\overparen}} % arc de cercle

%% FIGFIX-BEGIN v3 — correctifs globaux des figures (docs/onbuch-generation/tools/figfix)
\usepackage{adjustbox}\usepackage{etoolbox}
% toute figure TikZ/pgfplots plus large que la ligne est réduite à la largeur disponible
\BeforeBeginEnvironment{tikzpicture}{\begin{adjustbox}{max width=\linewidth}}
\AfterEndEnvironment{tikzpicture}{\end{adjustbox}}
% légendes pgfplots lisibles par-dessus les courbes
\pgfplotsset{every axis/.append style={legend style={fill=white,fill opacity=0.92,text opacity=1,draw=black!15,inner sep=3pt}}}
% étiquettes d'arêtes (syntaxe edge["..."]) avec fond blanc
\tikzset{every edge quotes/.append style={fill=white,inner sep=1.2pt}}
% bulles de cartes mentales : texte à la ligne dans la bulle, bulle un peu plus grande
\tikzset{every concept/.append style={text width=2.0cm,align=center,minimum size=2.3cm,inner sep=2pt}}
%% FIGFIX-END
\begin{document}

\pagegarde{Orthogonalité dans l'espace}{Mathématiques}{1ère D}{Géométrie dans l'espace}{N18}{8 h}{Cette leçon étend la notion d'orthogonalité du plan à l'espace : droites orthogonales, droite orthogonale à un plan, plans orthogonaux. Elle fournit les théorèmes fondamentaux et les méthodes de démonstration indispensables pour la suite du programme (produit scalaire, géométrie analytique) et pour les épreuves du Probatoire et du Baccalauréat.
\vspace{4pt}
\begin{multicols}{2}\small
\begin{itemize}
\item À la fin de cette leçon, je sais définir et reconnaître les positions relatives de deux droites de l'espace (coplanaires ou non).
\item À la fin de cette leçon, je sais montrer que deux droites sont coplanaires ou non coplanaires.
\item À la fin de cette leçon, je sais définir l'orthogonalité de deux droites de l'espace et la distinguer de la perpendicularité.
\item À la fin de cette leçon, je sais montrer que deux droites sont perpendiculaires ou orthogonales.
\item À la fin de cette leçon, je sais montrer qu'une droite est orthogonale à un plan en utilisant le théorème fondamental.
\item À la fin de cette leçon, je sais utiliser les propriétés reliant orthogonalité et parallélisme dans l'espace.
\end{itemize}\end{multicols}}

\begin{sommaire}
\begin{multicols}{2}
\begin{enumerate}
\item Positions relatives et coplanarité de deux droites
\item Droites orthogonales et droites perpendiculaires
\item Droite orthogonale à un plan
\item Orthogonalité et parallélisme : théorèmes de liaison
\item Plans orthogonaux
\item Méthodes de démonstration et applications
\item Activité d'intégration
\item Méthodes et exercices résolus
\item Exercices
\item Corrigés des exercices
\item Fiche bilan
\end{enumerate}
\end{multicols}
\end{sommaire}

\begin{objectifs}
\begin{itemize}
\item À la fin de cette leçon, je sais définir et reconnaître les positions relatives de deux droites de l'espace (coplanaires ou non).
\item À la fin de cette leçon, je sais montrer que deux droites sont coplanaires ou non coplanaires.
\item À la fin de cette leçon, je sais définir l'orthogonalité de deux droites de l'espace et la distinguer de la perpendicularité.
\item À la fin de cette leçon, je sais montrer que deux droites sont perpendiculaires ou orthogonales.
\item À la fin de cette leçon, je sais montrer qu'une droite est orthogonale à un plan en utilisant le théorème fondamental.
\item À la fin de cette leçon, je sais utiliser les propriétés reliant orthogonalité et parallélisme dans l'espace.
\item À la fin de cette leçon, je sais montrer que deux plans sont orthogonaux.
\item À la fin de cette leçon, je sais exploiter l'orthogonalité pour calculer des distances et résoudre des problèmes concrets de l'espace.
\end{itemize}
\end{objectifs}

\begin{prerequis}
\begin{itemize}
\item Positions relatives de droites et de plans dans l'espace, parallélisme (leçon sur le parallélisme dans l'espace)
\item Orthogonalité et perpendicularité dans le plan (classe de Seconde)
\item Vecteurs de l'espace et règles de calcul vectoriel (début de l'année de Première)
\item Propriétés des solides usuels : cube, pavé droit, pyramide, tétraèdre
\item Théorème de Pythagore et trigonométrie dans le triangle rectangle
\end{itemize}
\end{prerequis}

\newpage

\coursec{Positions relatives et coplanarité de deux droites}

\emph{Situation d'accroche.} Moussa, menuisier à Yaoundé, doit fixer une étagère murale parfaitement horizontale et vérifier que le pied de son meuble est bien perpendiculaire au sol. De son côté, un maçon de Douala contrôle que deux murs d'une maison se rencontrent à angle droit, sans rapporteur assez grand. Comment vérifier rigoureusement, dans l'espace, qu'une droite est perpendiculaire à un plan ou que deux plans sont orthogonaux ? La géométrie de l'espace, déjà entamée avec le parallélisme, fournit ici les outils de l'\cle{orthogonalité}. Cette leçon donne les définitions et théorèmes qui fondent ces vérifications.

\medskip
Mais avant de parler d'angles droits, il faut répondre à une question plus élémentaire : dans l'espace, comment deux droites peuvent-elles être placées l'une par rapport à l'autre ? Au collège, dans le plan, la réponse était simple : deux droites sont sécantes ou parallèles. Dans l'espace, une troisième possibilité apparaît, et c'est toute la richesse — et toute la difficulté — de ce chapitre.

\courssub{Les trois positions relatives de deux droites de l'espace}

\textbf{Intuition.} Prends deux crayons. Pose-les à plat sur ta table : ils se croisent, ou bien ils sont parallèles. Maintenant, tiens l'un des crayons horizontalement au-dessus de la table et l'autre posé sur la table, dans des directions différentes, sans qu'ils ne se touchent : ils ne se coupent pas, ils ne sont pas parallèles, et \emph{aucune feuille de papier} ne peut les contenir tous les deux à la fois. Cette troisième situation, impossible dans le plan, est quotidienne dans l'espace : pense à un fil électrique tendu le long d'un mur et à une latte posée au sol dans une autre direction.

\begin{definition}[Droites coplanaires et droites non coplanaires]
Deux droites de l'espace sont dites \cle{coplanaires} s'il existe un plan qui les contient toutes les deux. Dans le cas contraire, on dit qu'elles sont \cle{non coplanaires}.
\end{definition}

Dans un plan donné, deux droites sont soit sécantes (un unique point commun), soit parallèles (strictement parallèles : aucun point commun ; ou confondues). Les axiomes d'incidence de l'espace (par trois points non alignés passe un unique plan ; par une droite et un point extérieur passe un unique plan) permettent alors d'énoncer la caractérisation fondamentale.

\begin{propriete}[Caractérisation des droites coplanaires (admise)]
Deux droites de l'espace sont coplanaires si et seulement si elles sont \textbf{sécantes} ou \textbf{parallèles} (strictement parallèles ou confondues). Par conséquent, deux droites de l'espace sont dans l'une et une seule des trois situations suivantes :
\begin{itemize}
\item \textbf{sécantes} : elles ont un unique point commun ;
\item \textbf{parallèles} : elles sont dans un même plan et n'ont aucun point commun (ou sont confondues) ;
\item \textbf{non coplanaires} : elles ne se coupent pas et ne sont pas parallèles.
\end{itemize}
\end{propriete}

\begin{attention}[Le piège classique]
Dans l'espace, deux droites qui \emph{ne se coupent pas} ne sont \textbf{pas forcément parallèles} : elles peuvent être non coplanaires. Dire « les droites ne se rencontrent pas, donc elles sont parallèles » est une erreur de raisonnement qui coûte cher au Probatoire. Pour conclure au parallélisme, il faut en plus savoir que les deux droites sont \cle{coplanaires}.
\end{attention}

\begin{popfigure}
\begin{center}
\begin{tikzpicture}[scale=1.05, line join=round]
% Coordonnées du cube (standard : A(0,0,0), B(1,0,0), C(1,1,0), D(0,1,0), E(0,0,1), F(1,0,1), G(1,1,1), H(0,1,1))
% Projection cavalière : x -> (1,0), y -> (0.3,0.3), z -> (0,1) approx
% A(0,0), B(3,0), D(1,1), C(4,1), E(0,3), F(3,3), H(1,4), G(4,4)
\coordinate (A) at (0,0);
\coordinate (B) at (3,0);
\coordinate (D) at (1,1);
\coordinate (C) at (4,1);
\coordinate (E) at (0,3);
\coordinate (F) at (3,3);
\coordinate (H) at (1,4);
\coordinate (G) at (4,4);
% Arêtes cachées
\draw[dashed, popmuted] (A) -- (D); \draw[dashed, popmuted] (D) -- (C); \draw[dashed, popmuted] (E) -- (H); \draw[dashed, popmuted] (H) -- (D);
% Arêtes visibles
\draw[popdark] (A) -- (B) -- (C) -- (G) -- (H) -- (E) -- (A);
\draw[popdark] (B) -- (F); \draw[popdark] (C) -- (G);
\draw[popdark] (E) -- (F) -- (G); \draw[popdark] (F) -- (B);
% Sommets
\foreach \p/\pos in {A/below left, B/below right, C/right, D/above left, E/left, F/below right, G/right, H/above left}
  \fill[popdark] (\p) circle (1.3pt) node[\pos, font=\small] {$\p$};
% (AB) et (GH) parallèles (Bleu)
\draw[popblue, very thick] (A) -- (B);
\draw[popblue, very thick] (G) -- (H);
% (AC) et (EG) parallèles (Orange) -- CORRECTION: elles sont parallèles, pas sécantes
\draw[poporange, very thick] (A) -- (C);
\draw[poporange, very thick] (E) -- (G);
% (AB) et (CG) non coplanaires : (CG) en rose
\draw[poppink, very thick] (C) -- (G);
\end{tikzpicture}
\end{center}
\legende{Dans le cube $ABCDEFGH$ : en bleu, $(AB)$ et $(GH)$ sont parallèles ; en orange, $(AC)$ et $(EG)$ sont \textbf{parallèles} (diagonales de faces parallèles) ; en bleu et rose, $(AB)$ et $(CG)$ sont non coplanaires.}
\end{popfigure}

\courssub{Montrer que deux droites sont coplanaires}

\begin{methode}[Prouver la coplanarité de deux droites]
Pour montrer que deux droites $(d)$ et $(d')$ sont coplanaires, on peut :
\begin{enumerate}
\item \textbf{Exhiber un plan} dont on sait, par la configuration (face d'un solide, plan défini par trois points non alignés), qu'il contient les deux droites ;
\item montrer que $(d)$ et $(d')$ sont \textbf{parallèles} (par exemple par le théorème : deux droites parallèles à une même troisième sont parallèles entre elles) ;
\item montrer que $(d)$ et $(d')$ sont \textbf{sécantes} en un point, par exemple en prouvant qu'elles sont concourantes avec une troisième droite.
\end{enumerate}
Dans les cas 2 et 3, la coplanarité est automatique : deux droites parallèles (resp. sécantes) définissent toujours un plan unique.
\end{methode}

\begin{exemplebox}[Dans le cube $ABCDEFGH$]
\begin{itemize}
\item $(AB)$ et $(GH)$ : on a $(AB) \parallel (DC)$ (face $ABCD$) et $(DC) \parallel (HG)$ (face $DCGH$), donc $(AB) \parallel (GH)$ par transitivité. Elles sont coplanaires : le plan $(ABG)$ les contient toutes les deux (il contient $(AB)$ et le point $G$, donc la parallèle à $(AB)$ passant par $G$, c'est-à-dire $(GH)$).
\item $(AC)$ et $(EG)$ : $(AC)$ est une diagonale de la face $ABCD$, $(EG)$ une diagonale de la face $EFGH$. Les plans $(ABCD)$ et $(EFGH)$ sont parallèles ; le plan $(AEGC)$ les coupe selon $(AC)$ et $(EG)$, donc $(AC) \parallel (EG)$ (théorème de la coupe de deux plans parallèles par un plan). Elles sont coplanaires (plan $(AEGC)$).
\item $(AC)$ et $(BD)$, diagonales de la face $ABCD$, sont sécantes en le centre de cette face, donc coplanaires.
\end{itemize}
\end{exemplebox}

\courssub{Montrer que deux droites ne sont pas coplanaires}

C'est la situation réellement nouvelle de l'espace, et elle demande un raisonnement soigneux. L'outil roi est le \cle{raisonnement par l'absurde}.

\begin{methode}[Prouver la non-coplanarité par l'absurde]
Pour montrer que deux droites $(d)$ et $(d')$ ne sont pas coplanaires :
\begin{enumerate}
\item \textbf{Supposer} qu'il existe un plan $\mathcal{P}$ contenant $(d)$ et $(d')$ ;
\item en déduire que certains points (trois points non alignés, ou une droite et un point) déterminant un plan connu $\mathcal{Q}$ appartiennent à $\mathcal{P}$, donc que $\mathcal{P} = \mathcal{Q}$ (unicité du plan passant par trois points non alignés) ;
\item constater qu'un point de l'autre droite \emph{n'appartient pas} à $\mathcal{Q}$ : \textbf{contradiction} avec l'hypothèse ;
\item conclure : $(d)$ et $(d')$ sont non coplanaires.
\end{enumerate}
Variante équivalente : montrer directement qu'un point de l'une des droites n'appartient à \emph{aucun} plan contenant l'autre droite.
\end{methode}

\begin{attention}[Le piège des droites $(AD)$ et $(FG)$ dans un cube]
C'est l'erreur la plus fréquente au Probatoire ! Dans un cube $ABCDEFGH$ (ordre usuel : $ABCD$ base, $EFGH$ sommet, arêtes verticales $AE, BF, CG, DH$), les droites $(AD)$ et $(FG)$ \textbf{sont coplanaires}. Elles appartiennent au plan diagonal $(ADGF)$ (équation $x=z$ dans un repère adéquat). Elles ne sont pas parallèles, ne se coupent pas, mais sont dans un même plan.
\textbf{Exemple classique de droites non coplanaires :} $(AD)$ et $(CF)$, ou $(AB)$ et $(CG)$, ou $(AE)$ et $(DH)$.
\end{attention}

\begin{exemplebox}[$(AD)$ et $(CF)$ sont non coplanaires]
Dans le cube $ABCDEFGH$, montrons que les droites $(AD)$ et $(CF)$ ne sont pas coplanaires.

\emph{Supposons} qu'il existe un plan $\mathcal{P}$ contenant $(AD)$ et $(CF)$. Alors $\mathcal{P}$ contient les points $A$, $D$ et $C$ (puisque $C \in (CF)$). Or $A$, $D$, $C$ sont non alignés (sommets de la face $ABCD$), donc ils déterminent un unique plan : $\mathcal{P} = (ADC)$, qui est le plan de la face $ABCD$ (plan horizontal). Or $F \in (CF) \subset \mathcal{P}$, donc $F$ appartiendrait au plan $(ABCD)$. Mais $F$ est un sommet de la face supérieure $EFGH$, distincte de la base : $F \notin (ABCD)$. \textbf{Contradiction.}

Donc $(AD)$ et $(CF)$ sont \textbf{non coplanaires}.

Vérification rapide : $(AD)$ est dirigée selon la profondeur ($y$), $(CF)$ selon une diagonale verticale ($x$ et $z$) ; elles ne se coupent pas, ne sont pas parallèles, et aucun plan ne les contient toutes deux.
\end{exemplebox}

\begin{popfigure}
\begin{center}
\begin{tikzpicture}[scale=1.0]
% Plan P (horizontal)
\fill[popblueL, opacity=0.4] (0,0) -- (4,0) -- (4,2.5) -- (0,2.5) -- cycle;
\draw[popblue] (0,0) -- (4,0) -- (4,2.5) -- (0,2.5) -- cycle;
\node[popblue] at (0.5,2.1) {$\mathcal{P} = (ABCD)$};
% Droite d dans P : (AD)
\draw[popdark, very thick] (0.5,0.2) -- (0.5,2.3) node[above] {$(d)=(AD)$};
% Droite d' : (CF) traverse P en C mais sort
\draw[poppink, very thick] (3.5,0.2) -- (3.5,2.3) node[above] {$(d')=(CF)$};
% Point F au-dessus
\fill[poppink] (3.5,2.8) circle (2pt) node[above right] {$F \notin \mathcal{P}$};
\fill[popdark] (0.5,0.2) circle (2pt) node[below left] {$A$};
\fill[popdark] (0.5,2.3) circle (2pt) node[above left] {$D$};
\fill[popdark] (3.5,0.2) circle (2pt) node[below right] {$C$};
\end{tikzpicture}
\end{center}
\legende{$(AD) \subset \mathcal{P}=(ABCD)$ mais $(CF)$ contient $F \notin \mathcal{P}$ : les droites sont non coplanaires.}
\end{popfigure}

\courssub{Point de vue vectoriel (ouverture vers la Terminale)}

\begin{savaistu}[Un critère algébrique pour plus tard]
Soient $A, B, C, D$ quatre points de l'espace. Les droites $(AB)$ et $(CD)$ sont coplanaires si et seulement si les vecteurs $\overrightarrow{AB}$, $\overrightarrow{AC}$ et $\overrightarrow{AD}$ sont \cle{liés} (linéairement dépendants), c'est-à-dire s'il existe trois réels $\alpha, \beta, \gamma$ non tous nuls tels que
\[ \alpha\,\overrightarrow{AB} + \beta\,\overrightarrow{AC} + \gamma\,\overrightarrow{AD} = \vec{0}. \]
Ce critère deviendra un outil de calcul très efficace en Terminale, avec les coordonnées dans un repère de l'espace. En Première, on le signale simplement : retiens que la coplanarité de deux droites se traduit par une \emph{relation de dépendance} entre vecteurs directeurs et vecteur de liaison.
\end{savaistu}

\begin{exoresolu}[Coplanaires ou non dans le cube]
Soit $ABCDEFGH$ un cube. Déterminer, en justifiant, la position relative des droites : \textbf{a)} $(AB)$ et $(CG)$ ; \textbf{b)} $(HF)$ et $(DB)$.
\tcblower
\textbf{a)} $(AB)$ est une arête de la face $ABCD$, $(CG)$ est une arête verticale. Supposons qu'un plan $\mathcal{P}$ contienne $(AB)$ et $(CG)$. Alors $\mathcal{P}$ contient $A$, $B$ et $C$, trois points non alignés, donc $\mathcal{P} = (ABC)$, qui est le plan de la face $ABCD$. Mais $G \notin (ABCD)$ : contradiction. De plus $(AB)$ et $(CG)$ ne se coupent pas et ne sont pas parallèles. Donc $(AB)$ et $(CG)$ sont \textbf{non coplanaires}.

\textbf{b)} $(HF)$ et $(DB)$ sont des diagonales des faces $EFGH$ et $ABCD$. Ces deux faces sont parallèles, et le plan $(BDHF)$ (plan « vertical » contenant $B$, $D$, $H$, $F$) contient les deux droites. Donc $(HF)$ et $(DB)$ sont \textbf{coplanaires} ; comme elles sont dans des plans parallèles et incluses dans un même plan, elles sont \textbf{parallèles}.
\end{exoresolu}

\begin{exercice}[À toi de jouer]
Dans le cube $ABCDEFGH$, déterminer la position relative des droites suivantes, en justifiant soigneusement : $(EF)$ et $(BC)$ ; $(AF)$ et $(DG)$ ; $(EG)$ et $(FH)$.
\end{exercice}

\begin{corrige}[Exercice 1]
$(EF)$ et $(BC)$ : $(EF) \parallel (AB)$ (face $ABFE$) et $(BC) \parallel (AD)$ (face $ABCD$). $(AB)$ et $(AD)$ sont sécantes en $A$. Supposons un plan contenant $(EF)$ et $(BC)$ : il contiendrait $E$, $F$, $B$, $C$, donc le plan $(EFB)$ qui est le plan de la face $ABFE$ ; or $C \notin (ABFE)$ car $C$ appartient à la face opposée. Donc $(EF)$ et $(BC)$ sont \textbf{non coplanaires}.

$(AF)$ et $(DG)$ : $(AF)$ est une diagonale de la face $ABFE$, $(DG)$ une diagonale de la face $DCGH$. Ces faces sont parallèles, et le plan $(AFGD)$ contient les deux droites : elles sont \textbf{coplanaires et parallèles}.

$(EG)$ et $(FH)$ : ce sont les deux diagonales de la face $EFGH$ : elles sont \textbf{sécantes} en le centre de cette face, donc coplanaires.
\end{corrige}

\begin{aretenir}[L'essentiel de la section]
\begin{itemize}
\item Deux droites de l'espace sont \textbf{coplanaires} s'il existe un plan les contenant toutes les deux ; sinon elles sont \textbf{non coplanaires}.
\item Deux droites sont coplanaires si et seulement si elles sont \textbf{sécantes} ou \textbf{parallèles}.
\item Deux droites qui ne se coupent pas ne sont pas nécessairement parallèles : elles peuvent être non coplanaires.
\item Pour prouver la non-coplanarité, on raisonne \textbf{par l'absurde} : on suppose un plan commun, on l'identifie grâce à trois points non alignés, et on exhibe un point qui n'y appartient pas.
\item \textbf{Attention :} Dans un cube, $(AD)$ et $(FG)$ sont coplanaires (plan diagonal) ; les paires non coplanaires classiques sont $(AB)/(CG)$, $(AD)/(CF)$, $(AE)/(DH)$, etc.
\end{itemize}
\end{aretenir}

\coursec{Droites orthogonales et droites perpendiculaires}

Dans la section précédente, nous avons classé les couples de droites de l'espace selon leurs \emph{positions relatives} : droites coplanaires (sécantes ou parallèles) et droites non coplanaires. Nous savons désormais décider si deux droites se coupent, sont parallèles, ou ne sont dans aucun plan commun. Mais cette classification ne dit rien des \emph{angles}. Or, dans la vie quotidienne, l'angle entre deux directions est essentiel : le montant d'une porte doit faire un angle droit avec le linteau, les arêtes d'une caisse de savon OMO doivent se rencontrer perpendiculairement, et le poteau électrique de l'ENEOS doit être « droit » par rapport au sol.

Dans le plan, parler d'angle droit entre deux droites suppose qu'elles se coupent. Dans l'espace, deux droites qui ne se rencontrent jamais peuvent pourtant avoir des \emph{directions} qui font un angle droit : pense à l'arête horizontale du plafond de ta salle de classe et à l'arête verticale du mur d'en face. Elles ne se touchent pas, mais un menuisier dirait qu'elles sont « d'équerre ». Toute cette section consiste à donner un sens rigoureux à cette intuition.

\courssub{Une intuition pour commencer : l'angle de deux directions}

Prenons deux droites $(d)$ et $(d')$ de l'espace, a priori quelconques. Pour comparer leurs directions, l'idée naturelle est de les \emph{ramener au même point} : par un point $O$ choisi librement, menons la parallèle $(\Delta)$ à $(d)$ et la parallèle $(\Delta')$ à $(d')$. Ces deux droites $(\Delta)$ et $(\Delta')$ sont concourantes en $O$, donc coplanaires : nous sommes ramenés à une situation plane, que nous maîtrisons depuis la classe de Seconde. Nous pouvons alors mesurer l'angle qu'elles forment.

\begin{popfigure}
\begin{center}
\begin{tikzpicture}[scale=1.1, line cap=round]
  % droite d (horizontale, en bas)
  \draw[popblue, thick] (-3.2,-1.4) -- (3.2,-1.4) node[right] {$(d)$};
  % droite d' (verticale, décalée, en haut)
  \draw[poporange, thick] (1.6,0.6) -- (1.6,3.2) node[above] {$(d')$};
  % point O
  \fill[popdark] (-1.2,0.4) circle (1.6pt) node[below left] {$O$};
  % parallèle à d par O
  \draw[popblue, dashed, thick] (-3.2,0.4) -- (2.6,0.4) node[right] {$(\Delta)$};
  % parallèle à d' par O
  \draw[poporange, dashed, thick] (-1.2,-1.6) -- (-1.2,3.2) node[above] {$(\Delta')$};
  % angle droit en O
  \draw[popdark] (-1.2,0.4) ++(0.28,0) -- ++(0,0.28) -- ++(-0.28,0);
  % annotations
  \node[popmuted, align=center] at (-2.6,-2.2) {\small $(\Delta)\parallel(d)$ et $(\Delta')\parallel(d')$};
\end{tikzpicture}
\legende{Par un point $O$ quelconque, on mène les parallèles $(\Delta)$ et $(\Delta')$ aux droites $(d)$ et $(d')$. L'angle de $(d)$ et $(d')$ est, par définition, l'angle de $(\Delta)$ et $(\Delta')$. Ici cet angle est droit : $(d)$ et $(d')$ sont orthogonales, bien qu'elles ne se coupent pas.}
\end{center}
\end{popfigure}

\begin{attention}[Le point $O$ est-il important ?]
On pourrait craindre que le résultat dépende du point $O$ choisi. Il n'en est rien, et c'est une conséquence d'une propriété fondamentale : \textbf{deux angles dont les côtés sont respectivement parallèles et de même sens sont égaux}. Si l'on choisit un autre point $O'$, les nouvelles parallèles $(\Delta_1)$ et $(\Delta'_1)$ sont respectivement parallèles aux anciennes (transitivité du parallélisme), donc l'angle obtenu en $O'$ est égal à l'angle obtenu en $O$. La notion d'« angle de deux droites de l'espace » est donc bien définie, indépendamment du point choisi.
\end{attention}

\courssub{Définitions : orthogonales et perpendiculaires}

\begin{definition}[Droites orthogonales dans l'espace]
Deux droites $(d)$ et $(d')$ de l'espace sont dites \cle{orthogonales} lorsque les parallèles à ces droites menées par un point quelconque de l'espace sont perpendiculaires. On note $(d)\perp(d')$.
\end{definition}

\begin{definition}[Droites perpendiculaires]
Deux droites de l'espace sont dites \cle{perpendiculaires} lorsqu'elles sont \textbf{sécantes} et que leurs directions forment un angle droit.
\end{definition}

La distinction entre ces deux mots est subtile mais capitale au Probatoire :

\begin{aretenir}[Orthogonales ou perpendiculaires ?]
\begin{itemize}
\item Deux droites \textbf{perpendiculaires} sont toujours \textbf{orthogonales et sécantes}.
\item Deux droites \textbf{orthogonales} ne sont \textbf{pas nécessairement sécantes} : elles peuvent être non coplanaires.
\item Autrement dit : perpendiculaire $\Rightarrow$ orthogonale, mais la réciproque est fausse dans l'espace.
\end{itemize}
\end{aretenir}

\begin{attention}[Erreur fréquente n°1]
Dans le plan, « orthogonales » et « perpendiculaires » sont synonymes, car deux droites dont les directions font un angle droit se coupent forcément. \textbf{Dans l'espace, ce n'est plus vrai.} Écrire « les droites sont perpendiculaires » alors qu'elles ne se coupent pas est une faute de vocabulaire sanctionnée à l'examen. Réserve le mot \emph{perpendiculaires} aux droites qui se coupent à angle droit, et utilise \emph{orthogonales} dès que la question porte sur les directions.
\end{attention}

\courssub{Caractérisation vectorielle}

En Première, tu as appris que deux vecteurs non nuls $\vec{u}$ et $\vec{v}$ sont orthogonaux lorsque leur produit scalaire est nul : $\vec{u}\cdot\vec{v}=0$. Cette notion traduit exactement l'angle droit entre directions, ce qui donne une caractérisation très opérationnelle.

\begin{propriete}[Caractérisation vectorielle de l'orthogonalité]
Soit $(d)$ une droite de vecteur directeur $\vec{u}$ et $(d')$ une droite de vecteur directeur $\vec{u'}$. Alors :
\[ (d)\perp(d') \iff \vec{u}\perp\vec{u'} \iff \vec{u}\cdot\vec{u'}=0. \]
\end{propriete}

\textbf{Justification.} Mener par un point $O$ les parallèles à $(d)$ et $(d')$ revient à représenter les vecteurs $\vec{u}$ et $\vec{u'}$ à partir de $O$ : dire que ces parallèles sont perpendiculaires, c'est dire que $\vec{u}$ et $\vec{u'}$ sont orthogonaux. $\blacksquare$

Cette caractérisation est précieuse : elle transforme un problème de géométrie dans l'espace en un \emph{calcul} de produit scalaire, souvent mené dans un repère ou par décomposition de vecteurs sur les arêtes d'un solide.

\begin{exemplebox}[Lecture vectorielle dans un repère]
L'espace est muni d'un repère $(O\,;\vec{i},\vec{j},\vec{k})$. Soit $(d)$ dirigée par $\vec{u}\begin{pmatrix} 2 \\ -1 \\ 3 \end{pmatrix}$ et $(d')$ dirigée par $\vec{u'}\begin{pmatrix} 1 \\ 5 \\ 1 \end{pmatrix}$. On calcule :
\[ \vec{u}\cdot\vec{u'} = 2\times 1 + (-1)\times 5 + 3\times 1 = 2-5+3 = 0. \]
Donc $(d)$ et $(d')$ sont orthogonales — sans que nous ayons eu besoin de savoir si elles se coupent.
\end{exemplebox}

\courssub{Orthogonalité et parallélisme : les théorèmes de liaison}

Les deux propriétés qui suivent sont les outils de base de toutes les démonstrations d'orthogonalité dans l'espace. Elles disent que l'orthogonalité « voyage » le long des droites parallèles.

\begin{propriete}[Conservation par parallélisme — première forme]
Si deux droites $(d)$ et $(d')$ sont parallèles, alors toute droite orthogonale à l'une est orthogonale à l'autre.
\[ \begin{cases} (d)\parallel(d') \\ (\delta)\perp(d) \end{cases} \implies (\delta)\perp(d'). \]
\end{propriete}

\begin{experience}[Démonstration guidée]
Nous démontrons cette propriété à partir de la définition, étape par étape.
\begin{enumerate}
\item \textbf{Hypothèses.} $(d)\parallel(d')$ et $(\delta)\perp(d)$. Que signifie $(\delta)\perp(d)$ ? Par définition, les parallèles à $(\delta)$ et à $(d)$ menées par un point quelconque sont perpendiculaires.
\item \textbf{Choisissons un point $O$ quelconque} de l'espace. Menons par $O$ la parallèle $(\Delta)$ à $(\delta)$ et la parallèle $(D)$ à $(d)$. D'après l'étape 1, $(\Delta)\perp(D)$ : elles sont sécantes en $O$ et forment un angle droit.
\item \textbf{Faisons intervenir $(d')$.} Menons par $O$ la parallèle $(D')$ à $(d')$. Or $(d)\parallel(d')$, donc par transitivité du parallélisme, $(D)$ et $(D')$ sont toutes deux parallèles à $(d)$ et passent par $O$ : d'après l'unicité de la parallèle menée par un point à une droite donnée, $(D)=(D')$.
\item \textbf{Conclusion.} La parallèle à $(d')$ menée par $O$ est donc $(D)$, et $(\Delta)\perp(D)$. Ainsi les parallèles à $(\delta)$ et à $(d')$ menées par un point quelconque sont perpendiculaires : par définition, $(\delta)\perp(d')$. $\blacksquare$
\end{enumerate}
\end{experience}

\begin{propriete}[Conservation par parallélisme — deuxième forme]
Si deux droites $(d)$ et $(d')$ sont orthogonales, alors toute droite parallèle à l'une est orthogonale à l'autre.
\[ \begin{cases} (d)\perp(d') \\ (\delta)\parallel(d) \end{cases} \implies (\delta)\perp(d'). \]
\end{propriete}

\textbf{Démonstration.} Soit $O$ un point quelconque. La parallèle à $(\delta)$ menée par $O$ est aussi parallèle à $(d)$ (transitivité), donc c'est la parallèle à $(d)$ menée par $O$. Comme $(d)\perp(d')$, cette droite est perpendiculaire à la parallèle à $(d')$ menée par $O$. Donc $(\delta)\perp(d')$. $\blacksquare$

\begin{aretenir}[À retenir absolument]
L'orthogonalité est une propriété des \cle{directions}, pas des positions. Remplacer une droite par n'importe laquelle de ses parallèles ne change rien à l'orthogonalité. C'est exactement ce qui autorise la stratégie de démonstration suivante.
\end{aretenir}

\courssub{Méthode de démonstration}

\begin{methode}[Montrer que deux droites non sécantes sont orthogonales]
Pour prouver que $(d)$ et $(d')$, non sécantes, sont orthogonales :
\begin{enumerate}
\item \textbf{Repérer un plan commode} contenant l'une des droites, souvent une face d'un solide.
\item \textbf{Mener (ou repérer) par un point de l'une la parallèle à l'autre} : dans un pavé ou un cube, c'est souvent une arête parallèle déjà tracée.
\item \textbf{Se ramener au plan} : les deux droites obtenues sont concourantes, donc coplanaires.
\item \textbf{Prouver la perpendicularité dans ce plan} avec les outils de géométrie plane (angles du carré, du rectangle, Pythagore, etc.).
\item \textbf{Conclure} par la définition ou par les propriétés de conservation.
\end{enumerate}
\end{methode}

\begin{exoresolu}[Dans le cube $ABCDEFGH$]
On considère le cube $ABCDEFGH$ (les sommets $ABCD$ forment la face inférieure et $EFGH$ la face supérieure, $E$ au-dessus de $A$). Montrer que les droites $(AB)$ et $(FG)$ sont orthogonales, bien qu'elles ne soient pas sécantes.
\tcblower
\textbf{Solution détaillée.}

\emph{Étape 1 : les droites ne sont pas sécantes.} $(AB)$ est une arête de la face inférieure, $(FG)$ une arête de la face supérieure ; elles ne sont pas coplanaires (tout plan contenant $(AB)$ et un point de $(FG)$ hors de la face du dessus ne contient pas $(FG)$ entière ; plus simplement, $(AB)$ et $(FG)$ n'ont aucun point commun et ne sont pas parallèles, donc elles sont non coplanaires).

\emph{Étape 2 : on ramène au même point.} Par le point $F$, qui appartient à $(FG)$, la droite $(FE)$ est parallèle à $(AB)$. En effet, dans le cube, $(AB)\parallel(EF)$ car $ABFE$ est une face, donc un carré, et les côtés opposés d'un carré sont parallèles.

\emph{Étape 3 : perpendicularité dans un plan.} Les droites $(FE)$ et $(FG)$ sont deux arêtes de la face supérieure $EFGH$ issues du sommet $F$. Or $EFGH$ est un carré, donc ses côtés adjacents sont perpendiculaires : $(FE)\perp(FG)$ au point $F$.

\emph{Étape 4 : conclusion.} La parallèle à $(AB)$ menée par $F$ est perpendiculaire à $(FG)$. Par définition, $(AB)$ et $(FG)$ sont \textbf{orthogonales}, alors qu'elles ne se coupent pas. $\blacksquare$
\end{exoresolu}

\courssub{Le pavé droit, terrain de jeu de l'orthogonalité}

Le pavé droit (boîte de craie, brique de lait, carton de savon) est le solide roi de cette leçon : ses douze arêtes se répartissent en trois directions deux à deux orthogonales.

\begin{popfigure}
\begin{center}
\begin{tikzpicture}[scale=1.15, line join=round, line cap=round]
  % Pavé ABCDEFGH : ABCD face avant, EFGH face arrière
  \coordinate (A) at (0,0);
  \coordinate (B) at (4,0);
  \coordinate (C) at (4,2.4);
  \coordinate (D) at (0,2.4);
  \coordinate (E) at (1.3,1.1);
  \coordinate (F) at (5.3,1.1);
  \coordinate (G) at (5.3,3.5);
  \coordinate (H) at (1.3,3.5);
  % arêtes cachées en pointillés
  \draw[popmuted, dashed] (A) -- (E);
  \draw[popmuted, dashed] (E) -- (F);
  \draw[popmuted, dashed] (E) -- (H);
  % arêtes visibles
  \draw[popdark, thick] (A) -- (B) -- (C) -- (D) -- cycle;
  \draw[popdark, thick] (B) -- (F);
  \draw[popdark, thick] (C) -- (G);
  \draw[popdark, thick] (D) -- (H);
  \draw[popdark, thick] (F) -- (G) -- (H);
  % sommets
  \foreach \p/\pos in {A/below left, B/below right, C/right, D/left, E/above left, F/below right, G/above right, H/above left}
    \fill[popink] (\p) circle (1.4pt) node[\pos] {\small $\p$};
  % angle droit en B entre (BC) et (BF)
  \draw[poporange, thick] (B) ++(0,0.35) -- ++(0.2,0.175) -- ++(0,-0.35);
  % angle droit en C entre (CB) et (CG)
  \draw[popblue, thick] (C) ++(-0.35,0) -- ++(0,0.2) -- ++(0.26,0.22);
  % angle droit en G entre (GH) et (GC)
  \draw[popgreen, thick] (G) ++(-0.3,0) -- ++(0,-0.3) -- ++(0.3,0);
\end{tikzpicture}
\legende{Pavé droit $ABCDEFGH$. En chaque sommet, les trois arêtes sont deux à deux perpendiculaires (angles droits marqués en $B$, $C$ et $G$). De plus, deux arêtes de directions différentes qui ne se coupent pas — par exemple $(AB)$ et $(CG)$ — sont orthogonales sans être sécantes.}
\end{center}
\end{popfigure}

\begin{exemplebox}[Chasse aux orthogonalités dans le pavé]
Dans le pavé droit $ABCDEFGH$ de la figure :
\begin{itemize}
\item $(AB)$ et $(BC)$ sont \textbf{perpendiculaires} : sécantes en $B$ à angle droit.
\item $(AB)$ et $(CG)$ sont \textbf{orthogonales non sécantes} : en effet $(CG)\parallel(BF)$ et $(AB)\perp(BF)$, donc par conservation $(AB)\perp(CG)$.
\item $(AD)$ et $(HG)$ sont \textbf{orthogonales non sécantes} : $(HG)\parallel(DC)$ et $(AD)\perp(DC)$ dans la face $ABCD$.
\end{itemize}
\end{exemplebox}

\begin{attention}[Erreur fréquente n°2 : confondre orthogonalité et perpendicularité dans la rédaction]
Dire « $(AB)$ et $(CG)$ sont perpendiculaires » est \textbf{faux} : elles ne se coupent pas. La rédaction correcte est : « $(AB)$ et $(CG)$ sont orthogonales ». À l'inverse, si tu as montré que deux droites sont orthogonales \emph{et} qu'elles ont un point commun, tu peux conclure qu'elles sont perpendiculaires. Vérifie toujours la question de l'intersection avant de choisir le mot.
\end{attention}

\begin{attention}[Erreur fréquente n°3 : croire que deux droites orthogonales à une même troisième sont parallèles]
Dans le plan, deux droites perpendiculaires à une même droite sont parallèles. Dans l'espace, c'est \textbf{faux} : dans le pavé, $(AB)$ et $(AD)$ sont toutes deux orthogonales à $(AE)$, pourtant elles sont perpendiculaires entre elles, pas parallèles ! Retenons : dans l'espace, l'ensemble des droites orthogonales à une droite donnée en un point forme tout un \emph{plan} (ce sera l'objet de la section suivante).
\end{attention}

\begin{savaistu}[L'équerrage des charpentiers]
Les charpentiers traditionnels, de Foumban à Maroua, vérifient qu'une charpente ou un cadre de porte est bien « d'équerre » en comparant les \textbf{deux diagonales} du cadre : si les diagonales ont la même longueur, le quadrilatère est un rectangle, donc ses côtés sont perpendiculaires. C'est une application directe de la géométrie plane... mais pour assembler des fermes de toit dans l'espace, le charpentier raisonne en fait avec des directions orthogonales non sécantes : les pannes horizontales du toit doivent être orthogonales aux arêtiers, sans jamais se rencontrer. Les mathématiques de cette leçon sont littéralement au-dessus de ta tête chaque nuit !
\end{savaistu}

\begin{exercice}[Pour t'entraîner]
Dans le cube $ABCDEFGH$ (face inférieure $ABCD$, face supérieure $EFGH$, $E$ au-dessus de $A$) :
\begin{enumerate}
\item Montrer que $(BC)$ et $(EF)$ sont orthogonales. Sont-elles perpendiculaires ?
\item Montrer que les diagonales $(EG)$ et $(FH)$ de la face supérieure sont orthogonales.
\item Citer trois couples d'arêtes orthogonales non sécantes.
\end{enumerate}
\end{exercice}

\begin{corrige}[Exercice]
\begin{enumerate}
\item \textbf{Orthogonalité.} Dans la face latérale $BCGF$, $(BC)\parallel(FG)$ (carré). Dans la face supérieure $EFGH$, $(EF)\perp(FG)$ (carré). Par conservation du parallélisme (première forme), $(BC)\perp(EF)$.
\textbf{Perpendicularité ?} $(BC)$ et $(EF)$ n'ont pas de point commun (l'une est sur la face avant, l'autre sur la face supérieure, elles ne sont pas parallèles) : elles sont non coplanaires. Donc elles sont \textbf{orthogonales mais non perpendiculaires}.
\item Dans le carré $EFGH$, les diagonales $(EG)$ et $(FH)$ sont perpendiculaires (propriété du carré). Comme elles se coupent au centre du carré, elles sont \textbf{perpendiculaires}, donc a fortiori \textbf{orthogonales}.
\item Exemples (chaque couple correspond à deux directions distinctes parmi les trois directions du cube) : $(AB)$ et $(CG)$ ; $(AD)$ et $(BF)$ ; $(AE)$ et $(FG)$ ; $(BC)$ et $(DH)$ ; $(AB)$ et $(DH)$ ; $(AD)$ et $(EH)$ ; etc.
\end{enumerate}
\textbf{Moralité} : dans un cube, deux directions d'arêtes distinctes sont toujours orthogonales. Les diagonales d'une face carrée sont orthogonales entre elles. En revanche, une arête et une diagonale d'une face non adjacente ne sont généralement pas orthogonales (angle de $45^\circ$ ou $\arctan(\sqrt{2})$ pour la grande diagonale). Vérifie toujours par le raisonnement géométrique ou le calcul vectoriel !
\end{corrige}

\begin{aretenir}[Bilan de la section]
\begin{itemize}
\item $(d)\perp(d')$ signifie : les parallèles menées par un point quelconque sont perpendiculaires — notion indépendante du point choisi.
\item Perpendiculaires $=$ orthogonales $+$ sécantes.
\item Critère vectoriel : $(d)\perp(d') \iff \vec{u}\cdot\vec{u'}=0$.
\item L'orthogonalité se conserve par parallélisme : on peut toujours remplacer une droite par une parallèle pour se ramener à une configuration plane.
\end{itemize}
La section suivante exploitera ces outils pour définir la notion capitale de \emph{droite orthogonale à un plan}.
\end{aretenir}

\coursec{Droite orthogonale à un plan}

Dans la section précédente, nous avons appris à reconnaître deux droites orthogonales ou perpendiculaires dans l'espace. Mais regarde autour de toi : le poteau électrique planté dans la cour de l'école, le mât du drapeau camerounais devant la mairie, le pied d'un lampadaire à Yaoundé... Tous ces objets dressés verticalement semblent « perpendiculaires au sol tout entier », et pas seulement à une ligne tracée sur le sol. Comment traduire mathématiquement cette idée ? C'est exactement l'objet de cette section : définir ce que signifie « une droite orthogonale à un plan », établir le critère fondamental qui permet de le démontrer, et en tirer des conséquences capitales : le projeté orthogonal, la distance d'un point à un plan et le plan médiateur d'un segment.

\courssub{Définition et premières propriétés}

\textbf{Intuition.} Imagine le mât d'un drapeau planté dans une cour plane. Quelle que soit la direction que tu traces au sol en partant du pied du mât — vers le portail, vers la cantine, vers n'importe quel point de la cour — le mât fait un angle droit avec cette direction. C'est cette idée de « perpendiculaire à \emph{toutes} les directions du sol » qui fonde la définition.

\begin{definition}[Droite orthogonale à un plan]
Soit $(d)$ une droite et $(P)$ un plan de l'espace. On dit que la droite $(d)$ est \cle{orthogonale} au plan $(P)$, et on note $(d) \perp (P)$, lorsque $(d)$ est orthogonale à \textbf{toute} droite du plan $(P)$.
\end{definition}

\begin{attention}[Le mot « toute » est essentiel]
La définition exige l'orthogonalité avec \emph{toutes} les droites du plan, sans exception. Il ne suffit pas de vérifier l'orthogonalité avec une seule droite, ni même avec deux droites \emph{parallèles} du plan. Nous verrons plus loin le critère pratique qui rend cette vérification possible en un nombre fini d'étapes.
\end{attention}

\begin{propriete}[Conséquence immédiate : existence du pied]
Si une droite $(d)$ est orthogonale à un plan $(P)$, alors $(d)$ coupe $(P)$ en un \textbf{unique} point $H$, appelé \cle{pied de la perpendiculaire} (ou pied de la droite orthogonale).
\end{propriete}

\textbf{Justification.} Une droite et un plan dans l'espace sont dans l'une des trois situations suivantes : la droite est strictement parallèle au plan, elle est incluse dans le plan, ou elle le coupe en un unique point. Si $(d)$ était parallèle à $(P)$ ou incluse dans $(P)$, le plan $(P)$ contiendrait une droite $(\Delta)$ parallèle à $(d)$ ; or $(d)$ ne peut être orthogonale à une droite qui lui est parallèle (l'angle entre deux droites parallèles est nul, pas droit). Donc $(d)$ est sécante à $(P)$, et l'intersection d'une droite et d'un plan sécants est un point unique.

\begin{definition}[Projeté orthogonal d'un point sur un plan]
Soit $A$ un point et $(P)$ un plan. Le \cle{projeté orthogonal} de $A$ sur $(P)$ est le point d'intersection $H$ de $(P)$ avec la droite passant par $A$ et orthogonale à $(P)$. On dit aussi que $H$ est le pied de la perpendiculaire abaissée de $A$ sur $(P)$.
\end{definition}

\begin{popfigure}
\begin{tikzpicture}[scale=1.1]
  % Plan P en perspective
  \fill[popblueL] (-2.6,-0.6) -- (1.4,-0.6) -- (2.8,1.2) -- (-1.2,1.2) -- cycle;
  \draw[popblue] (-2.6,-0.6) -- (1.4,-0.6) -- (2.8,1.2) -- (-1.2,1.2) -- cycle;
  \node[popblue] at (-2.1,0.9) {$(P)$};
  % Point H
  \coordinate (H) at (0.3,0.15);
  % Droite (d) orthogonale
  \draw[poporange,very thick] (0.3,-0.45) -- (0.3,2.6);
  \node[poporange] at (0.75,2.5) {$(d)$};
  \fill[popink] (H) circle (1.6pt);
  \node[below right, popink] at (H) {$H$};
  % Deux droites sécantes de (P) passant par H
  \draw[popdark,thick] (-1.5,-0.25) -- (2.0,0.55);
  \node[popdark] at (2.15,0.6) {$(\Delta_1)$};
  \draw[popdark,thick] (-0.9,1.05) -- (1.5,-0.55);
  \node[popdark] at (1.65,-0.62) {$(\Delta_2)$};
  % Symboles angles droits
  \draw[popgreen,thick] (0.3,0.55) -- (0.62,0.60) -- (0.62,0.20);
  \draw[popgreen,thick] (0.3,0.55) -- (0.12,0.83) -- (-0.06,0.62);
\end{tikzpicture}
\legende{La droite $(d)$ est orthogonale au plan $(P)$ en $H$ : elle est orthogonale à toute droite de $(P)$, en particulier aux deux droites sécantes $(\Delta_1)$ et $(\Delta_2)$.}
\end{popfigure}

\courssub{Le théorème fondamental}

Vérifier l'orthogonalité avec \emph{toutes} les droites d'un plan semble impossible : il y en a une infinité ! Le théorème suivant, le plus important de cette leçon, ramène cette vérification à \textbf{deux} droites seulement, à condition qu'elles soient sécantes.

\begin{propriete}[Théorème fondamental de l'orthogonalité droite-plan]
Si une droite $(d)$ est orthogonale à \textbf{deux droites sécantes} d'un plan $(P)$, alors $(d)$ est orthogonale au plan $(P)$. \emph{(Démonstration exigible.)}
\end{propriete}

\begin{experience}[Démonstration du théorème fondamental]
\textbf{Hypothèses.} Soit $(P)$ un plan, $(\Delta_1)$ et $(\Delta_2)$ deux droites de $(P)$ sécantes en un point $O$, et $(d)$ une droite orthogonale à $(\Delta_1)$ et à $(\Delta_2)$. Quitte à translater $(d)$ parallèlement à elle-même (l'orthogonalité se conserve par parallélisme), on suppose que $(d)$ passe par $O$.

\textbf{Étape 1 : mise en place.} Sur $(d)$, de part et d'autre de $O$, plaçons deux points $M$ et $M'$ tels que $OM = OM'$. Sur $(\Delta_1)$, plaçons un point $U$ distinct de $O$, et sur $(\Delta_2)$ un point $V$ distinct de $O$.

\textbf{Étape 2 : $(\Delta_1)$ est médiatrice de $[MM']$ dans un plan auxiliaire.} Les droites $(d)$ et $(\Delta_1)$ sont orthogonales et sécantes en $O$, donc coplanaires et perpendiculaires. Dans le plan qu'elles déterminent, $(\Delta_1)$ est perpendiculaire au segment $[MM']$ en son milieu $O$ : c'est la médiatrice de $[MM']$. Tout point de $(\Delta_1)$ est donc équidistant de $M$ et $M'$ ; en particulier :
\[ UM = UM'. \]
Le même raisonnement avec $(\Delta_2)$ donne :
\[ VM = VM'. \]

\textbf{Étape 3 : la droite $(UV)$ est aussi médiatrice de $[MM']$.} Considérons les triangles $UMV$ et $UM'V$. Ils ont : $UM = UM'$ (étape 2), $VM = VM'$ (étape 2), et le côté $[UV]$ en commun. Ils sont donc isométriques (trois côtés égaux), d'où $\widehat{MUV} = \widehat{M'UV}$. Soit $N$ un point quelconque de la droite $(UV)$. Les triangles $UMN$ et $UM'N$ ont $UM = UM'$, le côté $[UN]$ commun et des angles égaux $\widehat{MUN} = \widehat{M'UN}$ : ils sont isométriques, donc $NM = NM'$.

\textbf{Étape 4 : conclusion.} Ainsi tout point $N$ de $(UV)$ est équidistant de $M$ et $M'$ : dans le plan déterminé par $(d)$ et $(UV)$, la droite $(UV)$ est la médiatrice de $[MM']$, donc $(UV) \perp (d)$ en $O$. Or toute droite de $(P)$ passant par $O$ peut jouer le rôle de $(UV)$ (il suffit de choisir $U$ et $V$ de sorte que $(UV)$ lui soit parallèle, l'orthogonalité se transmettant par parallélisme). Donc $(d)$ est orthogonale à toute droite de $(P)$ passant par $O$, puis à toute droite de $(P)$ par parallélisme. \hfill $\blacksquare$
\end{experience}

\begin{aretenir}[Le critère à retenir]
Pour prouver qu'une droite est orthogonale à un plan, il suffit de trouver \textbf{deux droites sécantes} de ce plan qui lui sont orthogonales. C'est le critère que tu utiliseras dans presque tous les exercices du Probatoire.
\end{aretenir}

\begin{propriete}[Théorème réciproque]
Si une droite $(d)$ est orthogonale à un plan $(P)$, alors $(d)$ est orthogonale à \textbf{toute} droite du plan $(P)$, même à celles qui ne passent pas par le pied $H$.
\end{propriete}

\textbf{Justification.} C'est une conséquence directe de la définition : l'orthogonalité d'une droite à un plan signifie précisément l'orthogonalité à toutes les droites du plan. Rappelons que deux droites orthogonales de l'espace ne sont pas nécessairement sécantes : si $(\Delta)$ est une droite de $(P)$ ne passant pas par $H$, la parallèle à $(\Delta)$ passant par $H$ est perpendiculaire à $(d)$, donc $(d)$ et $(\Delta)$ sont orthogonales.

\begin{methode}[Montrer qu'une droite est orthogonale à un plan]
\begin{enumerate}
  \item Identifier dans l'énoncé le plan $(P)$ et la droite $(d)$ concernés.
  \item Chercher dans $(P)$ deux droites \textbf{sécantes} $(\Delta_1)$ et $(\Delta_2)$ (souvent deux arêtes ou deux diagonales d'une face d'un solide).
  \item Démontrer séparément que $(d) \perp (\Delta_1)$ et que $(d) \perp (\Delta_2)$ (en utilisant les angles droits de la figure, les propriétés des faces carrées ou rectangulaires, etc.).
  \item Vérifier et \textbf{écrire explicitement} que $(\Delta_1)$ et $(\Delta_2)$ sont sécantes et contenues dans $(P)$.
  \item Conclure par le théorème fondamental : $(d) \perp (P)$.
\end{enumerate}
\end{methode}

\begin{attention}[Erreur classique n°1 : une seule droite ne suffit pas]
Être orthogonale à \textbf{une seule} droite d'un plan n'entraîne pas l'orthogonalité au plan. Contre-exemple : dans une salle de classe, une règle posée en biais sur le mur peut être perpendiculaire à l'arête verticale du mur sans être orthogonale au mur. De même, être orthogonale à \textbf{deux droites parallèles} du plan ne suffit pas : il faut impérativement deux droites \textbf{sécantes}. Au Probatoire, oublier de justifier le caractère sécant coûte des points.
\end{attention}

\begin{exemplebox}[Dans le cube $ABCDEFGH$]
Soit $ABCDEFGH$ un cube, avec la face $ABCD$ au sol et les arêtes verticales $[AE]$, $[BF]$, $[CG]$, $[DH]$. Montrons que $(AE) \perp (ABC)$.
\begin{itemize}
  \item La face $ABFE$ est un carré, donc $(AE) \perp (AB)$.
  \item La face $ADHE$ est un carré, donc $(AE) \perp (AD)$.
  \item Les droites $(AB)$ et $(AD)$ sont deux droites du plan $(ABC)$, sécantes en $A$.
\end{itemize}
D'après le théorème fondamental, $(AE)$ est orthogonale au plan $(ABC)$. Le pied de cette perpendiculaire est $A$.
\end{exemplebox}

\begin{exoresolu}[Pyramide à base carrée]
Soit $SABCD$ une pyramide de sommet $S$, dont la base $ABCD$ est un carré de centre $O$, et telle que la droite $(SO)$ soit orthogonale aux deux diagonales $(AC)$ et $(BD)$. Démontrer que $(SO)$ est orthogonale au plan $(ABC)$.
\tcblower
Les diagonales $(AC)$ et $(BD)$ du carré $ABCD$ sont deux droites du plan $(ABC)$, sécantes en $O$ (centre du carré). Par hypothèse, $(SO) \perp (AC)$ et $(SO) \perp (BD)$. La droite $(SO)$ est donc orthogonale à deux droites sécantes du plan $(ABC)$ : d'après le théorème fondamental, $(SO) \perp (ABC)$. On dit que $SABCD$ est une pyramide \emph{régulière} de hauteur $[SO]$.
\end{exoresolu}

\courssub{Existence et unicité de la perpendiculaire}

\begin{propriete}[Propriétés d'unicité (admises)]
On admet les deux résultats suivants :
\begin{enumerate}
  \item Par un point donné $A$, il passe une \textbf{unique} droite orthogonale à un plan donné $(P)$.
  \item Par un point donné $A$, il passe un \textbf{unique} plan orthogonal à une droite donnée $(d)$.
\end{enumerate}
\end{propriete}

\textbf{Justification intuitive.} Pose un livre bien à plat sur la table : il matérialise un plan. Par un point de la couverture, tu ne peux dresser qu'\emph{un seul} crayon « bien droit » par rapport au livre ; dès que tu l'inclines, il n'est plus perpendiculaire à toutes les directions de la couverture. Réciproquement, autour d'un crayon fixé, tu ne peux glisser qu'\emph{une seule} feuille de papier passant par un point donné et restant perpendiculaire au crayon : toute autre inclinaison rompt l'orthogonalité. Ces expériences concrètes rendent l'unicité très naturelle, même si sa démonstration rigoureuse dépasse le programme de Première.

\begin{aretenir}[Ce que l'unicité permet de faire]
Grâce à l'unicité, dès que tu as exhibé \emph{une} droite passant par $A$ et orthogonale à $(P)$, tu sais que c'est \emph{la} perpendiculaire cherchée : il est inutile d'en chercher une autre. C'est ce qui légitime la définition du projeté orthogonal et de la distance d'un point à un plan.
\end{aretenir}

\courssub{Distance d'un point à un plan}

\textbf{Intuition.} Un technicien de CAMTEL veut mesurer la hauteur d'un pylône au-dessus du sol : il ne mesure pas en diagonale, il mesure \emph{verticalement}, car c'est le chemin le plus court. De même, la distance d'un point à un plan se définit comme la plus courte distance possible entre ce point et les points du plan.

\begin{definition}[Distance d'un point à un plan]
Soit $A$ un point et $(P)$ un plan, $H$ le projeté orthogonal de $A$ sur $(P)$. La \cle{distance} du point $A$ au plan $(P)$, notée $d(A, (P))$, est la longueur $AH$ :
\[ d(A, (P)) = AH. \]
\end{definition}

\begin{propriete}[La perpendiculaire donne la plus courte distance]
Pour tout point $M$ du plan $(P)$ distinct de $H$, on a :
\[ AH < AM. \]
Autrement dit, le projeté orthogonal $H$ est le point de $(P)$ le plus proche de $A$. \emph{(Démonstration exigible.)}
\end{propriete}

\begin{experience}[Démonstration par le théorème de Pythagore]
Soit $M$ un point de $(P)$ distinct de $H$. La droite $(AH)$ est orthogonale au plan $(P)$, donc orthogonale à toute droite de $(P)$, en particulier à la droite $(HM)$. Comme $(AH)$ et $(HM)$ sont sécantes en $H$, elles sont perpendiculaires : le triangle $AHM$ est \textbf{rectangle en} $H$. D'après le théorème de Pythagore :
\[ AM^2 = AH^2 + HM^2. \]
Comme $M \neq H$, on a $HM > 0$, donc $HM^2 > 0$, d'où $AM^2 > AH^2$. Les longueurs étant positives, $AM > AH$. \hfill $\blacksquare$
\end{experience}

\begin{popfigure}
\begin{tikzpicture}[scale=1.15]
  % Plan P
  \fill[popblueL] (-2.8,-0.6) -- (1.6,-0.6) -- (3.0,1.2) -- (-1.4,1.2) -- cycle;
  \draw[popblue] (-2.8,-0.6) -- (1.6,-0.6) -- (3.0,1.2) -- (-1.4,1.2) -- cycle;
  \node[popblue] at (-2.3,0.95) {$(P)$};
  % Points
  \coordinate (H) at (0.2,0.1);
  \coordinate (A) at (0.2,2.6);
  \coordinate (M) at (2.0,0.75);
  % Segments
  \draw[poporange,very thick] (A) -- (H);
  \draw[poppurple,very thick] (A) -- (M);
  \draw[popdark,thick,dashed] (H) -- (M);
  % Angle droit en H
  \draw[popgreen,thick] (0.2,0.5) -- (0.47,0.55) -- (0.47,0.15);
  % Points et labels
  \fill[popink] (A) circle (1.6pt); \node[above,popink] at (A) {$A$};
  \fill[popink] (H) circle (1.6pt); \node[below left,popink] at (H) {$H$};
  \fill[popink] (M) circle (1.6pt); \node[right,popink] at (M) {$M$};
  \node[poporange,left] at (0.15,1.4) {$AH$};
  \node[poppurple,above] at (1.35,1.75) {$AM$};
\end{tikzpicture}
\legende{Le triangle $AHM$ est rectangle en $H$ : par Pythagore, $AM^2 = AH^2 + HM^2$, donc $AH < AM$. La perpendiculaire est le plus court chemin de $A$ au plan $(P)$.}
\end{popfigure}

\begin{exoresolu}[Calcul d'une distance à un plan]
Dans le cube $ABCDEFGH$ d'arête $a = 6$ cm, on a montré que $(AE) \perp (ABC)$. Déterminer la distance du point $E$ au plan $(ABC)$, puis la distance de $E$ au point $C$.
\tcblower
Comme $(AE) \perp (ABC)$, le projeté orthogonal de $E$ sur $(ABC)$ est $A$, donc :
\[ d(E, (ABC)) = EA = 6 \text{ cm}. \]
Pour $EC$ : le triangle $EAC$ est rectangle en $A$ (car $(EA) \perp (AC)$, puisque $(AC) \subset (ABC)$). La diagonale du carré $ABCD$ vaut $AC = a\sqrt{2} = 6\sqrt{2}$ cm. Par le théorème de Pythagore :
\[ EC^2 = EA^2 + AC^2 = 36 + 72 = 108, \quad \text{donc} \quad EC = \sqrt{108} = 6\sqrt{3} \approx 10{,}4 \text{ cm}. \]
On vérifie bien que $EA = 6 < EC \approx 10{,}4$ : la distance au plan est bien la plus courte distance.
\end{exoresolu}

\courssub{Plan médiateur d'un segment}

\textbf{Intuition.} Deux antennes relais, l'une MTN en un point $A$, l'autre Orange en un point $B$, couvrent une région. Où se trouvent les points de l'espace situés à \emph{égale distance} des deux antennes ? En géométrie plane, la réponse est la médiatrice du segment $[AB]$. Dans l'espace, la réponse est... tout un plan !

\begin{definition}[Plan médiateur]
Soit $[AB]$ un segment de l'espace. Le \cle{plan médiateur} de $[AB]$ est le plan passant par le milieu $I$ de $[AB]$ et orthogonal à la droite $(AB)$.
\end{definition}

\begin{propriete}[Caractérisation par l'équidistance]
Le plan médiateur de $[AB]$ est l'ensemble des points $M$ de l'espace équidistants de $A$ et de $B$ :
\[ M \in (\mathcal{P}) \iff MA = MB. \]
\emph{(Démonstration exigible.)}
\end{propriete}

\begin{experience}[Démonstration guidée]
Notons $(\mathcal{P})$ le plan passant par le milieu $I$ de $[AB]$ et orthogonal à $(AB)$.

\textbf{Sens direct.} Soit $M$ un point de $(\mathcal{P})$, distinct de $I$. La droite $(AB)$ est orthogonale au plan $(\mathcal{P})$, donc orthogonale à toute droite de $(\mathcal{P})$, en particulier à $(IM)$. Comme $(AB)$ et $(IM)$ sont sécantes en $I$, elles sont perpendiculaires : le triangle $AIM$ est rectangle en $I$, et le triangle $BIM$ aussi. Par le théorème de Pythagore :
\[ MA^2 = MI^2 + IA^2 \qquad \text{et} \qquad MB^2 = MI^2 + IB^2. \]
Or $I$ est le milieu de $[AB]$, donc $IA = IB$, d'où $MA^2 = MB^2$, puis $MA = MB$. (Si $M = I$, on a directement $IA = IB$.)

\textbf{Sens réciproque.} Soit $M$ un point de l'espace tel que $MA = MB$. Si $M$ appartient à la droite $(AB)$, alors $M = I \in (\mathcal{P})$. Sinon, les points $A$, $B$, $M$ définissent un plan, et dans ce plan le triangle $MAB$ est isocèle en $M$ : la médiane $(MI)$ est aussi hauteur, donc $(MI) \perp (AB)$. La droite $(AB)$ est ainsi orthogonale à la droite $(MI)$ en $I$. Par construction, $(AB)$ est orthogonale à toute droite de $(\mathcal{P})$ passant par $I$. Le plan passant par $I$ et orthogonal à $(AB)$ étant \emph{unique} (propriété d'unicité admise), et $(MI)$ étant une droite passant par $I$ et orthogonale à $(AB)$, la droite $(MI)$ est contenue dans $(\mathcal{P})$ : donc $M \in (\mathcal{P})$. \hfill $\blacksquare$
\end{experience}

\begin{popfigure}
\begin{tikzpicture}[scale=1.05]
  % Plan médiateur
  \fill[poppurpleL] (-1.2,-1.1) -- (1.8,-1.1) -- (2.9,0.9) -- (-0.1,0.9) -- cycle;
  \draw[poppurple] (-1.2,-1.1) -- (1.8,-1.1) -- (2.9,0.9) -- (-0.1,0.9) -- cycle;
  \node[poppurple] at (2.35,0.75) {$(\mathcal{P})$};
  % Segment AB traversant le plan
  \coordinate (I) at (0.85,-0.1);
  \coordinate (A) at (0.85,2.3);
  \coordinate (B) at (0.85,-2.1);
  \draw[popink,very thick] (A) -- (B);
  \fill[popink] (A) circle (1.8pt); \node[above,popink] at (A) {$A$};
  \fill[popink] (B) circle (1.8pt); \node[below,popink] at (B) {$B$};
  \fill[poporange] (I) circle (1.8pt); \node[above right,poporange] at (I) {$I$};
  % Point M du plan
  \coordinate (M) at (2.1,-0.55);
  \fill[popgreen] (M) circle (1.8pt); \node[right,popgreen] at (M) {$M$};
  \draw[popgreen,thick,dashed] (I) -- (M);
  \draw[popgreen,thick] (A) -- (M) -- (B);
  % Angle droit en I
  \draw[poporange,thick] (0.85,0.25) -- (1.12,0.20) -- (1.12,-0.12);
  % Codage milieu
  \node[popink,left] at (0.8,1.1) {\small $IA$};
  \node[popink,left] at (0.8,-1.15) {\small $IB$};
\end{tikzpicture}
\legende{Le plan médiateur $(\mathcal{P})$ de $[AB]$ passe par le milieu $I$ et est orthogonal à $(AB)$. Tout point $M$ de $(\mathcal{P})$ vérifie $MA = MB$.}
\end{popfigure}

\begin{attention}[Médiatrice ou plan médiateur ?]
Ne confonds pas les deux objets : dans le \textbf{plan}, l'ensemble des points équidistants de $A$ et $B$ est la \emph{médiatrice} (une droite) ; dans l'\textbf{espace}, c'est le \emph{plan médiateur} (un plan). Dans un exercice d'espace, répondre « la médiatrice de $[AB]$ » est une erreur : une droite ne peut pas contenir tous les points équidistants de $A$ et $B$ dans l'espace.
\end{attention}

\begin{savaistu}[Le plan médiateur autour de nous]
Le plan médiateur intervient partout : en téléphonie mobile, les ingénieurs de MTN ou Orange délimitent les zones de couverture de deux antennes voisines grâce aux surfaces d'équidistance ; en cristallographie, les plans médiateurs définissent les « cellules de Voronoï » qui décrivent la structure des cristaux ; et lors de la construction d'une route entre deux villes du Cameroun, le plan médiateur indique les lieux à égale distance des deux agglomérations — utile pour positionner une relais ou une station-service !
\end{savaistu}

\begin{exercice}[Pour t'entraîner]
Soit $ABCDEFGH$ un cube d'arête $5$ cm.
\begin{enumerate}
  \item Démontrer que la droite $(DH)$ est orthogonale au plan $(ABC)$.
  \item En déduire que $(DH) \perp (AC)$.
  \item Calculer la distance du point $H$ au plan $(ABC)$.
  \item Déterminer la nature de l'ensemble des points de l'espace équidistants de $A$ et de $G$, et préciser par quel point remarquable du cube il passe.
\end{enumerate}
\end{exercice}

\begin{corrige}[Exercice]
\begin{enumerate}
  \item Les faces $ADHE$ et $CDHG$ sont des carrés, donc $(DH) \perp (DA)$ et $(DH) \perp (DC)$. Les droites $(DA)$ et $(DC)$ sont deux droites du plan $(ABC)$, sécantes en $D$. D'après le théorème fondamental, $(DH) \perp (ABC)$.
  \item La droite $(AC)$ est contenue dans le plan $(ABC)$ ; or $(DH) \perp (ABC)$, donc par le théorème réciproque, $(DH)$ est orthogonale à toute droite de $(ABC)$, en particulier $(DH) \perp (AC)$.
  \item Le projeté orthogonal de $H$ sur $(ABC)$ est $D$, donc $d(H, (ABC)) = HD = 5$ cm.
  \item L'ensemble des points équidistants de $A$ et de $G$ est le plan médiateur du segment $[AG]$ (grande diagonale du cube). Il passe par le milieu de $[AG]$, qui est le centre du cube, et il est orthogonal à la droite $(AG)$.
\end{enumerate}
\end{corrige}

\begin{aretenir}[L'essentiel de la section]
\begin{itemize}
  \item $(d) \perp (P)$ signifie : $(d)$ orthogonale à \textbf{toute} droite de $(P)$ ; alors $(d)$ coupe $(P)$ en un unique point, le pied de la perpendiculaire.
  \item \textbf{Critère fondamental} : il suffit de deux droites \textbf{sécantes} de $(P)$ orthogonales à $(d)$.
  \item Par un point, il passe une unique perpendiculaire à un plan donné, et un unique plan orthogonal à une droite donnée.
  \item Distance de $A$ à $(P)$ : $d(A,(P)) = AH$, plus courte distance (Pythagore).
  \item Plan médiateur de $[AB]$ : plan orthogonal à $(AB)$ en son milieu $I$, ensemble des points $M$ tels que $MA = MB$.
\end{itemize}
\end{aretenir}

\coursec{Orthogonalité et parallélisme : théorèmes de liaison}

Dans la section précédente, tu as appris à reconnaître et à démontrer qu'une droite est orthogonale à un plan : il suffit qu'elle soit orthogonale à deux droites sécantes de ce plan. Cette notion, apparemment simple, va nous permettre de répondre à une question beaucoup plus riche : \emph{quel est le lien entre l'orthogonalité et le parallélisme dans l'espace ?}

Observe un instant la cour d'un lycée de Yaoundé un jour de cérémonie : les poteaux qui soutiennent la banderole sont tous plantés verticalement dans le sol, et ils sont visiblement parallèles entre eux. De même, les quatre arêtes verticales d'une armoire sont parallèles, et chacune est perpendiculaire au sol. Ce n'est pas un hasard : c'est un théorème de géométrie. Toute cette section est consacrée à quatre théorèmes qui font la navette entre deux mondes : le monde du \cle{parallélisme} et le monde de l'\cle{orthogonalité}. Ces théorèmes sont des outils de démonstration extrêmement puissants au Probatoire, car ils permettent de transformer un problème de parallélisme (souvent difficile à établir directement) en un problème d'orthogonalité (souvent plus facile), et réciproquement.

\courssub{Théorème 1 : deux droites parallèles et un plan orthogonal}

Commençons par l'intuition. Si deux poteaux sont parallèles et que le premier est bien vertical (orthogonal au sol), alors le second est forcément vertical lui aussi : un poteau parallèle à un poteau vertical ne peut pas être penché.

\begin{propriete}[Théorème 1 : transport de l'orthogonalité par parallélisme des droites]
Soient $(d)$ et $(d')$ deux droites \textbf{parallèles} et $(P)$ un plan. Si $(d)$ est orthogonale à $(P)$, alors $(d')$ est aussi orthogonale à $(P)$.
\[\left.\begin{array}{r} (d) \parallel (d') \\ (d) \perp (P) \end{array}\right\} \Longrightarrow (d') \perp (P)\]
\end{propriete}

\begin{experience}[Démonstration du théorème 1]
\textbf{Hypothèses :} $(d) \parallel (d')$ et $(d) \perp (P)$ en un point $H$.

\textbf{Étape 1 : $(d')$ coupe $(P)$.} Supposons par l'absurde que $(d')$ ne coupe pas $(P)$. Alors $(d') \parallel (P)$. Comme $(d) \parallel (d')$, on a aussi $(d) \parallel (P)$ (par transitivité du parallélisme droite-plan). Mais $(d) \perp (P)$ implique que $(d)$ coupe $(P)$ en $H$ : une droite ne peut être à la fois parallèle à un plan et le couper. Contradiction. Donc $(d')$ perce $(P)$ en un point $H'$.

\textbf{Étape 2 : orthogonalité avec toute droite de $(P)$.} Soit $(\Delta)$ une droite quelconque de $(P)$ passant par $H'$. Par $H$, menons la droite $(\Delta_1)$ parallèle à $(\Delta)$ ; elle est contenue dans $(P)$ (car $(P)$ contient $H$ et la direction de $(\Delta)$). Comme $(d) \perp (P)$, on a $(d) \perp (\Delta_1)$. Or deux droites parallèles conservent les angles : puisque $(d') \parallel (d)$ et $(\Delta) \parallel (\Delta_1)$, l'angle de $(d')$ et $(\Delta)$ égale l'angle de $(d)$ et $(\Delta_1)$, c'est-à-dire un angle droit. Donc $(d') \perp (\Delta)$.

\textbf{Étape 3 : conclusion.} $(d')$ est orthogonale à toute droite de $(P)$ passant par $H'$, en particulier à deux droites sécantes de $(P)$. Par le critère d'orthogonalité droite-plan, $(d') \perp (P)$. $\blacksquare$
\end{experience}

\courssub{Théorème 2 : deux plans parallèles et une droite orthogonale}

L'intuition est tout aussi naturelle : le sol et le plafond d'une salle de classe sont parallèles. Si un fil à plomb est perpendiculaire au sol, il est aussi perpendiculaire au plafond.

\begin{propriete}[Théorème 2 : transport de l'orthogonalité par parallélisme des plans]
Soient $(P)$ et $(Q)$ deux plans \textbf{parallèles} et $(d)$ une droite. Si $(d)$ est orthogonale à $(P)$, alors $(d)$ est aussi orthogonale à $(Q)$.
\[\left.\begin{array}{r} (P) \parallel (Q) \\ (d) \perp (P) \end{array}\right\} \Longrightarrow (d) \perp (Q)\]
\end{propriete}

\begin{experience}[Démonstration du théorème 2]
\textbf{Hypothèses :} $(P) \parallel (Q)$ et $(d) \perp (P)$ en $H$.

\textbf{Étape 1 : $(d)$ coupe $(Q)$.} Si $(d) \parallel (Q)$, comme $(Q) \parallel (P)$, on a $(d) \parallel (P)$ par transitivité. Mais $(d) \perp (P)$ implique que $(d)$ coupe $(P)$ : contradiction. Donc $(d)$ perce $(Q)$ en un point $K$.

\textbf{Étape 2 : transfert des directions.} Soient $(\Delta_1)$ et $(\Delta_2)$ deux droites sécantes de $(Q)$ passant par $K$. Comme $(P) \parallel (Q)$, par le point $H$ de $(P)$ on peut mener des droites $(\Delta'_1)$ et $(\Delta'_2)$ respectivement parallèles à $(\Delta_1)$ et $(\Delta_2)$, et ces droites sont contenues dans $(P)$.

\textbf{Étape 3 : conclusion.} Puisque $(d) \perp (P)$, on a $(d) \perp (\Delta'_1)$ et $(d) \perp (\Delta'_2)$. Le parallélisme conservant les angles, $(d) \perp (\Delta_1)$ et $(d) \perp (\Delta_2)$. La droite $(d)$ est donc orthogonale à deux droites sécantes de $(Q)$ : par le critère d'orthogonalité, $(d) \perp (Q)$. $\blacksquare$
\end{experience}

\begin{attention}[Sens des implications]
Dans les théorèmes 1 et 2, c'est le \textbf{parallélisme} qui est l'hypothèse et l'\textbf{orthogonalité} qui est la conclusion. Ne les confonds pas avec les théorèmes 3 et 4 ci-dessous, où c'est l'inverse ! Rédiger « $(d) \perp (P)$ et $(d') \perp (P)$ donc $(d) \parallel (d')$ » relève du théorème 3, pas du théorème 1. Cite toujours précisément les hypothèses utilisées.
\end{attention}

\courssub{Théorème 3 : deux droites orthogonales à un même plan}

Voici le théorème le plus utilisé de toute la leçon au Probatoire. Retour à notre cour de lycée : tous les poteaux verticaux sont orthogonaux au même plan (le sol), et c'est précisément \emph{pour cela} qu'ils sont parallèles entre eux.

\begin{propriete}[Théorème 3 — exigible au Probatoire]
Deux droites orthogonales à un même plan sont parallèles entre elles.
\[\left.\begin{array}{r} (d) \perp (P) \\ (d') \perp (P) \end{array}\right\} \Longrightarrow (d) \parallel (d')\]
\end{propriete}

\begin{experience}[Démonstration du théorème 3]
\textbf{Hypothèses :} $(d) \perp (P)$ en $H$ et $(d') \perp (P)$ en $H'$.

\textbf{Étape 1 : construction d'une auxiliaire.} Par le point $H'$, menons la droite $(d_1)$ parallèle à $(d)$. D'après le théorème 1, puisque $(d) \parallel (d_1)$ et $(d) \perp (P)$, on a $(d_1) \perp (P)$.

\textbf{Étape 2 : unicité de la perpendiculaire.} Ainsi, $(d')$ et $(d_1)$ sont deux droites passant par $H'$ et toutes deux orthogonales à $(P)$. Or on admet (propriété d'unicité vue en section 3) que \emph{par un point donné, il passe une unique droite orthogonale à un plan donné}. Donc $(d') = (d_1)$.

\textbf{Étape 3 : conclusion.} Comme $(d_1) \parallel (d)$ par construction et que $(d') = (d_1)$, on conclut $(d) \parallel (d')$. $\blacksquare$
\end{experience}

\begin{methode}[Pour montrer que deux droites sont parallèles]
Dans l'espace, on ne dispose pas des angles alternes-internes comme dans le plan. Une stratégie très efficace :
\begin{enumerate}
\item Repérer (ou construire) un plan $(P)$ auquel les deux droites semblent orthogonales.
\item Démontrer $(d) \perp (P)$ : montrer que $(d)$ est orthogonale à deux droites sécantes de $(P)$.
\item Démontrer de même $(d') \perp (P)$.
\item Conclure par le théorème 3 : $(d) \parallel (d')$.
\end{enumerate}
Variante : pour montrer qu'une droite est orthogonale à un plan, on peut aussi montrer qu'elle est parallèle à une droite déjà connue comme orthogonale à ce plan (théorème 1).
\end{methode}

\courssub{Théorème 4 : deux plans orthogonaux à une même droite}

Version « duale » du théorème 3 : les étagères d'une bibliothèque, toutes perpendiculaires au même montant vertical, sont parallèles entre elles.

\begin{propriete}[Théorème 4 — exigible au Probatoire]
Deux plans orthogonaux à une même droite sont parallèles entre eux.
\[\left.\begin{array}{r} (P) \perp (d) \\ (Q) \perp (d) \end{array}\right\} \Longrightarrow (P) \parallel (Q)\]
\end{propriete}

\begin{experience}[Démonstration du théorème 4]
\textbf{Hypothèses :} $(d) \perp (P)$ en $H$ et $(d) \perp (Q)$ en $K$.

\textbf{Étape 1 : construction d'une auxiliaire.} Par le point $K$, menons le plan $(P_1)$ parallèle à $(P)$. D'après le théorème 2, puisque $(P) \parallel (P_1)$ et $(d) \perp (P)$, on a $(d) \perp (P_1)$.

\textbf{Étape 2 : unicité du plan perpendiculaire.} Ainsi $(Q)$ et $(P_1)$ sont deux plans passant par $K$ et tous deux orthogonaux à $(d)$. Or \emph{par un point donné, il passe un unique plan orthogonal à une droite donnée}. Donc $(Q) = (P_1)$.

\textbf{Étape 3 : conclusion.} Comme $(P_1) \parallel (P)$ par construction et $(Q) = (P_1)$, on conclut $(P) \parallel (Q)$. $\blacksquare$
\end{experience}

\begin{popfigure}
\begin{tikzpicture}[scale=1.05]
% Sol (plan P) en perspective
\fill[popblueL,opacity=0.55] (0,0) -- (4.6,0) -- (5.8,1.7) -- (1.2,1.7) -- cycle;
\draw[popblue] (0,0) -- (4.6,0) -- (5.8,1.7) -- (1.2,1.7) -- cycle;
\node[popblue] at (5.15,1.35) {$(P)$};
% Poteau 1
\draw[line width=2.2pt,poporange] (1.6,0.55) -- (1.6,3.1);
\fill[popdark] (1.6,0.55) circle (1.6pt);
\node[below,popdark] at (1.6,0.5) {$H$};
\node[left,poporange] at (1.6,2.9) {$(d)$};
% Poteau 2
\draw[line width=2.2pt,popgreen] (3.9,1.15) -- (3.9,3.7);
\fill[popdark] (3.9,1.15) circle (1.6pt);
\node[below,popdark] at (3.9,1.1) {$H'$};
\node[right,popgreen] at (3.9,3.5) {$(d')$};
% Angles droits
\draw[popdark] (1.6,0.55) ++(0,0.28) -- ++(0.22,0.05) -- ++(0,-0.28);
\draw[popdark] (3.9,1.15) ++(0,0.28) -- ++(0.22,0.05) -- ++(0,-0.28);
% Flèche de conclusion
\draw[<->,>=Stealth,popink,thick] (1.85,3.35) -- (3.65,3.95) node[midway,above,sloped,popink] {\small parallèles};
\end{tikzpicture}
\legende{Deux poteaux $(d)$ et $(d')$ orthogonaux au même sol $(P)$ : le théorème 3 garantit qu'ils sont parallèles.}
\end{popfigure}

\courssub{Ce qui est vrai... et ce qui est faux : réciproques et contre-exemples}

Les quatre théorèmes précédents pourraient laisser croire que « orthogonalité à un même objet » entraîne toujours le parallélisme. C'est vrai quand l'objet commun est un \textbf{plan} (pour des droites) ou une \textbf{droite} (pour des plans). Mais attention :

\begin{attention}[La réciproque qui est FAUSSE dans l'espace]
Dans le \textbf{plan}, deux droites perpendiculaires à une même droite sont toujours parallèles. Dans l'\textbf{espace}, c'est faux : deux droites orthogonales à une même droite ne sont pas nécessairement parallèles. Elles peuvent être sécantes, ou même non coplanaires ! Retenir : l'orthogonalité à une même \emph{droite} ne suffit pas à conclure au parallélisme de deux droites ; il faut l'orthogonalité à un même \emph{plan}.
\end{attention}

\begin{exemplebox}[Contre-exemple sur le cube]
Soit $ABCDEFGH$ un cube, avec $ABCD$ la face du bas et $E$, $F$, $G$, $H$ les sommets correspondants de la face du haut.
\begin{itemize}
\item L'arête $(AB)$ est orthogonale à l'arête $(AE)$ (angle droit de la face $ABFE$).
\item L'arête $(AD)$ est aussi orthogonale à $(AE)$ (angle droit de la face $ADHE$).
\item Pourtant $(AB)$ et $(AD)$ sont \textbf{sécantes en $A$} et perpendiculaires entre elles : elles ne sont pas parallèles !
\end{itemize}
Ainsi $(AB) \perp (AE)$ et $(AD) \perp (AE)$, mais $(AB) \nparallel (AD)$. La « réciproque » naïve du théorème 3 est donc bien fausse dans l'espace.
\end{exemplebox}

\begin{popfigure}
\begin{tikzpicture}[scale=1.15]
% Cube
\coordinate (A) at (0,0); \coordinate (B) at (2.4,0);
\coordinate (D) at (1.1,1.0); \coordinate (C) at (3.5,1.0);
\coordinate (E) at (0,2.4); \coordinate (F) at (2.4,2.4);
\coordinate (H) at (1.1,3.4); \coordinate (G) at (3.5,3.4);
\draw[popmuted,dashed] (A)--(D) (D)--(C) (D)--(H);
\draw[popdark] (A)--(B)--(C) (B)--(F) (C)--(G) (E)--(F)--(G)--(H)--(E) (A)--(E) (F)--(G);
\draw[line width=2pt,poporange] (A)--(B);
\draw[line width=2pt,popgreen,dashed] (A)--(D);
\draw[line width=2pt,poppurple] (A)--(E);
\foreach \p/\n/\pos in {A/A/below left,B/B/below,D/D/above left,E/E/left}
  \node[\pos,popdark] at (\p) {$\n$};
\node[poporange,below] at (1.2,0) {\small $(AB)$};
\node[popgreen] at (0.35,0.62) {\small $(AD)$};
\node[poppurple,left] at (0,1.2) {\small $(AE)$};
\end{tikzpicture}
\legende{Dans le cube, $(AB) \perp (AE)$ et $(AD) \perp (AE)$, pourtant $(AB)$ et $(AD)$ sont sécantes : l'orthogonalité à une même droite n'implique pas le parallélisme.}
\end{popfigure}

\courssub{Tableau de synthèse des implications}

\begin{center}
\begin{tabularx}{\linewidth}{|l|X|X|}\hline
\rowcolor{popblueL} \textbf{Hypothèses} & \textbf{Conclusion} & \textbf{Statut} \\ \hline
$(d) \parallel (d')$ et $(d) \perp (P)$ & $(d') \perp (P)$ & Vrai (théorème 1) \\ \hline
$(P) \parallel (Q)$ et $(d) \perp (P)$ & $(d) \perp (Q)$ & Vrai (théorème 2) \\ \hline
$(d) \perp (P)$ et $(d') \perp (P)$ & $(d) \parallel (d')$ & Vrai (théorème 3) \\ \hline
$(P) \perp (d)$ et $(Q) \perp (d)$ & $(P) \parallel (Q)$ & Vrai (théorème 4) \\ \hline
$(d) \perp (\Delta)$ et $(d') \perp (\Delta)$ & $(d) \parallel (d')$ & \textbf{Faux} dans l'espace (contre-exemple du cube) \\ \hline
$(P) \perp (\pi)$ et $(Q) \perp (\pi)$ & $(P) \parallel (Q)$ & \textbf{Faux} : deux murs verticaux sont tous deux perpendiculaires au sol et peuvent se couper \\ \hline
\end{tabularx}
\end{center}

\begin{aretenir}[Le réflexe de synthèse]
Le parallélisme « transporte » l'orthogonalité (théorèmes 1 et 2), et l'orthogonalité à un même objet « fabrique » du parallélisme (théorèmes 3 et 4) — à condition que l'objet commun soit du \emph{bon type} : un plan pour deux droites, une droite pour deux plans.
\end{aretenir}

\courssub{Applications}

\textbf{Application 1 : les arêtes latérales d'un pavé droit sont parallèles.} Soit $ABCDEFGH$ un pavé droit de base $ABCD$. Par définition d'un pavé droit, les faces latérales sont des rectangles, donc $(AE) \perp (AB)$ et $(AE) \perp (AD)$. Comme $(AB)$ et $(AD)$ sont deux droites sécantes du plan $(ABC)$, on a $(AE) \perp (ABC)$. Le même raisonnement donne $(BF) \perp (ABC)$, $(CG) \perp (ABC)$ et $(DH) \perp (ABC)$. Par le théorème 3, les quatre arêtes $(AE)$, $(BF)$, $(CG)$, $(DH)$, orthogonales au même plan $(ABC)$, sont parallèles entre elles. C'est exactement pourquoi les montants d'une armoire ou d'un immeuble sont parallèles.

\textbf{Application 2 : la hauteur d'une pyramide.} Soit $SABCD$ une pyramide de sommet $S$ et de base le carré $ABCD$, et soit $O$ le centre de la base. On dit que $[SO]$ est la \cle{hauteur} de la pyramide lorsque $(SO) \perp (ABC)$. Conséquence précieuse : si une seconde pyramide $S'ABCD$ a aussi sa hauteur en $O$, alors $(S'O) \perp (ABC)$, et le théorème 3 donne $(SO) \parallel (S'O)$ : les deux sommets sont sur une même droite verticale. De plus, tout plan parallèle à la base $(ABC)$ — par exemple le plan d'une section de la pyramide — est, par le théorème 2, orthogonal à $(SO)$ : c'est ce qui justifie que les sections planes d'une pyramide par un plan parallèle à la base sont « horizontales » et homothétiques à la base.

\begin{exoresolu}[Dans un pavé droit]
\textbf{Énoncé.} Soit $ABCDEFGH$ un pavé droit tel que $AB = 4$~cm, $AD = 3$~cm et $AE = 5$~cm. On note $I$ le milieu de $[CG]$.
\begin{enumerate}
\item Démontrer que $(BF) \parallel (DH)$.
\item En déduire que les points $B$, $F$, $H$, $D$ sont coplanaires.
\item Démontrer que la droite $(AE)$ est orthogonale à la droite $(BD)$.
\end{enumerate}
\tcblower
\textbf{Solution.}
\begin{enumerate}
\item Les faces $ABFE$ et $ADHE$ sont des rectangles, donc $(BF) \perp (AB)$ et $(BF) \perp (AD)$ : ainsi $(BF) \perp (ABC)$. De même $(DH) \perp (ABC)$. Les droites $(BF)$ et $(DH)$ sont orthogonales au même plan $(ABC)$ : d'après le théorème 3, $(BF) \parallel (DH)$.
\item Deux droites parallèles sont toujours coplanaires : $(BF)$ et $(DH)$ définissent un plan, qui contient les quatre points $B$, $F$, $H$, $D$.
\item On a montré $(AE) \perp (ABC)$ (orthogonalité à $(AB)$ et $(AD)$, sécantes en $A$). Or la diagonale $(BD)$ est contenue dans le plan $(ABC)$. Une droite orthogonale à un plan est orthogonale à toute droite de ce plan : donc $(AE) \perp (BD)$. (Remarque : ces deux droites sont orthogonales sans être sécantes — elles ne sont pas perpendiculaires.)
\end{enumerate}
\end{exoresolu}

\begin{savaistu}[Pourquoi les immeubles ne s'écroulent pas (géométriquement parlant)]
Les ingénieurs du génie civil, à Douala comme ailleurs, exploitent en permanence le théorème 3 : pour garantir que les poteaux d'un immeuble sont rigoureusement parallèles, il suffit de vérifier que chacun est vertical, c'est-à-dire orthogonal au plan horizontal du sol — ce qu'on contrôle avec un simple fil à plomb. Une seule vérification par poteau, et le parallélisme global est assuré par le théorème. La géométrie de Première est littéralement dans les murs qui t'entourent !
\end{savaistu}

\coursec{Plans orthogonaux}

Dans la section précédente, nous avons établi des ponts précieux entre orthogonalité et parallélisme : deux droites parallèles à une même droite orthogonale à un plan sont orthogonales à ce plan, deux plans parallèles partagent les mêmes droites orthogonales, etc. Ces résultats vont maintenant nous servir d'outils pour aborder la dernière grande situation d'orthogonalité de cette leçon : celle qui concerne \textbf{deux plans}. Regarde autour de toi : le mur de ta salle de classe et le sol, l'écran de ton téléphone posé verticalement sur une table, les faces d'une boîte de craie... Partout, des plans se coupent « à angle droit ». Cette section donne un sens précis à cette expression et fournit les outils pour la démontrer.

\courssub{Intuition et définition}

\textbf{De l'observation à l'idée mathématique.}\par

Prenons une porte ouverte dans une maison de Yaoundé. Le mur est un plan vertical, le sol est un plan horizontal : ils semblent se couper « perpendiculairement ». Mais comment traduire cela mathématiquement ? Une première idée serait de dire que « toute droite du mur est orthogonale à toute droite du sol ». C'est faux ! La ligne de contact entre le mur et le sol appartient aux deux plans, et elle n'est pas orthogonale à elle-même. Il faut donc une définition plus subtile.

L'idée correcte est la suivante : le mur contient des droites \emph{verticales} (par exemple le montant de la porte), et ces droites verticales sont orthogonales au plan du sol tout entier. C'est l'existence d'\textbf{une seule} droite bien choisie — une droite d'un plan, orthogonale à l'autre plan — qui capture la notion.

\begin{definition}[Plans orthogonaux]
Deux plans $(P)$ et $(Q)$ distincts sont dits \cle{orthogonaux} (ou \cle{perpendiculaires}) lorsque l'un des deux contient une droite orthogonale à l'autre plan. On note alors :
\[ (P) \perp (Q). \]
Autrement dit : $(P) \perp (Q)$ s'il existe une droite $(d)$ telle que $(d) \subset (P)$ et $(d) \perp (Q)$.
\end{definition}

\begin{attention}[La définition porte sur les plans, pas sur toutes leurs droites]
Dire que $(P) \perp (Q)$ ne signifie \textbf{pas} que toute droite de $(P)$ est orthogonale à toute droite de $(Q)$. Cela signifie seulement qu'il \emph{existe} une droite de $(P)$ orthogonale au plan $(Q)$ tout entier. La droite d'intersection des deux plans, par exemple, appartient aux deux plans sans être orthogonale à quoi que ce soit d'intéressant ici.
\end{attention}

\courssub{Cohérence de la définition : la symétrie de la relation}

La définition semble donner un rôle privilégié à l'un des deux plans (« l'un contient une droite orthogonale à l'autre »). Est-ce que le second plan rend la pareille ? La réponse est oui, et c'est ce qui rend la définition cohérente.

\begin{propriete}[Symétrie de l'orthogonalité des plans]
Si le plan $(P)$ contient une droite orthogonale au plan $(Q)$, alors le plan $(Q)$ contient aussi une droite orthogonale au plan $(P)$. On peut donc dire sans ambiguïté que $(P)$ et $(Q)$ sont orthogonaux, la relation étant \cle{symétrique} :
\[ (P) \perp (Q) \iff (Q) \perp (P). \]
\end{propriete}

\textbf{Illustration de cette symétrie.}\par

Soient $(P)$ et $(Q)$ deux plans sécants selon une droite $(\Delta)$, et $(d)$ une droite de $(P)$ orthogonale à $(Q)$, perçant $(Q)$ en un point $H$. Menons dans $(Q)$ la droite $(d')$ passant par $H$ et orthogonale à $(\Delta)$. Comme $(d) \perp (Q)$, la droite $(d)$ est orthogonale à toute droite de $(Q)$, donc à $(d')$ et à $(\Delta)$. On montre alors que $(d')$, étant orthogonale à la fois à $(\Delta)$ et à $(d)$ qui sont deux droites sécantes de $(P)$, est orthogonale au plan $(P)$. Ainsi $(Q)$ contient bien une droite orthogonale à $(P)$.

\begin{popfigure}
\begin{center}
\begin{tikzpicture}[scale=1.05]
% Plan Q (horizontal, en perspective)
\fill[popblueL,opacity=0.55] (-3.2,-1.1) -- (3.6,-1.1) -- (2.4,1.0) -- (-4.4,1.0) -- cycle;
\draw[popblue,thick] (-3.2,-1.1) -- (3.6,-1.1) -- (2.4,1.0) -- (-4.4,1.0) -- cycle;
\node[popblue] at (-3.6,-0.7) {$(Q)$};
% Plan P (vertical)
\fill[poporangeL,opacity=0.5] (-1.9,-0.85) -- (1.7,-0.85) -- (1.7,2.6) -- (-1.9,2.6) -- cycle;
\draw[poporange,thick] (-1.9,-0.85) -- (1.7,-0.85) -- (1.7,2.6) -- (-1.9,2.6) -- cycle;
\node[poporange] at (1.35,2.35) {$(P)$};
% Intersection (Delta)
\draw[popdark,very thick] (-2.6,-0.85) -- (2.6,-0.85);
\node[popdark,below] at (2.35,-0.9) {$(\Delta)$};
% Droite d orthogonale à Q (rouge)
\draw[red,very thick] (0,-0.85) -- (0,2.3);
\fill[red] (0,-0.85) circle (1.6pt);
\node[red,right] at (0.08,1.9) {$(d)$};
\node[red,below left] at (-0.05,-0.9) {$H$};
% angle droit
\draw[red] (0,-0.55) -- (0.3,-0.55) -- (0.3,-0.85);
% Droite d' dans Q orthogonale à Delta
\draw[poppurple,very thick] (0,-0.85) -- (-1.1,0.55);
\node[poppurple,above left] at (-1.05,0.4) {$(d')$};
\end{tikzpicture}
\end{center}
\legende{Deux plans orthogonaux $(P)$ et $(Q)$ sécants selon $(\Delta)$. La droite rouge $(d) \subset (P)$ est orthogonale au plan $(Q)$ ; réciproquement, $(d') \subset (Q)$, orthogonale à $(\Delta)$ en $H$, est orthogonale à $(P)$.}
\end{popfigure}

\courssub{Théorème fondamental : critère d'orthogonalité de deux plans}

Voici le résultat central de cette section, celui que tu utiliseras dans presque toutes les démonstrations du Probatoire.

\begin{propriete}[Théorème — Critère d'orthogonalité de deux plans]
Deux plans $(P)$ et $(Q)$ sont orthogonaux \textbf{si et seulement si} l'un d'eux contient une droite orthogonale à l'autre :
\[ (P) \perp (Q) \iff \text{il existe une droite } (d) \subset (P) \text{ telle que } (d) \perp (Q). \]
En pratique, on utilise surtout le sens direct (le \emph{critère}) : \emph{si une droite est orthogonale à un plan, alors tout plan contenant cette droite est orthogonal à ce plan}.
\end{propriete}

\begin{experience}[Démonstration guidée du sens direct (critère)]
\textbf{Hypothèses.} Soit $(d)$ une droite orthogonale au plan $(Q)$, et soit $(P)$ un plan contenant $(d)$. On veut montrer que $(P) \perp (Q)$.

\begin{enumerate}
\item \textbf{$(P)$ et $(Q)$ sont sécants.} La droite $(d)$ perce $(Q)$ en un point $H$ (puisque $(d) \perp (Q)$). Or $H \in (d) \subset (P)$, donc $H$ appartient aux deux plans. Les plans $(P)$ et $(Q)$ ont au moins le point $H$ en commun, et ils ne sont pas confondus (sinon $(d)$, orthogonale à $(Q)$, serait contenue dans $(Q)$, ce qui est absurde). Deux plans distincts ayant un point commun sont sécants : notons $(\Delta) = (P) \cap (Q)$, avec $H \in (\Delta)$.
\item \textbf{Construction d'une droite témoin dans $(Q)$.} Dans le plan $(Q)$, menons par $H$ la droite $(d')$ orthogonale à $(\Delta)$.
\item \textbf{$(d')$ est orthogonale à deux droites sécantes de $(P)$.} D'une part $(d') \perp (\Delta)$ par construction. D'autre part, comme $(d) \perp (Q)$, la droite $(d)$ est orthogonale à \emph{toute} droite de $(Q)$, donc en particulier $(d) \perp (d')$.
\item \textbf{Conclusion.} Les droites $(d)$ et $(\Delta)$ sont deux droites sécantes du plan $(P)$ (elles se coupent en $H$), et $(d')$ leur est orthogonale à toutes les deux. D'après le théorème de la droite orthogonale à un plan (section 3), $(d') \perp (P)$. Le plan $(Q)$ contient donc une droite orthogonale à $(P)$ : par définition, $(P) \perp (Q)$. $\blacksquare$
\end{enumerate}
\end{experience}

\begin{methode}[Montrer que deux plans sont orthogonaux]
Pour démontrer que $(P) \perp (Q)$ :
\begin{enumerate}
\item Repérer (ou construire) dans l'un des deux plans, disons $(P)$, une droite $(d)$ candidate — souvent une arête, une hauteur ou une diagonale de la figure.
\item Montrer que $(d)$ est orthogonale à \textbf{deux droites sécantes} du plan $(Q)$ (c'est le critère d'orthogonalité droite-plan de la section 3). Très souvent, l'une de ces deux droites est l'intersection $(\Delta) = (P) \cap (Q)$.
\item Conclure : $(d) \perp (Q)$ et $(d) \subset (P)$, donc $(P) \perp (Q)$.
\end{enumerate}
Rédaction type : « La droite $(AE)$ est orthogonale aux droites $(AB)$ et $(AD)$, sécantes en $A$ dans le plan $(ABC)$ ; donc $(AE) \perp (ABC)$. Or $(AE) \subset (ABE)$, donc $(ABE) \perp (ABC)$. »
\end{methode}

\begin{exemplebox}[Faces d'un pavé droit]
Dans un pavé droit, chaque arête latérale est orthogonale au plan de la base. Comme chaque face latérale contient une arête latérale, \textbf{chaque face latérale est orthogonale au plan de la base}. De même, deux faces latérales adjacentes sont orthogonales entre elles : la face $(ADHE)$ contient $(AD)$ qui est orthogonale à $(AB)$ et à $(AE)$ (deux droites sécantes du plan $(ABFE)$), donc $(AD) \perp (ABFE)$, ce qui entraîne $(ADHE) \perp (ABFE)$.
\end{exemplebox}

\courssub{Une propriété utile : la droite orthogonale à l'intersection}

\begin{propriete}[Droite orthogonale à l'intersection de deux plans orthogonaux]
Soient $(P)$ et $(Q)$ deux plans orthogonaux, sécants selon une droite $(\Delta)$. Toute droite de $(P)$ orthogonale à $(\Delta)$ est orthogonale au plan $(Q)$.
\end{propriete}

\begin{experience}[Démonstration guidée]
\textbf{Hypothèses.} $(P) \perp (Q)$, $(\Delta) = (P) \cap (Q)$, et $(d) \subset (P)$ avec $(d) \perp (\Delta)$. Montrons que $(d) \perp (Q)$.
\begin{enumerate}
\item Puisque $(P) \perp (Q)$, il existe une droite $(\ell) \subset (P)$ orthogonale à $(Q)$ (définition).
\item Les droites $(d)$ et $(\ell)$ sont toutes deux contenues dans $(P)$ et toutes deux orthogonales à $(\Delta)$ (pour $(\ell)$, c'est parce que $(\ell) \perp (Q)$ et $(\Delta) \subset (Q)$). Dans le plan $(P)$, deux droites orthogonales à une même droite sont parallèles : $(d) \parallel (\ell)$.
\item Or $(\ell) \perp (Q)$. D'après le théorème de liaison orthogonalité-parallélisme (section 4), toute droite parallèle à une droite orthogonale à un plan est orthogonale à ce plan : donc $(d) \perp (Q)$. $\blacksquare$
\end{enumerate}
\end{experience}

Ce résultat est très pratique : dans deux plans orthogonaux, il suffit de repérer une droite perpendiculaire à la « charnière » $(\Delta)$ pour obtenir gratuitement une droite orthogonale à tout un plan. C'est exactement ce qui se passe avec un tableau mural : la ligne tracée au tableau perpendiculairement à la ligne du sol est verticale, donc orthogonale au sol.

\courssub{Exercice résolu dans le cube}

\begin{exoresolu}[Plans orthogonaux dans le cube]
Soit $ABCDEFGH$ un cube, de sommets disposés de façon usuelle ($ABCD$ face inférieure, $EFGH$ face supérieure, $E$ au-dessus de $A$). Démontrer que les plans $(ABE)$ et $(ABC)$ sont orthogonaux.
\tcblower
\textbf{Analyse.} Le plan $(ABC)$ est le plan de la face inférieure. Le plan $(ABE)$ contient la face avant. Leur intersection est la droite $(AB)$. Cherchons dans $(ABE)$ une droite orthogonale à $(ABC)$ : l'arête $(AE)$ est une candidate naturelle.

\textbf{Rédaction.}
\begin{enumerate}
\item $ABCD$ est un carré (face du cube). Dans la face latérale $ABFE$ (carré), $(AE) \perp (AB)$. Dans la face latérale $ADHE$ (carré), $(AE) \perp (AD)$.
\item Les droites $(AB)$ et $(AD)$ sont deux droites sécantes en $A$ du plan $(ABC)$. La droite $(AE)$ leur est orthogonale à toutes les deux, donc :
\[ (AE) \perp (ABC). \]
\item Or $(AE) \subset (ABE)$ puisque $A$ et $E$ définissent ce plan avec $B$. Le plan $(ABE)$ contient donc une droite orthogonale au plan $(ABC)$. Par le critère d'orthogonalité :
\[ (ABE) \perp (ABC). \]
\end{enumerate}
\textbf{Remarque.} On montrerait de même que $(ADE) \perp (ABC)$, $(BCF) \perp (ABC)$, $(DCG) \perp (ABC)$ : les quatre faces latérales du cube sont orthogonales à la base.
\end{exoresolu}

\courssub{Le coin d'une pièce : trois plans deux à deux orthogonaux}

Le coin d'une pièce de maison offre le modèle parfait de \textbf{trois plans deux à deux orthogonaux} : le sol, et deux murs adjacents. Chaque mur contient l'arête verticale du coin, qui est orthogonale au sol : donc chaque mur est orthogonal au sol. De plus, l'arête verticale du coin est l'intersection des deux murs, et chaque mur contient une ligne horizontale (au sol) orthogonale à cette arête et à l'autre mur : les deux murs sont orthogonaux entre eux.

\begin{popfigure}
\begin{center}
\begin{tikzpicture}[scale=1.1]
% Sol
\fill[popgreenL,opacity=0.6] (0,0) -- (3.6,0) -- (4.6,1.4) -- (1.0,1.4) -- cycle;
\draw[popgreen,thick] (0,0) -- (3.6,0) -- (4.6,1.4) -- (1.0,1.4) -- cycle;
\node[popgreen!60!black] at (3.9,0.5) {sol $(S)$};
% Mur gauche
\fill[popblueL,opacity=0.55] (0,0) -- (1.0,1.4) -- (1.0,3.6) -- (0,2.6) -- cycle;
\draw[popblue,thick] (0,0) -- (1.0,1.4) -- (1.0,3.6) -- (0,2.6) -- cycle;
\node[popblue] at (0.35,2.9) {mur $(M_1)$};
% Mur droit
\fill[poporangeL,opacity=0.55] (1.0,1.4) -- (4.6,1.4) -- (4.6,3.6) -- (1.0,3.6) -- cycle;
\draw[poporange,thick] (1.0,1.4) -- (4.6,1.4) -- (4.6,3.6) -- (1.0,3.6) -- cycle;
\node[poporange] at (3.9,3.2) {mur $(M_2)$};
% Arête verticale du coin
\draw[red,very thick] (1.0,1.4) -- (1.0,3.6);
\fill[red] (1.0,1.4) circle (1.8pt);
\node[red,left] at (0.95,2.6) {$(d)$};
% angles droits
\draw[popdark] (1.0,1.7) -- (1.3,1.7) -- (1.3,1.4);
\draw[popdark] (1.0,1.4) -- (1.25,1.55) -- (1.25,1.85);
\end{tikzpicture}
\end{center}
\legende{Le coin d'une pièce : le sol $(S)$ et les deux murs $(M_1)$, $(M_2)$ sont trois plans deux à deux orthogonaux. L'arête verticale $(d) = (M_1) \cap (M_2)$ est orthogonale au sol $(S)$.}
\end{popfigure}

\begin{attention}[Piège classique : deux plans orthogonaux à une même droite... ou contenant des droites orthogonales à une même droite]
\textbf{Erreur fréquente :} croire que si $(P)$ et $(Q)$ contiennent chacun une droite orthogonale à une même droite $(\ell)$, alors $(P) \perp (Q)$. C'est faux ! Contre-exemple : dans le cube, les plans $(ABC)$ (base) et $(EFG)$ (face du dessus) contiennent respectivement $(AB)$ et $(EF)$, toutes deux orthogonales à $(AE)$... et pourtant $(ABC) \parallel (EFG)$, ils ne sont pas orthogonaux du tout !

Retiens : pour conclure que $(P) \perp (Q)$, il faut une droite \textbf{contenue dans l'un} et \textbf{orthogonale à l'autre plan tout entier} — pas à une seule droite, pas à un troisième objet.
\end{attention}

\courssub{Complément pour la Terminale : l'angle diédral}

\begin{definition}[Angle diédral — notion d'ouverture]
Deux demi-plans $(\mathcal{P})$ et $(\mathcal{Q})$ de même frontière $(\Delta)$ définissent un \cle{angle diédral}. Pour le mesurer, on coupe la figure par un plan $(R)$ orthogonal à $(\Delta)$ : il trace sur les demi-plans deux demi-droites $[Ox)$ et $[Oy)$, et l'angle géométrique $\widehat{xOy}$ est la mesure du diédral. Ce plan $(R)$ s'appelle le \cle{plan d'angle droit} du diédral. Deux plans sont orthogonaux exactement lorsque le diédral qu'ils forment mesure $90~degrés$.
\end{definition}

Cette notion sera formalisée en Terminale C et D (géométrie dans l'espace, produit scalaire). Pour l'instant, retiens simplement l'image : l'angle entre deux plans se mesure « en tranche », dans un plan perpendiculaire à leur ligne de contact — comme on mesure l'ouverture d'un livre en regardant ses pages de profil.

\begin{savaistu}[Les écrans géants du stade de Japoma]
Lors de la CAN 2021, le stade de Japoma à Douala a déployé des écrans géants de plusieurs tonnes. Leurs structures de support sont calculées pour que les plans porteurs soient rigoureusement orthogonaux au sol : si un plan de support s'incline, le poids de l'écran crée un moment de basculement que les câbles doivent compenser en permanence. Les ingénieurs vérifient l'orthogonalité avec des niveaux laser, qui matérialisent exactement la droite orthogonale à un plan que tu étudies dans cette leçon. La géométrie de l'espace n'est pas un jeu abstrait : c'est ce qui empêche les écrans de tomber !
\end{savaistu}

\begin{exercice}[Pour t'entraîner]
Soit $SABC$ une pyramide de sommet $S$ dont la base $ABC$ est un triangle rectangle en $B$, et telle que $(SB) \perp (ABC)$.
\begin{enumerate}
\item Démontrer que les plans $(SAB)$ et $(ABC)$ sont orthogonaux.
\item Démontrer que $(BC) \perp (SAB)$, puis que les plans $(SBC)$ et $(SAB)$ sont orthogonaux.
\end{enumerate}
\end{exercice}

\begin{corrige}[Exercice : pyramide]
\begin{enumerate}
\item Par hypothèse $(SB) \perp (ABC)$. Or $S$ et $B$ appartiennent au plan $(SAB)$, donc $(SB) \subset (SAB)$. Le plan $(SAB)$ contient une droite orthogonale à $(ABC)$ : d'après le critère d'orthogonalité, $(SAB) \perp (ABC)$.
\item La droite $(BC)$ est orthogonale à $(SB)$ car $(SB) \perp (ABC)$ et $(BC) \subset (ABC)$. Elle est aussi orthogonale à $(BA)$ car le triangle $ABC$ est rectangle en $B$. Les droites $(SB)$ et $(BA)$ sont sécantes en $B$ dans le plan $(SAB)$, donc $(BC) \perp (SAB)$. Comme $(BC) \subset (SBC)$, le plan $(SBC)$ contient une droite orthogonale à $(SAB)$ : donc $(SBC) \perp (SAB)$.
\end{enumerate}
\end{corrige}

\begin{aretenir}[L'essentiel sur les plans orthogonaux]
\begin{itemize}
\item $(P) \perp (Q)$ signifie : l'un des plans contient une droite orthogonale à l'autre. La relation est symétrique.
\item \textbf{Méthode reine :} pour prouver $(P) \perp (Q)$, exhiber $(d) \subset (P)$ orthogonale à \emph{deux droites sécantes} de $(Q)$.
\item Si $(P) \perp (Q)$ avec $(\Delta) = (P) \cap (Q)$, alors toute droite de $(P)$ orthogonale à $(\Delta)$ est orthogonale à $(Q)$.
\item Piège : contenir chacun une droite orthogonale à une même droite ne suffit pas à rendre deux plans orthogonaux.
\item Modèles à connaître par cœur : faces du cube et du pavé droit, coin d'une pièce (trois plans deux à deux orthogonaux).
\end{itemize}
\end{aretenir}

\coursec{Méthodes de démonstration et applications}

Dans la section précédente, nous avons appris à reconnaître et à démontrer l'orthogonalité de deux plans grâce à un théorème puissant : si un plan contient une droite orthogonale à un autre plan, alors ces deux plans sont orthogonaux. Tu as maintenant en main \emph{tous} les outils théoriques de la leçon. Reste à apprendre à les \emph{utiliser efficacement} : c'est exactement ce que demande le Probatoire, où les questions de géométrie dans l'espace se succèdent selon une logique bien rodée. Cette dernière section te donne les stratégies de preuve, les réflexes de calcul et une étude complète guidée sur les solides classiques.

\courssub{Stratégies générales de démonstration}

\courssub{Prouver qu'une droite est orthogonale à un plan}

Rappelons le théorème fondamental : une droite est orthogonale à un plan si et seulement si elle est orthogonale à \textbf{deux droites sécantes} de ce plan. Toute la difficulté consiste donc à \emph{trouver} ces deux droites sécantes et à justifier chaque orthogonalité.

\begin{methode}[Schéma de preuve type : droite orthogonale à un plan]
Pour montrer qu'une droite $(d)$ est orthogonale à un plan $(P)$ :
\begin{enumerate}
\item \textbf{Repérer} dans $(P)$ deux droites sécantes $(\Delta_1)$ et $(\Delta_2)$, en un point $O$ de préférence situé sur $(d)$.
\item \textbf{Justifier} $(d) \perp (\Delta_1)$ : angle droit d'une face, diagonales d'un carré ou d'un losange, médiane d'un triangle isocèle, etc.
\item \textbf{Justifier} $(d) \perp (\Delta_2)$ par un argument analogue.
\item \textbf{Conclure} en citant le théorème : $(d)$ est orthogonale à deux droites sécantes de $(P)$, donc $(d) \perp (P)$.
\end{enumerate}
\end{methode}

\begin{attention}[Deux droites \emph{sécantes}, pas parallèles !]
L'hypothèse « sécantes » est indispensable : une droite peut être orthogonale à deux droites \emph{parallèles} d'un plan sans être orthogonale au plan. Par exemple, dans un pavé droit, une arête verticale est orthogonale à toutes les arêtes horizontales d'une même direction, pourtant elle n'est orthogonale qu'aux plans horizontaux. Oublier de vérifier que les deux droites se coupent est l'erreur la plus fréquente au Probatoire — et elle coûte cher.
\end{attention}

\courssub{Prouver que deux plans sont orthogonaux}

La stratégie est une \emph{chaîne} en deux maillons, que tu dois rédiger dans l'ordre :

\begin{methode}[Schéma de preuve type : plans orthogonaux]
Pour montrer que $(P) \perp (Q)$ :
\begin{enumerate}
\item \textbf{Chercher} une droite $(d)$ contenue dans $(P)$ qui soit orthogonale à $(Q)$. La démontrer par la méthode précédente (deux droites sécantes de $(Q)$ orthogonales à $(d)$).
\item \textbf{Conclure} : $(P)$ contient la droite $(d)$ orthogonale à $(Q)$, donc $(P) \perp (Q)$.
\end{enumerate}
En pratique, la droite $(d)$ est presque toujours une arête, une hauteur ou une diagonale du solide étudié : c'est l'énoncé (ou l'observation de la figure) qui te la désigne.
\end{methode}

\begin{methode}[Fiche récapitulative des quatre savoir-faire]
\begin{center}
\begin{tabularx}{\linewidth}{|l|X|}
\hline
\rowcolor{popblueL} \textbf{Savoir-faire} & \textbf{Théorème clé à mobiliser} \\
\hline
Droites coplanaires ou non & Deux droites sont coplanaires ssi elles sont sécantes ou parallèles ; sinon, elles ne le sont pas (exhiber un plan contenant l'une et parallèle à l'autre). \\
\hline
Droites orthogonales / perpendiculaires & $(d) \perp (\Delta)$ ssi toute parallèle à $(d)$ est perpendiculaire à toute parallèle à $(\Delta)$ ; en pratique, se ramener à un angle droit dans un plan bien choisi. \\
\hline
Droite orthogonale à un plan & $(d)$ orthogonale à \textbf{deux droites sécantes} de $(P)$ $\Rightarrow$ $(d) \perp (P)$. \\
\hline
Plans orthogonaux & $(P)$ contient une droite orthogonale à $(Q)$ $\Rightarrow$ $(P) \perp (Q)$. \\
\hline
\end{tabularx}
\end{center}
\end{methode}

\courssub{Exploiter le repère implicite d'un solide}

Voici une intuition précieuse : un pavé droit (une caisse de savon, une salle de classe, un carton de boissons) est un \emph{repère déguisé}. En chaque sommet, les trois arêtes sont deux à deux orthogonales, et chaque face est un rectangle : les angles droits sont \emph{partout}, il suffit de les voir.

\begin{propriete}[Angles droits d'un pavé droit]
Dans un pavé droit $ABCDEFGH$ :
\begin{itemize}
\item toute arête est orthogonale aux deux faces qu'elle ne traverse pas parallèlement ; par exemple $(AE) \perp (ABCD)$ et $(AE) \perp (EFGH)$ ;
\item les diagonales d'une face carrée sont perpendiculaires ;
\item toute diagonale de face est hypoténuse d'un triangle rectangle contenu dans cette face.
\end{itemize}
\end{propriete}

Lorsqu'un énoncé dit « soit $ABCDEFGH$ un cube » ou « un pavé droit », tu dois immédiatement traduire : $(AB) \perp (AD)$, $(AB) \perp (AE)$, $(AD) \perp (AE)$, et de même en chaque sommet. C'est ce stock d'orthogonalités gratuites qui alimente toutes les démonstrations.

\courssub{Calculs de distances : Pythagore dans l'espace}

\courssub{Le principe : toujours travailler dans un plan}

Le théorème de Pythagore est un théorème \emph{du plan}. Pour calculer une longueur dans l'espace, on procède donc toujours en deux temps : on \textbf{identifie un plan} contenant un triangle rectangle utile, on y applique Pythagore, puis on recommence si nécessaire dans un second plan.

\begin{exemplebox}[Grande diagonale d'un pavé droit]
Soit $ABCDEFGH$ un pavé droit de dimensions $AB = 12$~cm, $AD = 4$~cm, $AE = 3$~cm. Calculons la grande diagonale $AG$.

\textbf{Étape 1 — dans le plan $(ABCD)$.} La face $ABCD$ est un rectangle, donc le triangle $ABC$ est rectangle en $B$ :
\[ AC^2 = AB^2 + BC^2 = 12^2 + 4^2 = 144 + 16 = 160. \]

\textbf{Étape 2 — dans le plan $(ACG)$.} L'arête $(CG)$ est orthogonale au plan $(ABCD)$, donc perpendiculaire à la droite $(AC)$ qu'il contient : le triangle $ACG$ est rectangle en $C$. D'où :
\[ AG^2 = AC^2 + CG^2 = 160 + 3^2 = 169 \quad \text{donc} \quad AG = 13~\text{cm}. \]
\end{exemplebox}

\begin{savaistu}[Pythagore en dimension 3]
La grande diagonale d'un pavé droit de dimensions $a$, $b$, $c$ vaut
\[ d = \sqrt{a^2 + b^2 + c^2}. \]
C'est la généralisation spatiale du théorème de Pythagore : de même que la diagonale d'un rectangle vaut $\sqrt{a^2+b^2}$, celle du pavé vaut $\sqrt{a^2+b^2+c^2}$. Application concrète : pour vérifier qu'un carton de $60 \times 40 \times 30$~cm peut contenir un objet de $80$~cm, on calcule $\sqrt{60^2+40^2+30^2} = \sqrt{6100} \approx 78{,}1$~cm : l'objet ne rentre pas ! Les menuisiers de Douala utilisent ce calcul pour calibrer les caisses de transport.
\end{savaistu}

\begin{attention}[Préciser le plan avant d'appliquer Pythagore]
Erreur fréquente : écrire « d'après le théorème de Pythagore, $AG^2 = AB^2 + BG^2$ » sans vérifier que le triangle $ABG$ est rectangle. Or un triangle tracé sur un solide n'est rectangle que si l'on a \emph{prouvé} l'orthogonalité de deux de ses côtés. Rédige toujours ainsi : « $(CG) \perp (ABCD)$ et $(AC) \subset (ABCD)$, donc $(CG) \perp (AC)$ ; le triangle $ACG$ est rectangle en $C$ ; dans le \textbf{plan} $(ACG)$, le théorème de Pythagore donne... ». Sans cette justification, le calcul est faux même si le résultat est juste.
\end{attention}

\courssub{Distance d'un point à un plan, hauteur d'un tétraèdre}

Rappelons que la distance d'un point $S$ à un plan $(P)$ est la longueur $SO$, où $O$ est le projeté orthogonal de $S$ sur $(P)$. Pour une pyramide ou un tétraèdre, cette distance est la \emph{hauteur} du solide, et le volume s'écrit $V = \dfrac{1}{3} \times \mathcal{B} \times h$.

\begin{exoresolu}[Hauteur d'un tétraèdre régulier]
Soit $ABCD$ un tétraèdre régulier d'arête $a = 6$~cm. Soit $O$ le centre de gravité de la face $BCD$. On admet que $(AO) \perp (BCD)$. Calculer la hauteur $AO$ du tétraèdre.
\tcblower
\textbf{Étape 1 : longueur $BO$.} Le triangle $BCD$ est équilatéral de côté $6$. Sa médiane vaut $\dfrac{6\sqrt{3}}{2} = 3\sqrt{3}$. Le centre de gravité $O$ se situe aux deux tiers de la médiane issue de $B$ :
\[ BO = \frac{2}{3} \times 3\sqrt{3} = 2\sqrt{3}~\text{cm}. \]
\textbf{Étape 2 : Pythagore dans le plan $(ABO)$.} Comme $(AO) \perp (BCD)$ et $(BO) \subset (BCD)$, on a $(AO) \perp (BO)$ : le triangle $AOB$ est rectangle en $O$. Or $AB = 6$ (arête du tétraèdre). Donc :
\[ AO^2 = AB^2 - BO^2 = 36 - 12 = 24 \quad \text{d'où} \quad AO = 2\sqrt{6} \approx 4{,}9~\text{cm}. \]
\end{exoresolu}

\courssub{Étude complète guidée : la pyramide à base carrée}

Considérons la pyramide $SABCD$ à base carrée, de sommet $S$, telle que $(SO) \perp (ABCD)$ où $O$ est le centre du carré. C'est le modèle le plus fréquent au Probatoire.

\begin{popfigure}
\begin{center}
\begin{tikzpicture}[scale=1.1, line join=round]
\coordinate (A) at (0,0);
\coordinate (B) at (3.4,0);
\coordinate (C) at (4.4,1.6);
\coordinate (D) at (1,1.6);
\coordinate (O) at (2.2,0.8);
\coordinate (S) at (2.2,3.6);
% faces cachées
\draw[dashed, popmuted] (A)--(D) (D)--(C) (A)--(C) (B)--(D) (S)--(D);
% arêtes visibles
\draw[thick, popdark] (A)--(B)--(C);
\draw[thick, popblue] (S)--(A) (S)--(B) (S)--(C);
% hauteur
\draw[thick, poporange] (S)--(O);
% angle droit en O
\draw[poporange] (2.2,0.8) ++(0,0.22) -- ++(0.18,0.22) -- ++(0.18,0);
\fill (O) circle (1.2pt);
\foreach \p/\pos in {A/below left, B/below right, C/right, D/above left, S/above, O/below}
  \node[\pos, popdark, font=\small] at (\p) {$\p$};
\end{tikzpicture}
\end{center}
\legende{Pyramide $SABCD$ à base carrée : la hauteur $[SO]$ est orthogonale au plan de base, l'angle droit est marqué en $O$.}
\end{popfigure}

\begin{experience}[Démonstration guidée : une médiane orthogonale à un plan]
\textbf{Données.} $SABCD$ est une pyramide régulière : sa base $ABCD$ est un carré de centre $O$ et $(SO) \perp (ABCD)$. Soit $M$ le milieu de $[AB]$. On veut montrer que $(AB) \perp (SOM)$, puis en déduire que les plans $(SAB)$ et $(SOM)$ sont orthogonaux.

\textbf{Étape 1.} Dans le carré $ABCD$, $O$ est le milieu de $[AC]$ et $M$ le milieu de $[AB]$ ; la droite $(OM)$ est parallèle à $(BC)$ (droite des milieux dans le triangle $ABC$), donc $(OM) \perp (AB)$.

\textbf{Étape 2.} $(SO) \perp (ABCD)$ et $(AB) \subset (ABCD)$, donc $(SO) \perp (AB)$.

\textbf{Étape 3.} La droite $(AB)$ est orthogonale aux droites $(OM)$ et $(SO)$, qui sont deux droites \emph{sécantes} du plan $(SOM)$ (elles se coupent en $O$). Donc $(AB) \perp (SOM)$.

\textbf{Étape 4.} Le plan $(SAB)$ contient la droite $(AB)$ orthogonale au plan $(SOM)$ : donc $(SAB) \perp (SOM)$.

Retiens la mécanique : \emph{un angle droit dans la base} + \emph{la hauteur orthogonale à tout} = deux droites sécantes orthogonales, et tout s'enchaîne.
\end{experience}

\courssub{Arbre de décision : quelle méthode choisir ?}

Face à une question d'orthogonalité, le bon réflexe est de se demander d'abord \emph{ce que l'on doit prouver}, puis de remonter la chaîne des théorèmes. L'organigramme suivant résume la démarche.

\begin{popfigure}
\begin{center}
\begin{tikzpicture}[
  font=\small,
  quest/.style={rectangle, rounded corners, draw=popblue, fill=popblueL, text width=4.2cm, align=center, inner sep=4pt},
  act/.style={rectangle, rounded corners, draw=poporange, fill=poporangeL, text width=4.6cm, align=center, inner sep=4pt},
  lbl/.style={font=\footnotesize\itshape, popdark},
  arr/.style={-{Stealth[length=2.5mm]}, thick, popdark}
]
\node[quest] (q) at (0,0) {Que doit-on prouver ?};
\node[act, below left=1.1cm and 0.2cm of q] (a1) {Droites orthogonales : se ramener à un angle droit dans un plan (face, diagonales, triangle isocèle)};
\node[act, below=1.1cm of q] (a2) {Droite $\perp$ plan : deux droites \textbf{sécantes} du plan orthogonales à la droite};
\node[act, below right=1.1cm and 0.2cm of q] (a3) {Plans orthogonaux : trouver une droite de l'un orthogonale à l'autre};
\node[act, below=1.4cm of a3, draw=popgreen, fill=popgreenL] (a4) {Se ramener au cas « droite $\perp$ plan »};
\draw[arr] (q) -- (a1) node[lbl, midway, above left] {droites};
\draw[arr] (q) -- (a2) node[lbl, midway, right] {droite/plan};
\draw[arr] (q) -- (a3) node[lbl, midway, above right] {plans};
\draw[arr] (a3) -- (a4);
\draw[arr, dashed, popgreen] (a4.west) -| (a2.south);
\end{tikzpicture}
\end{center}
\legende{Arbre de décision : la preuve de l'orthogonalité de deux plans se ramène toujours à celle d'une droite orthogonale à un plan.}
\end{popfigure}

\courssub{Exercice de synthèse type Probatoire : le cube}

\begin{exoresolu}[Orthogonalités dans le cube $ABCDEFGH$]
Soit $ABCDEFGH$ un cube d'arête $a$ (face $ABCD$ en bas, $E$ au-dessus de $A$). Soit $O$ le centre de la face $ABCD$.
\begin{enumerate}
\item Montrer que $(BD) \perp (ACG)$.
\item En déduire que $(BD) \perp (AG)$.
\item Montrer que les plans $(ABCD)$ et $(ACGE)$ sont orthogonaux.
\item Calculer $AG$ en fonction de $a$.
\end{enumerate}
\tcblower
\textbf{1.} La face $ABCD$ est un carré, donc ses diagonales sont perpendiculaires : $(BD) \perp (AC)$. D'autre part $(AE) \perp (ABCD)$... utilisons plutôt $(CG)$ : l'arête $(CG)$ est orthogonale au plan $(ABCD)$, et $(BD) \subset (ABCD)$, donc $(CG) \perp (BD)$. Ainsi $(BD)$ est orthogonale aux droites $(AC)$ et $(CG)$, sécantes en $C$ dans le plan $(ACG)$. Donc $(BD) \perp (ACG)$.

\textbf{2.} La droite $(AG)$ est contenue dans le plan $(ACG)$ (car $A$ et $G$ appartiennent à ce plan). Or $(BD) \perp (ACG)$ : toute droite du plan $(ACG)$ est donc orthogonale à $(BD)$. En particulier $(BD) \perp (AG)$.

\textbf{3.} Le plan $(ACGE)$ contient l'arête $(AE)$, orthogonale au plan $(ABCD)$. Un plan contenant une droite orthogonale à un autre plan lui est orthogonal : $(ACGE) \perp (ABCD)$.

\textbf{4.} Dans le carré $ABCD$, triangle $ABC$ rectangle en $B$ : $AC^2 = AB^2 + BC^2 = 2a^2$. Comme $(CG) \perp (ABCD)$, le triangle $ACG$ est rectangle en $C$ ; dans le plan $(ACG)$ :
\[ AG^2 = AC^2 + CG^2 = 2a^2 + a^2 = 3a^2 \quad \text{donc} \quad AG = a\sqrt{3}. \]
\end{exoresolu}

\begin{exercice}[Pour t'entraîner]
Soit $SABC$ un tétraèdre tel que $(SA) \perp (ABC)$ et le triangle $ABC$ rectangle en $B$.
\begin{enumerate}
\item Montrer que $(BC) \perp (SAB)$.
\item En déduire que les plans $(SBC)$ et $(SAB)$ sont orthogonaux.
\item On donne $SA = 5$~cm, $AB = 6$~cm, $BC = 8$~cm. Calculer $SB$, $SC$ et le volume du tétraèdre.
\end{enumerate}
\end{exercice}

\begin{corrige}[Exercice : tétraèdre $SABC$]
\textbf{1.} $(SA) \perp (ABC)$ et $(BC) \subset (ABC)$, donc $(SA) \perp (BC)$. De plus $(AB) \perp (BC)$ car $ABC$ est rectangle en $B$. La droite $(BC)$ est orthogonale aux droites $(SA)$ et $(AB)$, sécantes en $A$ dans le plan $(SAB)$ : donc $(BC) \perp (SAB)$.

\textbf{2.} Le plan $(SBC)$ contient la droite $(BC)$ orthogonale au plan $(SAB)$, donc $(SBC) \perp (SAB)$.

\textbf{3.} $(SA) \perp (AB)$, triangle $SAB$ rectangle en $A$ : $SB^2 = 25 + 36 = 61$, $SB = \sqrt{61} \approx 7{,}8$~cm. $(SA) \perp (AC)$ ; d'abord $AC^2 = AB^2 + BC^2 = 36 + 64 = 100$, $AC = 10$~cm ; puis $SC^2 = SA^2 + AC^2 = 25 + 100 = 125$, $SC = 5\sqrt{5} \approx 11{,}2$~cm. Volume : $V = \dfrac{1}{3} \times \mathcal{A}_{ABC} \times SA = \dfrac{1}{3} \times \dfrac{6 \times 8}{2} \times 5 = 40$~cm$^3$.
\end{corrige}

\courssub{Ouverture vers la Terminale}

\begin{savaistu}[Le produit scalaire changera tout — section d'ouverture]
Toutes les démonstrations de cette leçon reposent sur des enchaînements géométriques parfois délicats : trouver le bon plan, la bonne paire de droites sécantes... En Terminale C et D, tu découvriras le \cle{produit scalaire} dans l'espace et la \cle{géométrie analytique}. Avec ces outils, prouver que $(BD) \perp (AG)$ dans le cube deviendra un calcul de trois lignes : $\overrightarrow{BD} \cdot \overrightarrow{AG} = 0$. De même, la distance d'un point à un plan s'obtiendra par une formule directe. Les résultats que tu as démontrés ici « à la main » seront alors retrouvés presque instantanément — mais c'est la vision géométrique acquise maintenant qui te permettra d'interpréter ces calculs. Ton travail d'aujourd'hui est donc un investissement, pas une corvée !
\end{savaistu}

\begin{aretenir}[L'essentiel de la section]
\begin{itemize}
\item Droite $\perp$ plan : exhiber \textbf{deux droites sécantes} du plan, orthogonales à la droite.
\item Plans orthogonaux : trouver une \textbf{droite de l'un orthogonale à l'autre}.
\item Tout calcul de longueur se fait dans un \textbf{plan identifié}, après avoir justifié l'angle droit.
\item Grande diagonale du pavé : $d = \sqrt{a^2+b^2+c^2}$ ; diagonale du cube : $a\sqrt{3}$.
\item Le pavé droit est un repère déguisé : exploite systématiquement les angles droits de ses faces.
\end{itemize}
\end{aretenir}

\newpage

\coursec{Activité d'intégration}

\coursec{Leçon 18 : Orthogonalité dans l'espace}

\courssub{Objectifs de l'activité}

À l'issue de cette activité d'intégration, tu seras capable de :

\begin{itemize}
\item modéliser une situation concrète de la vie quotidienne (un atelier de menuiserie) par un solide géométrique de l'espace ;
\item démontrer qu'une droite est \cle{orthogonale à un plan} en utilisant le théorème fondamental (orthogonalité à deux droites sécantes du plan) ;
\item distinguer et justifier les notions de droites \cle{perpendiculaires} et de droites \cle{orthogonales} non sécantes ;
\item démontrer que deux plans sont \cle{orthogonaux} en exhibant une droite de l'un orthogonale à l'autre ;
\item prouver que deux droites sont \cle{non coplanaires} par un raisonnement par l'absurde ;
\item calculer une \cle{distance} dans l'espace grâce au théorème de Pythagore utilisé en cascade ;
\item rédiger une démonstration complète et rigoureuse, puis évaluer la production d'un autre groupe.
\end{itemize}

\courssub{Situation}

Moussa est menuisier à Nkongsamba, dans la région du Littoral. Il vient d'aménager son nouvel atelier, une pièce en forme de \cle{pavé droit} que l'on modélise par le parallélépipède rectangle $ABCDEFGH$ de dimensions :
\[ AB = \SI{4}{m}, \qquad AD = \SI{3}{m}, \qquad AE = \SI{2,5}{m}. \]

Le sol de l'atelier est le plan $(ABC)$, le mur de devant est le plan $(ABF)$, et les arêtes $(AE)$, $(BF)$, $(CG)$, $(DH)$ sont les quatre arêtes verticales (les montants d'angle de la pièce).

Pour ranger ses planches, Moussa veut fixer une étagère le long du montant vertical $(AE)$, puis tendre un câble métallique entre deux coins opposés de l'atelier, $A$ et $G$, pour y suspendre son petit poste radio. Son apprenti, Junior, a aussi tendu deux fils : l'un le long du segment $[AC]$ (diagonale du sol, pour guider la découpe), l'autre le long du segment $[HB]$ (pour accrocher un éclairage). Junior affirme : « Ces deux fils se coupent forcément quelque part ! » Moussa n'en est pas si sûr\ldots

Avant de percer le moindre trou, Moussa veut vérifier mathématiquement la solidité et la cohérence de son installation. Vous êtes les membres de son équipe d'apprentis : à vous de répondre à toutes ses questions !

\begin{popfigure}
\begin{center}
\begin{tikzpicture}[scale=1.1, line join=round]
\coordinate (A) at (0,0);
\coordinate (B) at (4,0);
\coordinate (C) at (5.4,1.4);
\coordinate (D) at (1.4,1.4);
\coordinate (E) at (0,2.5);
\coordinate (F) at (4,2.5);
\coordinate (G) at (5.4,3.9);
\coordinate (H) at (1.4,3.9);
% faces cachées en pointillés
\draw[dashed, popmuted] (A) -- (D) -- (C);
\draw[dashed, popmuted] (D) -- (H);
% arêtes visibles
\draw[thick, popdark] (A) -- (B) -- (C) -- (G) -- (H) -- (E) -- cycle;
\draw[thick, popdark] (B) -- (F) -- (G);
\draw[thick, popdark] (E) -- (F);
% éléments de la situation
\draw[very thick, popblue] (A) -- (E);
\draw[very thick, poporange] (A) -- (G);
\draw[thick, popgreen, dashed] (A) -- (C);
\draw[thick, poppurple] (H) -- (B);
% points
\foreach \p/\pos in {A/below left, B/below, C/right, D/above left, E/left, F/right, G/above right, H/above}
  \fill[popink] (\p) circle (1.3pt) node[\pos, inner sep=2pt] {\small $\p$};
\node[popblue, left] at (-0.15,1.25) {\small montant};
\node[poporange, above right, rotate=33] at (2.7,2.1) {\small câble};
\end{tikzpicture}
\end{center}
\legende{Modélisation de l'atelier de Moussa : pavé droit $ABCDEFGH$, montant $(AE)$ en bleu, câble $[AG]$ en orange, fils $[AC]$ et $[HB]$ en vert et violet.}
\end{popfigure}

\courssub{Consignes}

Travail en groupes de quatre élèves. Chaque groupe rédige sur une feuille une démonstration complète et soignée de chaque question, puis deux groupes échangent leurs copies pour une correction croisée avant la synthèse collective au tableau.

\begin{enumerate}
\item \textbf{Le montant est-il bien vertical ?} Démontrer que la droite $(AE)$ est orthogonale au plan du sol $(ABC)$. On précisera le théorème utilisé et toutes ses hypothèses.
\item \textbf{Des arêtes orthogonales sans contact.} On considère les arêtes $(AD)$ du sol et $(BF)$ du mur de devant. Démontrer que ces deux droites sont orthogonales. Sont-elles perpendiculaires ? Justifier la différence.
\item \textbf{Le mur est-il bien droit ?} Démontrer que le plan du mur de devant $(ABF)$ est orthogonal au plan du sol $(ABC)$.
\item \textbf{Le débat entre Moussa et Junior.} Démontrer que les droites $(AC)$ et $(HB)$ ne sont pas coplanaires. Que penser alors de l'affirmation de Junior ?
\item \textbf{La longueur du câble.} Calculer la longueur exacte du câble $[AG]$ tendu entre les deux coins opposés de l'atelier, puis en donner une valeur approchée au centimètre près. Le rouleau de câble de $\SI{6}{m}$ acheté par Moussa au marché suffira-t-il ?
\item \textbf{Pour aller plus loin.} Soit $I$ le milieu de $[AB]$. Démontrer que la droite $(HG)$ est orthogonale au plan $(BCG)$, puis en déduire que $(HG)$ est orthogonale à toute droite du plan $(BCG)$ passant par $G$.
\end{enumerate}

\courssub{Ressources à mobiliser}

\begin{aretenir}[Boîte à outils de la leçon]
\begin{itemize}
\item \textbf{Droite orthogonale à un plan :} une droite $(d)$ est orthogonale à un plan $(P)$ si et seulement si elle est orthogonale à \emph{deux droites sécantes} de $(P)$. Elle est alors orthogonale à \emph{toute} droite de $(P)$.
\item \textbf{Dans un pavé droit,} toute arête est orthogonale aux faces qu'elle ne traverse pas parallèlement ; en particulier chaque arête verticale est orthogonale aux faces horizontales.
\item \textbf{Plans orthogonaux :} si une droite $(d)$ est orthogonale à un plan $(P)$, alors tout plan contenant $(d)$ est orthogonal à $(P)$.
\item \textbf{Droites non coplanaires :} deux droites de l'espace sont soit coplanaires (sécantes ou parallèles), soit non coplanaires. Pour prouver la non-coplanarité, on raisonne souvent par l'absurde.
\item \textbf{Théorème de Pythagore en cascade :} si $(d) \perp (P)$ en $H$, alors pour tout point $M$ de $(P)$, le triangle formé est rectangle en $H$.
\end{itemize}
\end{aretenir}

\courssub{Démarche et corrigé commenté}

\begin{exoresolu}[L'étagère murale de l'atelier de menuiserie]
On reprend les six questions de la situation : pavé droit $ABCDEFGH$ avec $AB = \SI{4}{m}$, $AD = \SI{3}{m}$, $AE = \SI{2,5}{m}$.
\tcblower
\textbf{Question 1 : $(AE)$ orthogonale au sol $(ABC)$.}

$ABCDEFGH$ est un pavé droit, donc la face $ABCD$ est un rectangle et les arêtes latérales sont perpendiculaires aux faces de base. Ainsi :
\begin{itemize}
\item $(AE) \perp (AB)$ car l'angle $\widehat{EAB}$ est un angle droit du rectangle $ABFE$ ;
\item $(AE) \perp (AD)$ car l'angle $\widehat{EAD}$ est un angle droit du rectangle $ADHE$.
\end{itemize}
Or $(AB)$ et $(AD)$ sont deux droites du plan $(ABC)$, \emph{sécantes en $A$}. D'après le théorème fondamental (caractérisation de l'orthogonalité droite-plan), $(AE)$ est orthogonale au plan $(ABC)$. Le montant est bien perpendiculaire au sol : l'étagère sera stable. $\square$

\medskip
\textbf{Question 2 : $(AD)$ et $(BF)$ orthogonales mais non perpendiculaires.}

D'une part, $(BF)$ est parallèle à $(AE)$ (arêtes verticales d'un pavé droit). Or $(AE) \perp (ABC)$ d'après la question 1, donc $(BF) \perp (ABC)$ (propriété : si deux droites sont parallèles et si l'une est orthogonale à un plan, l'autre l'est aussi). Comme $(AD)$ est une droite du plan $(ABC)$, on a $(BF) \perp (AD)$ : les droites $(AD)$ et $(BF)$ sont \cle{orthogonales}.

D'autre part, $(AD)$ passe par $A$ et $D$ (arrière du sol) tandis que $(BF)$ passe par $B$ et $F$ (coin avant droit) : ces droites ne se coupent pas. Elles sont donc orthogonales \emph{sans être perpendiculaires}, car la perpendicularité exige l'intersection. $\square$

\medskip
\textbf{Question 3 : plans $(ABF)$ et $(ABC)$ orthogonaux.}

On a montré à la question 1 que $(AE) \perp (ABC)$. Or le plan $(ABF)$ contient la droite $(AE)$ (puisque $A \in (ABF)$ et $E \in (ABF)$ car $E \in (BF)$ ? Non, $E \notin (BF)$ mais $A \in (ABF)$ et $F \in (ABF)$ et $B \in (ABF)$. Le plan $(ABF)$ contient la droite $(AE)$ car $A$ et $E$ appartiennent au plan $(ABF)$ : $A$ est sur l'arête $AB$, $E$ est sur l'arête $AE$ qui est dans le plan $ABF$ ? Attendez : Le plan $(ABF)$ est défini par les points $A, B, F$. La droite $(AE)$ passe par $A$ et $E$. Est-ce que $E \in (ABF)$ ? Dans un pavé droit, $A, B, F, E$ sont les sommets de la face latérale $ABFE$, qui est un rectangle. Donc $A, B, F, E$ sont coplanaires. Ainsi $E \in (ABF)$. Donc $(AE) \subset (ABF)$.
D'après le théorème des plans orthogonaux : \emph{si un plan contient une droite orthogonale à un autre plan, alors les deux plans sont orthogonaux}. Donc $(ABF) \perp (ABC)$. Le mur est bien dressé perpendiculairement au sol. $\square$

\medskip
\textbf{Question 4 : $(AC)$ et $(HB)$ non coplanaires.}

Raisonnons par l'absurde : supposons que $(AC)$ et $(HB)$ soient coplanaires, c'est-à-dire contenues dans un même plan $(Q)$. Alors $A$, $C$, $H$, $B$ appartiendraient à $(Q)$. Or $A$, $B$, $C$ sont trois points non alignés du sol, donc $(Q)$ serait le plan $(ABC)$. Mais $H \notin (ABC)$ car $H$ est à la hauteur $\SI{2,5}{m}$ au-dessus de $D$ (arête verticale $(DH)$). Contradiction.

Donc $(AC)$ et $(HB)$ ne sont pas coplanaires : elles ne peuvent être ni sécantes ni parallèles. Junior a tort : les deux fils ne se couperont jamais, même s'ils semblent se croiser vus de la porte (effet de perspective). $\square$

\medskip
\textbf{Question 5 : longueur du câble $[AG]$.}

\emph{Étape 1 : diagonale du sol.} Le triangle $ABC$ est rectangle en $B$ (face rectangulaire $ABCD$). D'après le théorème de Pythagore :
\[ AC^2 = AB^2 + BC^2 = 4^2 + 3^2 = 16 + 9 = 25, \quad \text{donc } AC = \SI{5}{m}. \]

\emph{Étape 2 : triangle rectangle dans l'espace.} Comme $(CG)$ est une arête verticale, $(CG) \perp (ABC)$, donc $(CG) \perp (AC)$ : le triangle $ACG$ est rectangle en $C$. On applique à nouveau Pythagore :
\begin{align*}
AG^2 &= AC^2 + CG^2 \\
&= 25 + (\SI{2,5}{m})^2 \\
&= 25 + 6,25 = 31,25.
\end{align*}
Donc $AG = \sqrt{31,25} = \frac{5\sqrt{5}}{2} \approx \SI{5,59}{m}$, soit environ $\SI{5}{m}\,\SI{59}{cm}$.

Le rouleau de $\SI{6}{m}$ suffit : il restera même environ $\SI{41}{cm}$ de marge pour les nœuds de fixation. $\square$

\medskip
\textbf{Question 6 : pour aller plus loin.}

$(HG)$ est parallèle à $(DC)$ et à $(AB)$ (arêtes opposées d'un pavé droit). Or $(AB) \perp (BC)$ (rectangle $ABCD$) et $(AB) \perp (BF)$ (arête perpendiculaire à la face latérale $BCGF$), et $(BC)$ et $(BF)$ sont deux droites sécantes du plan $(BCG)$ en $B$ (la face $BCGF$ est le plan $(BCG)$). Donc $(AB) \perp (BCG)$, et par parallélisme $(HG) \perp (BCG)$. Par définition de l'orthogonalité droite-plan, $(HG)$ est alors orthogonale à \emph{toute} droite du plan $(BCG)$, en particulier à toute droite de ce plan passant par $G$, comme $(GC)$ ou $(GB)$. $\square$
\end{exoresolu}

\begin{attention}[Pièges relevés lors de la correction croisée]
\begin{itemize}
\item Pour prouver qu'une droite est orthogonale à un plan, il faut \textbf{deux} droites sécantes de ce plan : une seule ne suffit pas !
\item Ne pas confondre « orthogonales » et « perpendiculaires » : des droites orthogonales ne se rencontrent pas nécessairement.
\item Dans le calcul de $AG$, on ne peut pas additionner directement $AB + AD + AE$ : il faut deux applications successives du théorème de Pythagore, en justifiant chaque angle droit.
\item Deux droites qui semblent se croiser sur un dessin en perspective ne sont pas forcément sécantes dans l'espace : seule la démonstration tranche.
\item Attention au vocabulaire : « $(AE)$ est orthogonale au plan $(ABC)$ » et non « $(AE)$ est perpendiculaire au plan $(ABC)$ » (on réserve « perpendiculaire » pour les droites sécantes).
\end{itemize}
\end{attention}

\courssub{Bilan de l'activité}

Cette activité a permis de mobiliser, sur une même situation concrète et authentique, l'ensemble des savoir-faire de la leçon :

\begin{itemize}
\item la \cle{caractérisation fondamentale} de l'orthogonalité d'une droite et d'un plan (questions 1 et 6), pierre angulaire de tout le chapitre ;
\item la distinction subtile entre droites \cle{orthogonales} et droites \cle{perpendiculaires} (question 2) ;
\item le théorème des \cle{plans orthogonaux}, appliqué au couple mur-sol (question 3) ;
\item le raisonnement par l'absurde pour établir la \cle{non-coplanarité} de deux droites (question 4) ;
\item le calcul de \cle{distance} dans l'espace par double application du théorème de Pythagore (question 5), qui donne au passage la formule de la diagonale d'un pavé : $AG^2 = AB^2 + AD^2 + AE^2$.
\end{itemize}

La correction croisée a en outre développé des compétences transversales : lire et évaluer la démonstration d'un camarade, repérer une hypothèse manquante, exiger la justification de chaque angle droit. Retiens l'essentiel : dans l'espace, \emph{rien ne se devine, tout se démontre} — et c'est précisément cette rigueur qui garantit que l'étagère de Moussa tiendra droit !

\begin{savaistu}[Le saviez-vous ?]
La formule $d^2 = a^2 + b^2 + c^2$ de la grande diagonale du pavé droit est la généralisation à trois dimensions du théorème de Pythagore. C'est elle qu'utilisent les techniciens de MTN ou de CAMTEL pour estimer la longueur de câble nécessaire dans un bâtiment, et les maçons camerounais pour vérifier l'équerrage d'une pièce : en mesurant les deux diagonales du sol, elles doivent être égales si les murs forment bien un rectangle !
\end{savaistu}

\newpage

\coursec{Méthodes et exercices résolus}

\coursec{Orthogonalité dans l'espace}

\courssub{Positions relatives et coplanarité de deux droites}

\begin{aretenir}[Droites coplanaires et droites gauches]
Dans l'espace, deux droites $(d_1)$ et $(d_2)$ sont \cle{coplanaires} s'il existe un plan qui les contient toutes les deux. Cela se produit dans deux cas exclusifs :
\begin{itemize}
\item elles sont \textbf{sécantes} (elles ont un point commun) ;
\item elles sont \textbf{parallèles} (elles ont même direction).
\end{itemize}
Si elles ne sont ni sécantes ni parallèles, on dit qu'elles sont \cle{non coplanaires} (ou \emph{gauches}).
\end{aretenir}

\begin{methode}[Coplanarité de deux droites]
Pour déterminer si deux droites $(d_1)$ et $(d_2)$ de l'espace sont coplanaires ou non, on procède ainsi :
\begin{enumerate}
\item \textbf{Repérer les supports.} Identifier deux points de $(d_1)$ et deux points de $(d_2)$ (souvent des arêtes ou des diagonales d'un solide : cube, tétraèdre, pyramide, pavé droit).
\item \textbf{Tester les cas favorables.} Chercher si les deux droites sont \emph{sécantes} (point commun) ou \emph{parallèles}. Dans l'un ou l'autre cas, elles sont coplanaires.
\item \textbf{Pour prouver la non-coplanarité}, on raisonne souvent par l'absurde ou par l'appartenance à des plans :
\begin{itemize}
\item on exhibe un plan $(P)$ contenant $(d_1)$ ;
\item on montre qu'un point de $(d_2)$ n'appartient pas à $(P)$ et que $(d_2)$ n'est pas parallèle à $(P)$ ;
\item ou bien : si les droites étaient coplanaires, quatre points (deux sur chaque droite) seraient coplanaires, ce qui contredit une hypothèse (par exemple un tétraèdre, dont les quatre sommets sont non coplanaires).
\end{itemize}
\item \textbf{Conclure proprement} : « les droites sont coplanaires » ou « les droites sont non coplanaires (gauches) ».
\end{enumerate}
\textbf{Réflexe} : dans un tétraèdre $ABCD$, deux arêtes \emph{opposées} (comme $(AB)$ et $(CD)$) sont toujours non coplanaires, car sinon $A$, $B$, $C$, $D$ seraient coplanaires.
\end{methode}

\begin{popfigure}
\begin{tikzpicture}[scale=1.2]
\coordinate (A) at (0,0);
\coordinate (B) at (3,0);
\coordinate (C) at (1.5,2.6);
\coordinate (D) at (1.5,0.8);
\draw[popdark, thick] (A) -- (B) -- (C) -- cycle;
\draw[popdark, thick] (A) -- (D) (B) -- (D) (C) -- (D);
\draw[popdark, dashed] (A) -- (B); % Base hidden? No, standard tetrahedron.
% Let's draw a nice tetrahedron
\coordinate (A) at (0,0);
\coordinate (B) at (3,0);
\coordinate (C) at (1.5,2.6);
\coordinate (D) at (1.5,0.8);
\draw[popblue, thick] (A) -- (B) -- (C) -- cycle;
\draw[popblue, thick] (A) -- (D) (B) -- (D) (C) -- (D);
\foreach \p/\pos in {A/left, B/right, C/above, D/below} \node[\pos] at (\p) {$\p$};
\node at (0.75,1.3) {$I$};
\node at (1.5,1.7) {$J$};
\draw[poporange, thick, dashed] (0.75,1.3) -- (1.5,1.7);
\end{tikzpicture}
\legende{Tétraèdre $ABCD$ : arêtes opposées $(AB)$ et $(CD)$ non coplanaires ; $(IJ)$ et $(AB)$ coplanaires}
\end{popfigure}

\begin{exoresolu}[Coplanarité dans un tétraèdre]
Soit $ABCD$ un tétraèdre (les quatre points $A$, $B$, $C$, $D$ sont non coplanaires). On note $I$ le milieu de $[AB]$ et $J$ le milieu de $[CD]$.
\begin{enumerate}
\item Démontrer que les droites $(AB)$ et $(CD)$ sont non coplanaires.
\item Les droites $(IJ)$ et $(AB)$ sont-elles coplanaires ? Justifier.
\end{enumerate}
\tcblower
\textbf{1.} Raisonnons par l'absurde. Supposons que $(AB)$ et $(CD)$ soient coplanaires : il existerait alors un plan $(P)$ contenant à la fois $(AB)$ et $(CD)$. Or une droite est entièrement contenue dans un plan dès que deux de ses points appartiennent à ce plan. Donc $A, B \in (P)$ (car $A, B \in (AB)$) et $C, D \in (P)$ (car $C, D \in (CD)$). Les quatre points $A$, $B$, $C$, $D$ seraient coplanaires, ce qui \textbf{contredit l'hypothèse} « $ABCD$ est un tétraèdre ».

\textbf{Conclusion} : les droites $(AB)$ et $(CD)$ sont \cle{non coplanaires} (on dit aussi qu'elles sont \emph{gauches}).

\medskip
\textbf{2.} Les droites $(IJ)$ et $(AB)$ ont le point $I$ en commun : en effet $I \in [AB]$ donc $I \in (AB)$, et $I \in (IJ)$ par définition. Elles sont donc \textbf{sécantes en $I$} (elles ne sont pas confondues car $J \notin (AB)$ : si $J \in (AB)$, alors les droites $(AB)$ et $(CD)$ se rencontreraient en $J$ puisque $J \in [CD]$, ce qui contredit la question 1 qui affirme qu'elles sont gauches).

Or deux droites sécantes sont toujours coplanaires : elles déterminent un unique plan.

\textbf{Conclusion} : les droites $(IJ)$ et $(AB)$ sont \cle{coplanaires}, sécantes en $I$.
\end{exoresolu}

\courssub{Droites orthogonales et droites perpendiculaires}

\begin{aretenir}[Orthogonalité et perpendicularité]
Deux droites $(d_1)$ et $(d_2)$ sont \cle{orthogonales} si leurs directions sont orthogonales, c'est-à-dire si les parallèles menées à ces droites par un même point quelconque sont perpendiculaires.
Elles sont \cle{perpendiculaires} si, \textbf{en plus}, elles sont \emph{sécantes} (elles se coupent en formant un angle droit).
\begin{center}
\textbf{Perpendiculaires $\implies$ Orthogonales \quad mais \quad Orthogonales $\nRightarrow$ Perpendiculaires}
\end{center}
\end{aretenir}

\begin{methode}[Orthogonalité de deux droites]
Pour démontrer que deux droites sont orthogonales (ou perpendiculaires) :
\begin{enumerate}
\item \textbf{Choisir l'outil adapté à la situation} :
\begin{itemize}
\item \emph{Géométrie plane} : si les deux droites sont coplanaires (souvent dans une face d'un solide), utiliser les propriétés du plan (angles d'un carré, d'un rectangle, théorème de Pythagore, médiane d'un triangle isocèle, hauteur d'un triangle équilatéral).
\item \emph{Via un plan} : montrer que l'une des droites est orthogonale à un plan contenant l'autre (voir fiche 3) — c'est la méthode la plus puissante en géométrie dans l'espace.
\item \emph{Produit scalaire} (si un repère ou des vecteurs sont disponibles) : $(d_1) \perp (d_2) \iff \vec{u}_1 \cdot \vec{u}_2 = 0$ où $\vec{u}_1, \vec{u}_2$ sont des vecteurs directeurs.
\end{itemize}
\item \textbf{Distinguer les vocabulaires} : les droites sont \emph{perpendiculaires} si elles sont orthogonales \textbf{et} sécantes ; elles peuvent être orthogonales sans se rencontrer (droites gauches).
\item \textbf{Rédiger la conclusion} en précisant : « orthogonales » (toujours vrai si les directions le sont) ou « perpendiculaires » (seulement si on a aussi prouvé l'intersection).
\end{enumerate}
\textbf{Réflexe dans un cube} : une arête est orthogonale à toute droite de la face qu'elle ne perce pas parallèlement ; les diagonales d'une face carrée sont perpendiculaires entre elles.
\end{methode}

\begin{popfigure}
\begin{tikzpicture}[scale=1.2]
\coordinate (A) at (0,0);
\coordinate (B) at (3,0);
\coordinate (C) at (3,3);
\coordinate (D) at (0,3);
\coordinate (E) at (0.5,3.5);
\coordinate (F) at (3.5,3.5);
\coordinate (G) at (3.5,6.5);
\coordinate (H) at (0.5,6.5);
% Base
\draw[popdark, thick] (A) -- (B) -- (C) -- (D) -- cycle;
% Top
\draw[popdark, thick] (E) -- (F) -- (G) -- (H) -- cycle;
% Vertical edges
\draw[popdark, thick] (A) -- (E) (B) -- (F) (C) -- (G) (D) -- (H);
% Hidden edges dashed
\draw[popdark, dashed] (D) -- (H) -- (E) -- (A); % Back left
\draw[popdark, dashed] (C) -- (D); % Back bottom? No, D-C is visible if viewed from front-right. Let's assume standard view.
% Diagonals on base
\draw[poporange, thick] (A) -- (C) (B) -- (D);
\node at (1.5,1.5) {$O$};
% Labels
\foreach \p/\pos in {A/below left, B/below right, C/above right, D/above left, E/above left, F/above right, G/above right, H/above left}
  \node[\pos] at (\p) {$\p$};
% Arrow for AE
\draw[popgreen, thick, ->] (0.25,3.25) -- (0.25,4.5) node[left] {$(AE)$};
\end{tikzpicture}
\legende{Cube $ABCDEFGH$ : $(AC)\perp(BD)$ dans le plan $(ABCD)$ ; $(AE)\perp(ABCD)$ donc $(AE)\perp(BD)$}
\end{popfigure}

\begin{exoresolu}[Orthogonalité dans un cube]
Un menuisier de Yaoundé confectionne un coffre cubique $ABCDEFGH$ d'arête $a$ (avec $ABCD$ la face du dessous et $EFGH$ la face du dessus, $E$ au-dessus de $A$).
\begin{enumerate}
\item Démontrer que les droites $(AC)$ et $(BD)$ sont perpendiculaires.
\item Démontrer que les droites $(AE)$ et $(BD)$ sont orthogonales. Sont-elles perpendiculaires ?
\end{enumerate}
\tcblower
\textbf{1.} Les droites $(AC)$ et $(BD)$ sont les deux diagonales de la face $ABCD$, qui est un \textbf{carré}. Or les diagonales d'un carré sont perpendiculaires et se coupent en leur milieu. Donc $(AC)$ et $(BD)$ sont sécantes et forment un angle droit : elles sont \cle{perpendiculaires}.

\medskip
\textbf{2.} Montrons d'abord que $(AE)$ est orthogonale au plan $(ABCD)$. L'arête $[AE]$ est perpendiculaire aux deux droites $(AB)$ et $(AD)$ de la face $ABCD$ (angles droits du carré aux sommets des arêtes du cube). Or $(AB)$ et $(AD)$ sont deux droites \emph{sécantes} du plan $(ABCD)$. D'après le théorème de la droite orthogonale à un plan (fiche 3), $(AE) \perp (ABCD)$.

Toute droite orthogonale à un plan est orthogonale à \emph{toute} droite de ce plan. Comme $(BD) \subset (ABCD)$, on en déduit que $(AE)$ et $(BD)$ sont \cle{orthogonales}.

Enfin, $(AE)$ et $(BD)$ ne se rencontrent pas : $(AE)$ coupe le plan $(ABCD)$ en $A$ seul, et $A \notin (BD)$ (les diagonales du carré ne passent pas par les sommets autres que leurs extrémités). Elles sont donc orthogonales \textbf{mais non perpendiculaires}.

\textbf{Conclusion} : $(AE) \perp (BD)$ en direction (droites orthogonales), sans être sécantes.
\end{exoresolu}

\courssub{Droite orthogonale à un plan}

\begin{aretenir}[Théorème fondamental : Droite orthogonale à un plan]
Une droite $(d)$ est orthogonale à un plan $(P)$ \textbf{si et seulement si} elle est orthogonale à \cle{deux droites sécantes} de $(P)$.
\begin{itemize}
\item Si $(d) \perp (P)$, alors $(d)$ est orthogonale à \emph{toute} droite de $(P)$.
\item La réciproque est le critère de démonstration : il suffit de trouver deux droites sécantes dans $(P)$ orthogonales à $(d)$.
\end{itemize}
\end{aretenir}

\begin{methode}[Droite orthogonale à un plan]
Démarche pour prouver $(d) \perp (P)$ :
\begin{enumerate}
\item \textbf{Identifier le plan $(P)$} et y choisir deux droites \emph{sécantes} $(\Delta_1)$ et $(\Delta_2)$ (souvent deux côtés d'une face, ou un côté et une diagonale, ou deux médianes).
\item \textbf{Prouver séparément} $(d) \perp (\Delta_1)$ et $(d) \perp (\Delta_2)$, en justifiant chaque orthogonalité (face carrée/rectangulaire, triangle isocèle, théorème de Pythagore, droite déjà connue orthogonale à un plan contenant $(\Delta_i)$).
\item \textbf{Vérifier les conditions} : $(\Delta_1)$ et $(\Delta_2)$ sont bien \emph{dans} $(P)$ et bien \emph{sécantes} (non parallèles !).
\item \textbf{Conclure} par le théorème : $(d) \perp (P)$.
\item \textbf{Exploiter} : $(d)$ est alors orthogonale à \emph{toute} droite de $(P)$, ce qui fournit gratuitement de nouvelles orthogonalités.
\end{enumerate}
\textbf{Variante utile} : si deux plans sont parallèles, toute droite orthogonale à l'un est orthogonale à l'autre.
\end{methode}

\begin{popfigure}
\begin{tikzpicture}[scale=1.2]
\coordinate (A) at (-2,-1.5);
\coordinate (B) at (2,-1.5);
\coordinate (C) at (2,1.5);
\coordinate (D) at (-2,1.5);
\coordinate (O) at (0,0);
\coordinate (S) at (0,3);
\draw[popdark, thick] (A) -- (B) -- (C) -- (D) -- cycle;
\draw[popdark, thick] (S) -- (A) (S) -- (B) (S) -- (C) (S) -- (D);
\draw[popdark, dashed] (A) -- (C) (B) -- (D);
\draw[popblue, thick, ->] (S) -- (O) node[midway, right] {$SO$};
\foreach \p/\pos in {A/left, B/right, C/right, D/left, O/below, S/above} \node[\pos] at (\p) {$\p$};
% Right angle marks
\draw[popgreen] (0,-0.15) -- (0.15,-0.15) -- (0.15,0);
\draw[popgreen] (-0.15,0) -- (-0.15,0.15) -- (0,0.15);
\end{tikzpicture}
\legende{Pyramide régulière $SABCD$ : $(SO)$ hauteur, orthogonale à la base $(ABCD)$}
\end{popfigure}

\begin{exoresolu}[Hauteur d'une pyramide régulière]
Soit $SABCD$ une pyramide régulière de sommet $S$, dont la base $ABCD$ est un carré de centre $O$ (c'est le modèle du toit d'une case traditionnelle). On donne $AB = \SI{4}{m}$ et $SA = SB = SC = SD = \SI{6}{m}$.
\begin{enumerate}
\item Démontrer que la droite $(SO)$ est orthogonale au plan $(ABC)$.
\item En déduire que $(SO)$ est orthogonale à $(BD)$.
\end{enumerate}
\tcblower
\textbf{1.} Travaillons dans les faces latérales. Le triangle $SAC$ est isocèle en $S$ (car $SA = SC$) et $O$ est le milieu de $[AC]$ (centre du carré, intersection des diagonales). Dans un triangle isocèle, la médiane issue du sommet principal est aussi hauteur : donc $(SO) \perp (AC)$.

De même, le triangle $SBD$ est isocèle en $S$ (car $SB = SD$) et $O$ est le milieu de $[BD]$, donc $(SO) \perp (BD)$.

Les droites $(AC)$ et $(BD)$ sont :
\begin{itemize}
\item contenues dans le plan $(ABC)$ (ce sont les diagonales de la base) ;
\item \textbf{sécantes en $O$} (diagonales du carré).
\end{itemize}
La droite $(SO)$ est donc orthogonale à deux droites sécantes du plan $(ABC)$ : d'après le théorème fondamental,
\[ (SO) \perp (ABC). \]

\medskip
\textbf{2.} Une droite orthogonale à un plan est orthogonale à toute droite de ce plan. Comme $(BD) \subset (ABC)$, on a immédiatement $(SO) \perp (BD)$. (On le savait déjà par le triangle isocèle $SBD$, mais la propriété générale le redonne sans calcul.)

\textbf{Conclusion} : $(SO)$ est la hauteur de la pyramide, orthogonale à toute droite de la base.
\end{exoresolu}

\courssub{Orthogonalité et parallélisme : théorèmes de liaison}

\begin{aretenir}[Théorèmes de liaison orthogonalité--parallélisme]
\begin{enumerate}
\item Si \textbf{deux droites sont parallèles}, toute droite orthogonale à l'une est orthogonale à l'autre. \quad ($d \parallel d'$ et $d \perp P \implies d' \perp P$)
\item Si \textbf{deux droites sont orthogonales à un même plan}, alors elles sont parallèles entre elles. \quad ($d \perp P$ et $d' \perp P \implies d \parallel d'$)
\item Si \textbf{deux plans sont parallèles}, toute droite orthogonale à l'un est orthogonale à l'autre ; et deux plans orthogonaux à une même droite sont parallèles entre eux. \quad ($P \parallel Q$ et $d \perp P \implies d \perp Q$)
\end{enumerate}
\end{aretenir}

\begin{methode}[Orthogonalité et parallélisme]
Démarche type pour « transporter » une orthogonalité :
\begin{enumerate}
\item Repérer dans la figure les parallélismes « gratuits » (faces opposées d'un parallélépipède, droites des milieux, arêtes latérales d'un prisme droit).
\item Établir une première orthogonalité $(d) \perp (P)$ par la fiche 3 (deux droites sécantes).
\item \emph{Transporter} cette orthogonalité par parallélisme : soit vers une droite parallèle à $(d)$, soit vers un plan parallèle à $(P)$.
\item Conclure en citant précisément le théorème utilisé (numéro 1, 2 ou 3 ci-dessus).
\end{enumerate}
\textbf{Réflexe rédaction} : toujours écrire la chaîne complète « $(d) \parallel (d')$ et $(d) \perp (P)$, donc $(d') \perp (P)$ » ou « $(P) \parallel (Q)$ et $(d) \perp (P)$, donc $(d) \perp (Q)$ ».
\end{methode}

\begin{exoresolu}[Transport d'orthogonalité dans un pavé droit]
Soit $ABCDEFGH$ un pavé droit (parallélépipède rectangle), modélisant une salle de classe : $ABCD$ est le sol, $EFGH$ le plafond.
\begin{enumerate}
\item Justifier que $(AE) \perp (ABC)$.
\item En déduire que $(CG) \perp (EFG)$.
\end{enumerate}
\tcblower
\textbf{1.} Dans le pavé droit, l'arête $[AE]$ est perpendiculaire aux arêtes $[AB]$ et $[AD]$ (les faces $ABFE$ et $ADHE$ sont des rectangles, donc les angles en $A$ sont droits). Les droites $(AB)$ et $(AD)$ sont deux droites du plan $(ABC)$, sécantes en $A$. Donc, par le théorème fondamental :
\[ (AE) \perp (ABC). \]

\medskip
\textbf{2.} On procède en deux transports.

\emph{Étape 1 : transport sur une droite parallèle.} Dans un parallélépipède, les arêtes latérales sont parallèles : $(AE) \parallel (CG)$. Or toute droite parallèle à une droite orthogonale à un plan est elle-même orthogonale à ce plan (théorème 1). Donc :
\[ (CG) \perp (ABC). \]

\emph{Étape 2 : transport sur un plan parallèle.} Les faces du sol et du plafond sont parallèles : $(ABC) \parallel (EFG)$. Or toute droite orthogonale à un plan est orthogonale à tout plan qui lui est parallèle (théorème 3). Donc :
\[ (CG) \perp (EFG). \]

\textbf{Conclusion} : l'arête verticale $(CG)$ est orthogonale au plan du plafond $(EFG)$, ce qui est bien l'image concrète d'un mur vertical rencontrant le plafond à angle droit.
\end{exoresolu}

\courssub{Plans orthogonaux}

\begin{aretenir}[Plans orthogonaux]
Deux plans $(P)$ et $(Q)$ sont \cle{orthogonaux} si l'un d'eux contient une droite orthogonale à l'autre.
\[ (P) \perp (Q) \iff \exists\, (d) \subset (P) \text{ tel que } (d) \perp (Q) \]
\textbf{Propriété réciproque utile} : si $(P) \perp (Q)$ et si une droite de $(P)$ est perpendiculaire à la \emph{droite d'intersection} de $(P)$ et $(Q)$, alors elle est orthogonale à $(Q)$.
\end{aretenir}

\begin{methode}[Plans orthogonaux]
Démarche pour prouver $(P) \perp (Q)$ :
\begin{enumerate}
\item \textbf{Chercher la « droite pivot »} : une droite $(d)$ contenue dans $(P)$ et orthogonale au plan $(Q)$ (ou l'inverse). C'est presque toujours une arête, une hauteur ou une médiane de la figure.
\item \textbf{Prouver $(d) \perp (Q)$} par la fiche 3 (orthogonalité à deux droites sécantes de $(Q)$).
\item \textbf{Vérifier l'appartenance} : $(d) \subset (P)$.
\item \textbf{Conclure} : $(P) \perp (Q)$.
\end{enumerate}
\textbf{Piège} : il ne suffit pas que deux plans se coupent « à angle droit en apparence » sur la figure ; il faut exhiber la droite orthogonale et le prouver.
\end{methode}

\begin{exoresolu}[Plans orthogonaux dans une pyramide]
On reprend la pyramide régulière $SABCD$ de l'exercice précédent : base carrée $ABCD$ de centre $O$, avec $(SO) \perp (ABC)$.
\begin{enumerate}
\item Démontrer que les plans $(SAC)$ et $(ABC)$ sont orthogonaux.
\item Démontrer que les plans $(SAC)$ et $(SBD)$ sont orthogonaux.
\end{enumerate}
\tcblower
\textbf{1.} La droite $(SO)$ est orthogonale au plan $(ABC)$ (démontré précédemment). Or $S \in (SAC)$ et $O \in [AC] \subset (SAC)$, donc la droite $(SO)$ est \textbf{contenue dans le plan $(SAC)$}.

Le plan $(SAC)$ contient donc une droite orthogonale au plan $(ABC)$ : par définition,
\[ (SAC) \perp (ABC). \]

\medskip
\textbf{2.} Cherchons une droite de l'un des plans orthogonale à l'autre. Montrons que $(BD) \perp (SAC)$ :
\begin{itemize}
\item $(BD) \perp (AC)$ : les diagonales du carré $ABCD$ sont perpendiculaires ;
\item $(BD) \perp (SO)$ : car $(SO) \perp (ABC)$ et $(BD) \subset (ABC)$, donc $(SO)$ est orthogonale à toute droite de la base.
\end{itemize}
Les droites $(AC)$ et $(SO)$ sont deux droites du plan $(SAC)$, \textbf{sécantes en $O$} (car $O \in [AC]$ et $O \in (SO)$). La droite $(BD)$ est donc orthogonale à deux droites sécantes de $(SAC)$ : d'où
\[ (BD) \perp (SAC). \]
Or $(BD) \subset (SBD)$ (puisque $B$ et $D$ sont des points du plan $(SBD)$). Le plan $(SBD)$ contient donc une droite orthogonale au plan $(SAC)$ :
\[ (SBD) \perp (SAC). \]

\textbf{Conclusion} : les deux plans diagonaux de la pyramide sont orthogonaux entre eux, et chacun est orthogonal à la base.
\end{exoresolu}

\courssub{Méthodes de démonstration et applications métriques}

\begin{aretenir}[Distance, hauteur et volume]
\begin{itemize}
\item \textbf{Distance d'un point $M$ à un plan $(P)$} : c'est la longueur $MH$, où $H$ est le \emph{projeté orthogonal} de $M$ sur $(P)$ (c'est-à-dire $(MH) \perp (P)$). Identifier ou construire $H$ est la première étape.
\item \textbf{Hauteur d'un solide} : pour une pyramide ou un cône, la hauteur est la distance du sommet au plan de base. Dans une pyramide régulière, c'est $SO$ où $O$ est le centre de la base.
\item \textbf{Calculer par Pythagore} : une fois la hauteur identifiée, on se place dans un triangle rectangle bien choisi (souvent formé par le sommet, le centre de la base et un sommet de la base) pour trouver la longueur manquante.
\item \textbf{Volume} : $V = \dfrac{1}{3} \times \mathcal{B} \times h$ pour une pyramide, où $\mathcal{B}$ est l'aire de la base.
\end{itemize}
\textbf{Réflexe} : toujours nommer et justifier le caractère rectangle du triangle avant d'appliquer Pythagore.
\end{aretenir}

\begin{methode}[Exploiter l'orthogonalité pour calculer]
\begin{enumerate}
\item Identifier la droite orthogonale au plan de référence (hauteur, arête verticale, etc.).
\item Repérer le triangle rectangle qui contient la longueur cherchée (hauteur, distance, longueur d'arête).
\item Écrire la relation de Pythagore \textbf{après avoir justifié l'angle droit} (par orthogonalité droite/plan ou propriété du solide).
\item Calculer la valeur exacte (forme radicaux) puis, si demandé, la valeur approchée.
\item Pour un volume : calculer l'aire de la base $\mathcal{B}$, puis appliquer $V = \frac{1}{3}\mathcal{B}h$.
\end{enumerate}
\end{methode}

\begin{exoresolu}[Hauteur et volume d'un toit pyramidal]
Un artisan de Bafoussam construit un clocheton en forme de pyramide régulière $SABCD$ : base carrée de côté $AB = \SI{4}{m}$, arêtes latérales $SA = \SI{6}{m}$. On note $O$ le centre de la base.
\begin{enumerate}
\item Calculer la longueur exacte $OA$.
\item En déduire la hauteur $SO$ du clocheton.
\item Calculer le volume du clocheton (on donnera la valeur exacte, puis une valeur approchée à $\SI{0,1}{m^3}$ près).
\end{enumerate}
\tcblower
\textbf{1.} Les diagonales d'un carré de côté $a$ ont pour longueur $a\sqrt{2}$. Ici $AC = 4\sqrt{2}$, et $O$ est le milieu de $[AC]$, donc :
\[ OA = \frac{AC}{2} = \frac{4\sqrt{2}}{2} = 2\sqrt{2} \text{ m}. \]

\medskip
\textbf{2.} On a démontré (fiche 3) que $(SO) \perp (ABC)$. En particulier $(SO) \perp (OA)$, donc le triangle $SOA$ est \textbf{rectangle en $O$}. Le théorème de Pythagore donne :
\begin{align*}
SO^2 + OA^2 &= SA^2 \\
SO^2 &= SA^2 - OA^2 = 6^2 - (2\sqrt{2})^2 = 36 - 8 = 28 \\
SO &= \sqrt{28} = 2\sqrt{7} \text{ m}.
\end{align*}

\medskip
\textbf{3.} L'aire de la base carrée est $\mathcal{B} = AB^2 = \SI{16}{m^2}$. Le volume de la pyramide est :
\[ V = \frac{1}{3} \times \mathcal{B} \times SO = \frac{1}{3} \times 16 \times 2\sqrt{7} = \frac{32\sqrt{7}}{3} \text{ m}^3 \approx \SI{28,2}{m^3}. \]

\textbf{Conclusion} : le clocheton a une hauteur de $2\sqrt{7} \approx \SI{5,3}{m}$ et un volume d'environ $\SI{28,2}{m^3}$.
\end{exoresolu}

\begin{savaistu}[L'orthogonalité dans l'architecture camerounaise]
Les toits traditionnels des cases (rondes ou carrées) du Grassland ou du Sahel illustrent parfaitement l'orthogonalité dans l'espace. La charpente en bois repose sur des poteaux verticaux (orthogonaux au plan du sol) et des poutres horizontales (parallèles au sol). La faîtière, ligne de rencontre des deux pans de toit, est souvent orthogonale aux murs pignons. Les architectes modernes (comme pour le Palais des Congrès de Yaoundé ou les nouveaux stades) utilisent ces mêmes principes : les poteaux sont orthogonaux aux dalles, les poutres orthogonales aux poteaux, garantissant la stabilité par la transmission verticale des charges. Le calcul des longueurs de chevrons fait appel au théorème de Pythagore dans les triangles rectangles formés par la hauteur du toit, la demi-largeur de la maison et la pente du toit.
\end{savaistu}

\begin{attention}[Les 5 erreurs qui coûtent des points au Probatoire]
\begin{enumerate}
\item \textbf{Oublier la condition « deux droites sécantes ».} Dire « $(d)$ est perpendiculaire à deux droites du plan $(P)$, donc $(d) \perp (P)$ » est \emph{faux} si les deux droites sont parallèles ! Il faut impérativement écrire : « $(\Delta_1)$ et $(\Delta_2)$ sont deux droites \textbf{sécantes} de $(P)$ » et le justifier (par exemple : « sécantes en $A$ » ou « non parallèles car... »).
\item \textbf{Confondre orthogonales et perpendiculaires.} Deux droites orthogonales ne se rencontrent pas forcément. N'écris « perpendiculaires » que si tu as prouvé que les droites sont \emph{sécantes} ; sinon écris « orthogonales ». Au Probatoire, ce vocabulaire est sanctionné.
\item \textbf{Affirmer une orthogonalité à partir du dessin.} Une figure en perspective peut montrer un angle qui « semble » droit sans l'être. Toute orthogonalité doit être \emph{démontrée} (théorème, propriété du solide, calcul), jamais lue sur la figure.
\item \textbf{Appliquer Pythagore sans justifier le triangle rectangle.} Avant d'écrire $SO^2 + OA^2 = SA^2$, il faut avoir établi (et rédigé !) que le triangle $SOA$ est rectangle en $O$, ce qui repose sur $(SO) \perp (ABC)$. Une application de Pythagore non justifiée perd la moitié des points de la question.
\item \textbf{Mal conclure sur les plans orthogonaux.} Pour prouver $(P) \perp (Q)$, il faut exhiber une droite $(d)$ avec \emph{deux} appartenances vérifiées : $(d) \subset (P)$ \textbf{et} $(d) \perp (Q)$. Oublier de mentionner que la droite est contenue dans l'un des plans rend la démonstration incomplète.
\end{enumerate}
\end{attention}

\newpage

\coursec{Exercices}

\courssub{Je vérifie mes connaissances}

\begin{exercice}[QCM : une seule bonne réponse]
Pour chaque question, coche la case correspondant à l'unique réponse exacte.
\begin{enumerate}
\item Deux droites orthogonales dans l'espace sont :
\begin{enumerate}
\item \textbf{a)} toujours sécantes ;
\item \textbf{b)} toujours coplanaires ;
\item \textbf{c)} telles que la parallèle à l'une menée par un point de l'autre est perpendiculaire à cette autre ;
\item \textbf{d)} toujours parallèles.
\end{enumerate}
\item Une droite $(d)$ est orthogonale à un plan $(P)$ si et seulement si $(d)$ est orthogonale à :
\begin{enumerate}
\item \textbf{a)} une droite de $(P)$ ;
\item \textbf{b)} deux droites sécantes de $(P)$ ;
\item \textbf{c)} toutes les droites de l'espace ;
\item \textbf{d)} un plan parallèle à $(P)$.
\end{enumerate}
\item Si deux plans $(P)$ et $(Q)$ sont orthogonaux, alors toute droite de $(P)$ est :
\begin{enumerate}
\item \textbf{a)} orthogonale à $(Q)$ ;
\item \textbf{b)} parallèle à $(Q)$ ;
\item \textbf{c)} orthogonale à la droite d'intersection de $(P)$ et $(Q)$ ;
\item \textbf{d)} aucune des réponses précédentes n'est toujours vraie.
\end{enumerate}
\item Dans un cube $ABCDEFGH$ (avec $ABCD$ face inférieure, $EFGH$ face supérieure, $E$ au-dessus de $A$), les droites $(AB)$ et $(CG)$ sont :
\begin{enumerate}
\item \textbf{a)} sécantes ; \quad \textbf{b)} parallèles ; \quad \textbf{c)} orthogonales non sécantes ; \quad \textbf{d)} confondues.
\end{enumerate}
\end{enumerate}
\end{exercice}

\begin{exercice}[Vrai ou faux ? Justifie]
Réponds par VRAI ou FAUX en justifiant chaque réponse par une propriété du cours ou un contre-exemple (tu pourras t'appuyer sur un cube).
\begin{enumerate}
\item Deux droites orthogonales à une même droite sont parallèles.
\item Deux plans orthogonaux à un même plan sont parallèles.
\item Deux droites non sécantes sont parallèles.
\item Si une droite est orthogonale à deux droites d'un plan, alors elle est orthogonale à ce plan.
\item Deux droites perpendiculaires sont orthogonales.
\end{enumerate}
\end{exercice}

\begin{exercice}[Définitions à compléter]
Recopie et complète les phrases suivantes avec les mots ou expressions qui conviennent.
\begin{enumerate}
\item Deux droites de l'espace sont dites \cle{coplanaires} lorsqu'il existe \textbf{un plan} les contenant toutes les deux.
\item Deux droites sont \cle{orthogonales} lorsque les parallèles menées à ces droites par un point quelconque de l'espace sont \textbf{perpendiculaires}.
\item Une droite $(d)$ est orthogonale à un plan $(P)$ lorsqu'elle est orthogonale à \textbf{deux droites sécantes} de $(P)$.
\item Deux plans sont \cle{orthogonaux} lorsque l'un d'eux contient une droite \textbf{orthogonale} à l'autre.
\end{enumerate}
\end{exercice}

\begin{exercice}[Lecture de figure dans un cube]
Soit $ABCDEFGH$ un cube ($ABCD$ face inférieure, $EFGH$ face supérieure, $E$ au-dessus de $A$). Sans démonstration longue, donne la position relative (sécantes, parallèles, non coplanaires) des couples de droites suivants, puis précise si elles sont orthogonales :
\begin{enumerate}
\item $(AB)$ et $(BC)$ ;
\item $(AB)$ et $(EF)$ ;
\item $(AB)$ et $(FG)$ ;
\item $(AG)$ et $(BC)$ ;
\item $(AC)$ et $(EG)$.
\end{enumerate}
\end{exercice}

\courssub{Je m'entraîne}

\begin{exercice}[Positions relatives dans un cube]
Soit $ABCDEFGH$ un cube. Pour chacun des couples de droites suivants, détermine en justifiant si les droites sont coplanaires ou non, et si elles sont orthogonales ou non :
\begin{enumerate}
\item $(EH)$ et $(BC)$ ;
\item $(AC)$ et $(BF)$ ;
\item $(HB)$ et $(AD)$ ;
\item $(EC)$ et $(AB)$.
\end{enumerate}
\emph{Indication : pour la non-coplanarité, on pourra raisonner par l'absurde en exhibant une contradiction avec les propriétés du cube.}
\end{exercice}

\begin{exercice}[Tétraèdre à arête orthogonale à une face]
Soit $ABCD$ un tétraèdre tel que la droite $(AB)$ soit orthogonale au plan $(BCD)$.
\begin{enumerate}
\item Fais une figure en perspective cavalière.
\begin{popfigure}
\begin{tikzpicture}[scale=1.2, >=Stealth]
% Base BCD horizontal (cavalier: x right, y up-left, z up)
\coordinate (B) at (0,0);
\coordinate (C) at (3,0);
\coordinate (D) at (1.5,2.6);
% A above B (since AB ortho BCD)
\coordinate (A) at (0,3);
% Draw base
\draw[popblue, thick] (B) -- (C) -- (D) -- cycle;
\draw[popblue, dashed] (B) -- (D);
% Draw vertical AB
\draw[poporange, very thick] (A) -- (B) node[midway, left] {$AB$};
% Draw other edges
\draw[popblue, thick] (A) -- (C);
\draw[popblue, thick] (A) -- (D);
% Labels
\node[below left] at (B) {$B$};
\node[below right] at (C) {$C$};
\node[above] at (D) {$D$};
\node[above left] at (A) {$A$};
% Right angle mark at B in plane BCD? No, AB ortho plane.
% Mark orthogonality AB / plane
\draw[poporange] (0.2,0.2) -- (0.4,0) -- (0.2,-0.2); % symbol at B
\end{tikzpicture}
\legende{Tétraèdre $ABCD$ avec $(AB) \perp (BCD)$ (perspective cavalière).}
\end{popfigure}
\item Montre que $(AB)$ est orthogonale à $(CD)$.
\item On suppose de plus que le triangle $BCD$ est rectangle en $B$. Montre que $(CD)$ est orthogonale au plan $(ABC)$.
\item Déduis-en que les plans $(ABC)$ et $(ACD)$ sont orthogonaux.
\end{enumerate}
\end{exercice}

\begin{exercice}[Pyramide à base carrée]
Soit $SABCD$ une pyramide de sommet $S$, dont la base $ABCD$ est un carré de centre $O$, telle que la droite $(SO)$ soit orthogonale au plan $(ABCD)$. On donne $AB = 6$ cm et $SO = 4$ cm.
\begin{popfigure}
\begin{tikzpicture}[scale=0.8, >=Stealth]
% Base square ABCD
\coordinate (A) at (-3,-3);
\coordinate (B) at (3,-3);
\coordinate (C) at (3,3);
\coordinate (D) at (-3,3);
\coordinate (O) at (0,0);
% Summit S above O
\coordinate (S) at (0,5);
% Draw base
\draw[popblue, thick] (A) -- (B) -- (C) -- (D) -- cycle;
\draw[popblue, dashed] (A) -- (C) (B) -- (D);
% Draw lateral edges
\draw[popblue, thick] (S) -- (A) (S) -- (B) (S) -- (C) (S) -- (D);
% Draw height SO
\draw[poporange, very thick, dashed] (S) -- (O) node[midway, right] {$SO=4$};
% Labels
\node[below left] at (A) {$A$};
\node[below right] at (B) {$B$};
\node[above right] at (C) {$C$};
\node[above left] at (D) {$D$};
\node[above] at (S) {$S$};
\node[below] at (O) {$O$};
% Dimensions
\node[below] at (0,-3) {$6$ cm};
\end{tikzpicture}
\legende{Pyramide $SABCD$ à base carrée, $(SO) \perp (ABCD)$.}
\end{popfigure}
\begin{enumerate}
\item Justifie que $(SO)$ est orthogonale aux droites $(AC)$ et $(BD)$.
\item Calcule la longueur $OA$, puis la longueur $SA$.
\item Montre que $(BD)$ est orthogonale au plan $(SAC)$.
\item Déduis-en que les plans $(SAC)$ et $(SBD)$ sont orthogonaux.
\end{enumerate}
\end{exercice}

\begin{exercice}[Raisonnement par l'absurde]
Soit $ABCDEFGH$ un cube. On veut montrer que les droites $(AC)$ et $(BH)$ ne sont pas coplanaires.
\begin{enumerate}
\item Suppose, par l'absurde, qu'il existe un plan $(P)$ contenant $(AC)$ et $(BH)$. Montre que $(P)$ contiendrait alors les points $A$, $B$, $C$ et $H$.
\item Montre que le plan $(ABC)$ est le plan $(ABCD)$ et que $H \notin (ABCD)$.
\item Conclus.
\item Les droites $(AC)$ et $(BH)$ sont-elles orthogonales ? Justifie.
\end{enumerate}
\end{exercice}

\courssub{Je me prépare au Probatoire}

\begin{exercice}[Pavé droit : la diagonale orthogonale à un plan]
Lors de la construction d'un entrepôt à Nkongsamba, le maçon modélise une poutre métallique par la diagonale $[HF]$ d'un pavé droit $ABCDEFGH$ tel que $AB = a$, $AD = b$ et $AE = c$ ($ABCD$ face inférieure, $EFGH$ face supérieure, $E$ au-dessus de $A$). On suppose dans tout l'exercice que $a = b$ (la face $ABCD$ est un carré).
\begin{popfigure}
\begin{tikzpicture}[scale=0.7, >=Stealth]
% Coordinates for cavalier perspective (alpha=45, k=0.5 approx)
% A(0,0), B(a,0), D(0.5b, 0.5b), E(0,c) -> but a=b
\def\a{3} \def\c{4}
\coordinate (A) at (0,0);
\coordinate (B) at (\a,0);
\coordinate (D) at (0.5*\a, 0.5*\a);
\coordinate (E) at (0,\c);
\coordinate (C) at ($ (B) + (D) - (A) $);
\coordinate (F) at ($ (B) + (E) - (A) $);
\coordinate (H) at ($ (D) + (E) - (A) $);
\coordinate (G) at ($ (C) + (E) - (A) $);
% Hidden edges
\draw[popmuted, dashed] (A) -- (D) -- (H) -- (E) -- cycle;
\draw[popmuted, dashed] (D) -- (C);
\draw[popmuted, dashed] (A) -- (B);
\draw[popmuted, dashed] (A) -- (E);
% Visible edges
\draw[popblue, thick] (B) -- (C) -- (G) -- (F) -- cycle;
\draw[popblue, thick] (B) -- (F);
\draw[popblue, thick] (C) -- (G);
\draw[popblue, thick] (H) -- (G);
\draw[popblue, thick] (E) -- (F);
\draw[popblue, thick] (E) -- (H);
% Diagonals
\draw[poporange, very thick] (H) -- (F) node[midway, above left] {$HF$};
\draw[popgreen, thick] (E) -- (G) node[midway, above right] {$EG$};
\draw[poppurple, thick] (A) -- (G) node[midway, below right] {$AG$};
% Labels
\node[below left] at (A) {$A$};
\node[below right] at (B) {$B$};
\node[above right] at (C) {$C$};
\node[above left] at (D) {$D$};
\node[below left] at (E) {$E$};
\node[below right] at (F) {$F$};
\node[above right] at (G) {$G$};
\node[above left] at (H) {$H$};
% Dimensions
\node[below] at ($ (A)!0.5!(B) $) {$a$};
\node[left] at ($ (A)!0.5!(E) $) {$c$};
\end{tikzpicture}
\legende{Pavé droit $ABCDEFGH$ à base carrée ($a=b$).}
\end{popfigure}

\textbf{Partie A : orthogonalités intermédiaires}
\begin{enumerate}
\item Montre que $(AE)$ est orthogonale au plan $(EFGH)$, puis que $(AE)$ est orthogonale à $(HF)$.
\item Montre que la droite $(HF)$ est orthogonale à la droite $(EG)$.
\end{enumerate}

\textbf{Partie B : orthogonalité d'une droite et d'un plan}
\begin{enumerate}
\setcounter{enumi}{2}
\item Déduis de la partie A que la droite $(HF)$ est orthogonale au plan $(AEG)$.
\item Montre que les plans $(AEG)$ et $(EFGH)$ sont orthogonaux.
\end{enumerate}

\textbf{Partie C : calculs}
\begin{enumerate}
\setcounter{enumi}{4}
\item On donne $a = b = 3$ m et $c = 4$ m. Calcule les longueurs $EG$ et $AG$.
\item Soit $I$ le point d'intersection de $(HF)$ et $(EG)$. Calcule la distance $AI$.
\end{enumerate}
\end{exercice}

\begin{exercice}[Le poteau et ses haubans]
Dans le quartier Ndogpassi à Douala, un poteau électrique vertical est maintenu par deux haubans (câbles tendus) fixés au sol. On modélise la situation ainsi : le sol est un plan horizontal $(P)$, le poteau est le segment $[SO]$ avec $O \in (P)$ et $S$ le sommet du poteau, et les haubans sont les segments $[SA]$ et $[SB]$ où $A$ et $B$ sont deux points du sol tels que $O$, $A$, $B$ ne soient pas alignés. Le poteau est vertical, ce qui signifie que $(SO)$ est perpendiculaire en $O$ aux droites $(OA)$ et $(OB)$. On donne $SO = 6$ m, $OA = 2{,}5$ m et $OB = 4$ m.
\begin{popfigure}
\begin{tikzpicture}[scale=0.6, >=Stealth]
% Ground plane P (horizontal line)
\draw[popgreen, thick] (-5,0) -- (5,0) node[right] {$(P)$};
% Points on ground
\coordinate (O) at (0,0);
\coordinate (A) at (-2.5,0);
\coordinate (B) at (4,0);
% Pole SO vertical
\coordinate (S) at (0,6);
% Draw pole
\draw[poporange, very thick] (S) -- (O) node[midway, right] {$SO=6$ m};
% Draw cables
\draw[popblue, thick] (S) -- (A) node[midway, left] {$SA$};
\draw[popblue, thick] (S) -- (B) node[midway, right] {$SB$};
% Ground distances
\draw[popmuted, <->] (O) -- (A) node[midway, below] {$OA=2,5$ m};
\draw[popmuted, <->] (O) -- (B) node[midway, below] {$OB=4$ m};
% Right angle marks
\draw[poporange] (0.2,0) -- (0.2,0.2) -- (0,0.2);
\draw[poporange] (-0.2,0) -- (-0.2,0.2) -- (0,0.2); % at O for OA? OA is horizontal.
% Actually right angle is between SO (vertical) and OA (horizontal).
% Mark at O between SO and OA
\draw[poporange] (0,0.2) -- (0.2,0.2) -- (0.2,0);
\draw[poporange] (0,0.2) -- (-0.2,0.2) -- (-0.2,0);
% Labels
\node[below left] at (O) {$O$};
\node[below] at (A) {$A$};
\node[below] at (B) {$B$};
\node[above] at (S) {$S$};
\end{tikzpicture}
\legende{Modélisation du poteau et de ses haubans.}
\end{popfigure}

\textbf{Partie A : modélisation et orthogonalité}
\begin{enumerate}
\item Justifie que la droite $(SO)$ est orthogonale au plan $(P)$.
\item En déduire que les triangles $SOA$ et $SOB$ sont rectangles en $O$.
\item Montre que le plan $(SOA)$ est orthogonal au plan $(P)$.
\end{enumerate}

\textbf{Partie B : calcul des longueurs de câble}
\begin{enumerate}
\setcounter{enumi}{3}
\item Calcule la longueur de chacun des haubans $SA$ et $SB$ (on donnera les valeurs exactes, puis des valeurs approchées à $10^{-2}$ m près).
\item L'entreprise d'électrification rurale dispose de câbles vendus par rouleaux de $15$ m. Un rouleau suffit-il pour les deux haubans ? Justifie.
\item On suppose de plus que $\widehat{AOB} = 90^{\circ}$. Calcule la distance $AB$ entre les deux points d'ancrage au sol.
\end{enumerate}
\end{exercice}

\begin{exercice}[Tétraèdre régulier et orthogonalité]
Un artisan de Foumban taille un objet décoratif ayant la forme d'un tétraèdre $ABCD$ dont toutes les arêtes ont la même longueur $a$ (tétraèdre régulier). On note $I$ le milieu de $[CD]$ et $H$ le pied de la hauteur issue de $A$ dans le triangle $BCD$.
\begin{popfigure}
\begin{tikzpicture}[scale=1.2, >=Stealth]
% Regular tetrahedron coordinates (approx cavalier)
% Base BCD equilateral horizontal
\coordinate (B) at (0,0);
\coordinate (C) at (3,0);
\coordinate (D) at (1.5,2.6);
\coordinate (I) at ($ (C)!0.5!(D) $); % Midpoint CD
% A above centroid of BCD
\coordinate (G) at (1.5,0.866); % Centroid approx
\coordinate (A) at (1.5,3.5);
% H is foot of height from A to BCD -> H = G (centroid)
\coordinate (H) at (G);
% Draw base
\draw[popblue, thick] (B) -- (C) -- (D) -- cycle;
% Draw edges from A
\draw[popblue, thick] (A) -- (B);
\draw[popblue, thick] (A) -- (C);
\draw[popblue, thick] (A) -- (D);
% Medians in base
\draw[popmuted, dashed] (B) -- (I);
\draw[popmuted, dashed] (C) -- ($ (B)!0.5!(D) $);
\draw[popmuted, dashed] (D) -- ($ (B)!0.5!(C) $);
% Height AH
\draw[poporange, very thick, dashed] (A) -- (H) node[midway, right] {$AH$};
% AI and BI
\draw[poppurple, thick] (A) -- (I) node[midway, left] {$AI$};
\draw[poppurple, thick] (B) -- (I) node[midway, left] {$BI$};
% CD
\draw[poporange, very thick] (C) -- (D) node[midway, above] {$CD$};
% Right angle marks AI/CD and BI/CD
\draw[poporange] ($ (I)!0.2!(C) $) -- ($ (I)!0.2!(C) + (0.1,0.1) $) -- ($ (I)!0.2!(A) $); % approx
% Labels
\node[below left] at (B) {$B$};
\node[below right] at (C) {$C$};
\node[above] at (D) {$D$};
\node[above] at (A) {$A$};
\node[below] at (I) {$I$};
\node[below left] at (H) {$H$};
\end{tikzpicture}
\legende{Tétraèdre régulier $ABCD$, $I$ milieu de $[CD]$, $H$ centre de gravité de $BCD$.}
\end{popfigure}

\textbf{Partie A : orthogonalités}
\begin{enumerate}
\item Montre que les triangles $ACD$ et $BCD$ sont équilatéraux, et que $(CD)$ est perpendiculaire à $(AI)$ et à $(BI)$.
\item Déduis-en que la droite $(CD)$ est orthogonale au plan $(ABI)$.
\item Montre alors que $(CD)$ est orthogonale à $(AB)$. Que peut-on dire de deux arêtes opposées d'un tétraèdre régulier ?
\item Montre que les plans $(ABI)$ et $(BCD)$ sont orthogonaux.
\end{enumerate}

\textbf{Partie B : calculs}
\begin{enumerate}
\setcounter{enumi}{4}
\item On donne $a = 6$ cm. Calcule $AI$ et $BI$.
\item Calcule l'aire du triangle $ABI$ (on pourra remarquer que $I$ est le milieu de $[CD]$ et utiliser la hauteur issue de $I$ dans le triangle $ABI$, ou montrer que ce triangle est isocèle et calculer sa hauteur).
\item Calcule la distance $BH$, sachant que $H$ est le centre de gravité du triangle équilatéral $BCD$ (on rappelle que le centre de gravité se trouve aux deux tiers de chaque médiane à partir du sommet).
\end{enumerate}
\end{exercice}

\newpage

\coursec{Corrigés des exercices}

\begin{corrige}[Exercice 1 — QCM : une seule bonne réponse]
\begin{enumerate}
\item Réponse \textbf{c)}. Par définition, deux droites $(d)$ et $(\Delta)$ sont orthogonales lorsque la parallèle à $(d)$ menée par un point quelconque de $(\Delta)$ est perpendiculaire à $(\Delta)$. Les réponses a) et b) sont fausses : deux droites orthogonales peuvent être non coplanaires (donc non sécantes). La réponse d) est absurde.
\item Réponse \textbf{b)}. C'est le théorème fondamental : $(d) \perp (P)$ si et seulement si $(d)$ est orthogonale à deux droites \emph{sécantes} de $(P)$. Une seule droite ne suffit pas.
\item Réponse \textbf{d)}. Contre-exemples dans un cube $ABCDEFGH$ : les plans $(ABCD)$ et $(ABFE)$ sont orthogonaux, pourtant $(AB) \subset (ABCD)$ est contenue dans $(ABFE)$ (ni orthogonale, ni parallèle à ce plan au sens strict), et $(AD)$ est parallèle à $(BC) \subset (ABFE)$ sans être orthogonale à l'intersection. Aucune des trois affirmations n'est toujours vraie.
\item Réponse \textbf{c)}. $(AB)$ et $(CG)$ ne sont pas coplanaires (sinon $A, B, C, G$ seraient coplanaires, ce qui est faux car $G \notin (ABCD)$ et $G \notin (ABC)$). Or $(CG) \parallel (BF)$ et $(BF) \perp (AB)$ dans la face $ABFE$ : donc $(AB) \perp (CG)$ au sens des droites orthogonales, sans être sécantes.
\end{enumerate}
\end{corrige}

\begin{corrige}[Exercice 2 — Vrai ou faux ? Justifie]
On s'appuie sur le cube $ABCDEFGH$.
\begin{enumerate}
\item \textbf{FAUX.} Contre-exemple : $(AD) \perp (AB)$ et $(AE) \perp (AB)$, pourtant $(AD)$ et $(AE)$ sont sécantes en $A$, non parallèles. Dans l'espace, l'orthogonalité à une même droite ne donne pas le parallélisme (contrairement au plan).
\item \textbf{FAUX.} Contre-exemple : les plans $(ABFE)$ et $(ADHE)$ sont tous deux orthogonaux au plan $(ABCD)$ (chacun contient une arête verticale orthogonale à $(ABCD)$), pourtant ils sont sécants suivant $(AE)$.
\item \textbf{FAUX.} Deux droites non sécantes peuvent être non coplanaires : $(AB)$ et $(CG)$ ne se coupent pas et ne sont pas parallèles. « Non sécantes » n'implique « parallèles » que si les droites sont coplanaires.
\item \textbf{FAUX.} Il manque l'hypothèse que les deux droites du plan sont \emph{sécantes}. Contre-exemple : dans le plan $(ABCD)$ du cube, les droites $(AB)$ et $(CD)$ sont parallèles. La droite $(BC)$, contenue dans $(ABCD)$, est perpendiculaire à $(AB)$ et à $(CD)$, donc orthogonale à ces deux droites parallèles du plan. Pourtant $(BC)$ n'est pas orthogonale au plan $(ABCD)$ puisqu'elle y est contenue. La condition « sécantes » est donc indispensable.
\item \textbf{VRAI.} Deux droites perpendiculaires sont sécantes et forment un angle droit ; la parallèle à l'une menée par un point de l'autre est alors confondue avec elle-même (ou lui est parallèle), donc perpendiculaire à l'autre : elles sont orthogonales. La réciproque est fausse en général (droites orthogonales non coplanaires).
\end{enumerate}
\end{corrige}

\begin{corrige}[Exercice 3 — Définitions à compléter]
\begin{enumerate}
\item Deux droites de l'espace sont dites \cle{coplanaires} lorsqu'il existe \textbf{un plan} les contenant toutes les deux.
\item Deux droites sont \cle{orthogonales} lorsque les parallèles menées à ces droites par un point quelconque de l'espace sont \textbf{perpendiculaires}.
\item Une droite $(d)$ est orthogonale à un plan $(P)$ lorsqu'elle est orthogonale à \textbf{deux droites sécantes} de $(P)$. (Elle est alors orthogonale à \emph{toute} droite de $(P)$.)
\item Deux plans sont \cle{orthogonaux} lorsque l'un d'eux contient une droite \textbf{orthogonale} à l'autre.
\end{enumerate}
\end{corrige}

\begin{corrige}[Exercice 4 — Lecture de figure dans un cube]
\begin{enumerate}
\item $(AB)$ et $(BC)$ : sécantes en $B$, perpendiculaires (côtés consécutifs du carré $ABCD$), donc orthogonales.
\item $(AB)$ et $(EF)$ : parallèles (arêtes correspondantes des faces opposées), donc coplanaires (plan $(ABFE)$) ; non orthogonales.
\item $(AB)$ et $(FG)$ : $(FG) \parallel (BC)$ (face $EFGH$) et $(BC) \perp (AB)$, donc $(AB)$ et $(FG)$ sont orthogonales. Elles ne sont pas coplanaires (sinon $A, B, F, G$ coplanaires, or $G \notin (ABFE)$) : orthogonales non sécantes.
\item $(AG)$ et $(BC)$ : 
  \begin{itemize}
  \item \textbf{Coplanarité :} Si un plan contenait $(AG)$ et $(BC)$, il contiendrait $A, B, C$ donc serait le plan $(ABC) = (ABCD)$. Or $G \notin (ABCD)$. Donc \textbf{non coplanaires}.
  \item \textbf{Orthogonalité :} On utilise un repère : $A(0,0,0)$, $B(1,0,0)$, $C(1,1,0)$, $G(1,1,1)$. Vecteur directeur de $(BC)$ : $\vec{u}=(0,1,0)$. Vecteur directeur de $(AG)$ : $\vec{v}=(1,1,1)$. Produit scalaire $\vec{u}\cdot\vec{v}=1 \neq 0$. Donc \textbf{non orthogonales}.
  \end{itemize}
  Conclusion : non coplanaires, non orthogonales.
\item $(AC)$ et $(EG)$ : $(EG) \parallel (AC)$ (diagonales correspondantes des faces parallèles $ABCD$ et $EFGH$), donc coplanaires (plan $(AEGC)$) et parallèles ; non orthogonales.
\end{enumerate}
\end{corrige}

\begin{corrige}[Exercice 5 — Positions relatives dans un cube]
\begin{enumerate}
\item $(EH)$ et $(BC)$ : $(EH) \parallel (AD)$ (face $ADHE$) et $(AD) \parallel (BC)$ (face $ABCD$), donc $(EH) \parallel (BC)$ : elles sont coplanaires (plan $(EHCB)$, qui est une section rectangulaire du cube) et non orthogonales.
\item $(AC)$ et $(BF)$ : 
  \begin{itemize}
  \item \textbf{Coplanarité :} Supposons qu'un plan $(P)$ contienne $(AC)$ et $(BF)$. Alors $(P)$ contient $A$, $B$, $C$, donc $(P) = (ABC) = (ABCD)$ ; or $F \notin (ABCD)$ : contradiction. Donc \textbf{non coplanaires}.
  \item \textbf{Orthogonalité :} La parallèle à $(BF)$ menée par $A$ est $(AE)$ ; $(AE) \perp (ABCD)$ donc $(AE) \perp (AC)$ : les droites $(AC)$ et $(BF)$ sont \textbf{orthogonales} (non sécantes).
  \end{itemize}
\item $(HB)$ et $(AD)$ : 
  \begin{itemize}
  \item \textbf{Coplanarité :} Si un plan contenait $(HB)$ et $(AD)$, il contiendrait $A$, $D$, $B$, donc serait $(ABD) = (ABCD)$, or $H \notin (ABCD)$ : \textbf{non coplanaires}.
  \item \textbf{Orthogonalité :} Repère $A(0,0,0)$, $B(1,0,0)$, $D(0,1,0)$, $H(0,1,1)$. $\overrightarrow{HB}=(1,-1,-1)$, $\overrightarrow{AD}=(0,1,0)$. Produit scalaire $-1 \neq 0$ : \textbf{non orthogonales}.
  \end{itemize}
\item $(EC)$ et $(AB)$ : 
  \begin{itemize}
  \item \textbf{Coplanarité :} Un plan contenant $(EC)$ et $(AB)$ contiendrait $A$, $B$, $E$, donc serait $(ABE) = (ABFE)$, or $C \notin (ABFE)$ : \textbf{non coplanaires}.
  \item \textbf{Orthogonalité :} $E(0,0,1)$, $C(1,1,0)$ : $\overrightarrow{EC}=(1,1,-1)$. $\overrightarrow{AB}=(1,0,0)$. Produit scalaire $1 \neq 0$ : \textbf{non orthogonales}.
  \end{itemize}
\end{enumerate}
\begin{attention}[Point de vigilance]
Pour conclure « non coplanaires », il faut exhiber la contradiction : le plan candidat serait déterminé par trois points non alignés (deux sur une droite, un sur l'autre), et montrer que le quatrième point n'y appartient pas.
\end{attention}
\end{corrige}

\begin{corrige}[Exercice 6 — Tétraèdre à arête orthogonale à une face]
\begin{enumerate}
\item Figure : voir l'énoncé (perspective cavalière, base $BCD$ en vraie grandeur apparente, $(AB)$ verticale).
\item $(AB) \perp (BCD)$ par hypothèse, donc $(AB)$ est orthogonale à \emph{toute} droite du plan $(BCD)$, en particulier à $(CD) \subset (BCD)$.
\item On sait que $(CD) \perp (AB)$ (question 2) et $(CD) \perp (BC)$ car le triangle $BCD$ est rectangle en $B$. Les droites $(AB)$ et $(BC)$ sont deux droites \emph{sécantes en $B$} du plan $(ABC)$. Donc $(CD)$ est orthogonale au plan $(ABC)$.
\item Le plan $(ACD)$ contient la droite $(CD)$, qui est orthogonale au plan $(ABC)$ (question 3). Or deux plans sont orthogonaux lorsque l'un contient une droite orthogonale à l'autre. Donc $(ACD) \perp (ABC)$.
\end{enumerate}
\end{corrige}

\begin{corrige}[Exercice 7 — Pyramide à base carrée]
\begin{enumerate}
\item $(SO) \perp (ABCD)$ par hypothèse ; or $(AC) \subset (ABCD)$ et $(BD) \subset (ABCD)$, donc $(SO) \perp (AC)$ et $(SO) \perp (BD)$.
\item La diagonale du carré de côté $6$ cm vaut $AC = 6\sqrt{2}$ cm, donc $OA = \dfrac{AC}{2} = 3\sqrt{2}$ cm. Le triangle $SOA$ est rectangle en $O$ (car $(SO) \perp (OA)$) ; par le théorème de Pythagore :
\[ SA^2 = SO^2 + OA^2 = 16 + 18 = 34 \quad \text{donc} \quad SA = \sqrt{34} \approx 5{,}83 \text{ cm}. \]
\item $(BD) \perp (SO)$ (question 1) et $(BD) \perp (AC)$ (diagonales du carré $ABCD$). Or $(SO)$ et $(AC)$ sont deux droites sécantes en $O$ du plan $(SAC)$. Donc $(BD) \perp (SAC)$.
\item Le plan $(SBD)$ contient la droite $(BD)$, orthogonale au plan $(SAC)$. Donc $(SBD) \perp (SAC)$.
\end{enumerate}
\end{corrige}

\begin{corrige}[Exercice 8 — Raisonnement par l'absurde]
\begin{enumerate}
\item Si un plan $(P)$ contient $(AC)$ et $(BH)$, alors il contient les points $A$ et $C$ (points de $(AC)$) ainsi que $B$ et $H$ (points de $(BH)$), donc les quatre points $A$, $B$, $C$, $H$.
\item $(P)$ contiendrait les points $A$, $B$, $C$ non alignés : or trois points non alignés déterminent un unique plan, donc $(P) = (ABC) = (ABCD)$ (plan de la face inférieure). Mais $H$ est un sommet de la face supérieure, $H \notin (ABCD)$ : contradiction.
\item L'hypothèse est donc fausse : il n'existe aucun plan contenant $(AC)$ et $(BH)$. Les droites $(AC)$ et $(BH)$ sont \textbf{non coplanaires}.
\item Oui, elles sont orthogonales. D'une part $(AC) \perp (BD)$ (diagonales du carré $ABCD$). D'autre part $(DH) \perp (ABCD)$ (arête verticale du cube), donc $(DH) \perp (AC)$. Ainsi $(AC)$ est orthogonale aux deux droites sécantes $(BD)$ et $(DH)$ du plan $(BDH)$, donc $(AC) \perp (BDH)$. Or $(BH) \subset (BDH)$, donc $(AC) \perp (BH)$ : elles sont orthogonales (sans être sécantes).
\end{enumerate}
\end{corrige}

\begin{corrige}[Exercice 9 — Pavé droit : la diagonale orthogonale à un plan]
\textbf{Partie A.}
\begin{enumerate}
\item Le pavé est droit : l'arête $(AE)$ est orthogonale à la face $(ABCD)$, donc aussi au plan parallèle $(EFGH)$ (une droite orthogonale à un plan est orthogonale à tout plan parallèle ; ici $(AE) \perp (EF)$ et $(AE) \perp (EH)$, deux droites sécantes de $(EFGH)$). Comme $(HF) \subset (EFGH)$, on en déduit $(AE) \perp (HF)$.
\item Comme $a = b$, la face $EFGH$ est un \emph{carré} ; ses diagonales $(HF)$ et $(EG)$ sont donc perpendiculaires : $(HF) \perp (EG)$.
\end{enumerate}
\textbf{Partie B.}
\begin{enumerate}
\setcounter{enumi}{2}
\item $(HF)$ est orthogonale à $(AE)$ (question 1) et à $(EG)$ (question 2). Or $(AE)$ et $(EG)$ sont deux droites sécantes en $E$ du plan $(AEG)$. Donc $(HF) \perp (AEG)$.
\item Le plan $(EFGH)$ contient la droite $(HF)$, orthogonale au plan $(AEG)$. Donc $(EFGH) \perp (AEG)$.
\end{enumerate}
\textbf{Partie C.}
\begin{enumerate}
\setcounter{enumi}{4}
\item $EG$ est la diagonale du carré $EFGH$ de côté $3$ m : $EG = 3\sqrt{2}$ m. Pour $AG$ : dans le triangle rectangle $AEG$ (rectangle en $E$ car $(AE) \perp (EFGH)$) :
\[ AG^2 = AE^2 + EG^2 = 16 + 18 = 34 \quad \text{donc} \quad AG = \sqrt{34} \approx 5{,}83 \text{ m}. \]
\item $I$ est le centre du carré $EFGH$ (intersection des diagonales), donc $EI = \dfrac{EG}{2} = \dfrac{3\sqrt{2}}{2}$ m. Comme $(AE) \perp (EFGH)$, le triangle $AEI$ est rectangle en $E$ :
\[ AI^2 = AE^2 + EI^2 = 16 + \frac{18}{4} = 16 + 4{,}5 = 20{,}5 \quad \text{donc} \quad AI = \sqrt{20{,}5} = \frac{\sqrt{82}}{2} \approx 4{,}53 \text{ m}. \]
\end{enumerate}
\textbf{Barème indicatif (sur 10)} : Q1 : 1,5 pt ; Q2 : 1 pt ; Q3 : 2 pts ; Q4 : 1 pt ; Q5 : 2,5 pts ; Q6 : 2 pts.
\begin{attention}[Points de vigilance]
Citer explicitement que $(AE)$ et $(EG)$ sont \emph{sécantes en $E$} dans le plan $(AEG)$ : c'est la condition d'application du théorème. Pour la question 4, préciser que $(HF) \subset (EFGH)$. Dans les calculs, justifier le caractère rectangle des triangles avant d'appliquer Pythagore.
\end{attention}
\end{corrige}

\begin{corrige}[Exercice 10 — Le poteau et ses haubans]
\textbf{Partie A.}
\begin{enumerate}
\item $(SO)$ est perpendiculaire en $O$ aux droites $(OA)$ et $(OB)$. Comme $O$, $A$, $B$ ne sont pas alignés, $(OA)$ et $(OB)$ sont deux droites \emph{sécantes en $O$} du plan $(P)$. Donc $(SO) \perp (P)$.
\item $(SO) \perp (P)$ implique $(SO)$ orthogonale à toute droite de $(P)$ passant par $O$, en particulier $(SO) \perp (OA)$ et $(SO) \perp (OB)$ : les triangles $SOA$ et $SOB$ sont rectangles en $O$.
\item Le plan $(SOA)$ contient la droite $(SO)$, orthogonale au plan $(P)$. Donc $(SOA) \perp (P)$.
\end{enumerate}
\textbf{Partie B.}
\begin{enumerate}
\setcounter{enumi}{3}
\item Pythagore dans $SOA$ rectangle en $O$ : $SA^2 = 6^2 + 2{,}5^2 = 36 + 6{,}25 = 42{,}25$, donc $SA = 6{,}5$ m (valeur exacte, car $6{,}5^2 = 42{,}25$). Pythagore dans $SOB$ : $SB^2 = 36 + 16 = 52$, donc $SB = \sqrt{52} = 2\sqrt{13} \approx 7{,}21$ m.
\item Longueur totale : $SA + SB = 6{,}5 + 2\sqrt{13} \approx 6{,}5 + 7{,}21 = 13{,}71$ m $< 15$ m. Oui, un rouleau de $15$ m suffit (il reste environ $1{,}29$ m).
\item Si $\widehat{AOB} = 90^{\circ}$, le triangle $AOB$ est rectangle en $O$ :
\[ AB^2 = OA^2 + OB^2 = 6{,}25 + 16 = 22{,}25 \quad \text{donc} \quad AB = \sqrt{22{,}25} \approx 4{,}72 \text{ m}. \]
\end{enumerate}
\textbf{Barème indicatif (sur 10)} : Q1 : 2 pts ; Q2 : 1 pt ; Q3 : 1,5 pt ; Q4 : 2,5 pts ; Q5 : 1,5 pt ; Q6 : 1,5 pt.
\begin{attention}[Points de vigilance]
La non-alignement de $O$, $A$, $B$ est la condition qui garantit que $(OA)$ et $(OB)$ sont sécantes : il faut le citer. Pour la question 5, comparer la longueur totale à $15$ m avec un calcul explicite et une conclusion rédigée.
\end{attention}
\end{corrige}

\begin{corrige}[Exercice 11 — Tétraèdre régulier et orthogonalité]
\textbf{Partie A.}
\begin{enumerate}
\item Toutes les arêtes valent $a$, donc $AC = CD = AD = a$ et $BC = CD = BD = a$ : les triangles $ACD$ et $BCD$ sont équilatéraux. Dans un triangle équilatéral, la médiane issue d'un sommet est aussi hauteur : $I$ milieu de $[CD]$ donne $(AI) \perp (CD)$ dans $ACD$ et $(BI) \perp (CD)$ dans $BCD$.
\item $(CD)$ est perpendiculaire à $(AI)$ et à $(BI)$, deux droites sécantes en $I$ du plan $(ABI)$. Donc $(CD) \perp (ABI)$.
\item $(AB) \subset (ABI)$, donc $(CD) \perp (AB)$. Conclusion : dans un tétraèdre régulier, deux arêtes opposées quelconques sont orthogonales (non sécantes).
\item Le plan $(BCD)$ contient la droite $(CD)$, orthogonale au plan $(ABI)$. Donc $(BCD) \perp (ABI)$.
\end{enumerate}
\textbf{Partie B.}
\begin{enumerate}
\setcounter{enumi}{4}
\item $AI$ est la hauteur d'un triangle équilatéral de côté $a = 6$ cm : $AI = \dfrac{a\sqrt{3}}{2} = 3\sqrt{3}$ cm. De même $BI = 3\sqrt{3}$ cm. (Vérification par Pythagore : $AI^2 = AC^2 - CI^2 = 36 - 9 = 27$, $AI = \sqrt{27} = 3\sqrt{3}$.)
\item Le triangle $ABI$ a $AI = BI = 3\sqrt{3}$ et $AB = 6$. Soit $J$ le milieu de $[AB]$ : la hauteur $(IJ)$ vérifie
\[ IJ^2 = AI^2 - AJ^2 = 27 - 9 = 18 \quad \text{donc} \quad IJ = 3\sqrt{2} \text{ cm}. \]
Aire : $\mathcal{A} = \dfrac{AB \times IJ}{2} = \dfrac{6 \times 3\sqrt{2}}{2} = 9\sqrt{2} \approx 12{,}73 \text{ cm}^2$.
\item $H$ est le centre de gravité du triangle équilatéral $BCD$ ; il est situé aux deux tiers de la médiane $[BI]$ à partir de $B$ :
\[ BH = \frac{2}{3} BI = \frac{2}{3} \times 3\sqrt{3} = 2\sqrt{3} \approx 3{,}46 \text{ cm}. \]
\end{enumerate}
\textbf{Barème indicatif (sur 10)} : Q1 : 2 pts ; Q2 : 1,5 pt ; Q3 : 1,5 pt ; Q4 : 1 pt ; Q5 : 1,5 pt ; Q6 : 1,5 pt ; Q7 : 1 pt.
\begin{attention}[Points de vigilance]
Bien distinguer « perpendiculaires » (droites sécantes, dans un même plan) et « orthogonales » (sens général de l'espace). Pour la question 2, citer la sécance de $(AI)$ et $(BI)$ en $I$. Pour la question 7, le centre de gravité est aux $\frac{2}{3}$ de la médiane \emph{à partir du sommet} : c'est bien $BH = \frac{2}{3}BI$, et non $\frac{1}{3}$.
\end{attention}
\end{corrige}

\newpage

\coursec{Fiche bilan}

\begin{popfigure}
\begin{tikzpicture}[mindmap, grow cyclic, every node/.style={concept, text=white, font=\small\bfseries},
  root concept/.append style={concept color=popdark, font=\bfseries},
  level 1/.append style={level distance=3.4cm, sibling angle=60},
  level 2/.append style={level distance=2.1cm, sibling angle=45, font=\scriptsize}]
\node[root concept]{Orthogonalité dans l'espace}
  child[concept color=popblue]{ node {Positions relatives}
    child { node {coplanaires ?} }
    child { node {sécantes, parallèles, non coplanaires} }
  }
  child[concept color=poporange]{ node {Droites orthogonales}
    child { node {parallèles à deux perpendiculaires} }
    child { node {perp. $\Rightarrow$ orth.} }
  }
  child[concept color=popgreen]{ node {Droite $\perp$ plan}
    child { node {orth. à 2 droites sécantes du plan} }
    child { node {orth. à toute droite du plan} }
  }
  child[concept color=poppurple]{ node {Orthogonalité et parallélisme}
    child { node {deux droites $\perp$ même plan $\Rightarrow$ parallèles} }
    child { node {deux plans $\perp$ même droite $\Rightarrow$ parallèles} }
  }
  child[concept color=poppink]{ node {Plans orthogonaux}
    child { node {l'un contient une droite $\perp$ à l'autre} }
  }
  child[concept color=popgold]{ node {Méthodes}
    child { node {réciprocité, rédaction} }
  };
\end{tikzpicture}
\legende{Carte mentale de la leçon : l'orthogonalité dans l'espace en un coup d'\oe il.}
\end{popfigure}

\begin{bilan}
\textbf{Définitions clés.}
\begin{itemize}
  \item Deux droites sont \cle{coplanaires} si elles appartiennent à un même plan ; sinon elles sont \cle{non coplanaires}.
  \item Deux droites sont \cle{orthogonales} si les parallèles menées par un point quelconque à chacune d'elles sont perpendiculaires. Notation : $(d) \perp (\Delta)$.
  \item Deux droites \cle{perpendiculaires} sont orthogonales \emph{et sécantes}.
  \item Une droite $(d)$ est \cle{orthogonale à un plan} $(P)$ si elle est orthogonale à \textbf{toute} droite de $(P)$. Notation : $(d) \perp (P)$.
  \item Deux plans sont \cle{orthogonaux} (ou perpendiculaires) si l'un contient une droite orthogonale à l'autre. Notation : $(P) \perp (Q)$.
\end{itemize}

\textbf{Théorèmes et propriétés à connaître par c\oe ur.}
\begin{center}
\begin{tabularx}{\linewidth}{|l|X|}
\hline
\rowcolor{popblueL} \textbf{Théorème} & \textbf{Énoncé} \\
\hline
Critère droite $\perp$ plan & $(d)$ orthogonale à \textbf{deux droites sécantes} de $(P)$ $\Rightarrow$ $(d) \perp (P)$. \\
\hline
Conséquence & Si $(d) \perp (P)$, alors $(d)$ est orthogonale à toute droite de $(P)$. \\
\hline
Unicité & Par un point donné passe une seule droite orthogonale à un plan, et un seul plan orthogonal à une droite. \\
\hline
Liaison 1 & Deux droites orthogonales à un même plan sont parallèles entre elles. \\
\hline
Liaison 2 & Deux plans orthogonaux à une même droite sont parallèles entre eux. \\
\hline
Liaison 3 & Si $(d) \perp (P)$ et $(d) \parallel (\Delta)$, alors $(\Delta) \perp (P)$. \\
\hline
Plans orthogonaux & $(P) \perp (Q)$ $\iff$ $(P)$ contient une droite orthogonale à $(Q)$. \\
\hline
Intersection & Si $(P) \perp (Q)$, toute droite de $(P)$ orthogonale à $(P)\cap(Q)$ est orthogonale à $(Q)$. \\
\hline
\end{tabularx}
\end{center}

\textbf{Méthodes express.}
\begin{itemize}
  \item \textbf{Montrer que deux droites sont orthogonales} : trouver un plan contenant l'une et orthogonal à l'autre, ou exhiber deux perpendiculaires respectivement parallèles.
  \item \textbf{Montrer que $(d) \perp (P)$} : exhiber deux droites \emph{sécantes} de $(P)$ orthogonales à $(d)$ (souvent en utilisant des triangles isocèles, des diagonales de carrés, des médianes).
  \item \textbf{Montrer que $(P) \perp (Q)$} : trouver dans $(P)$ une droite orthogonale à $(Q)$.
  \item \textbf{Montrer que deux droites sont non coplanaires} : raisonner par l'absurde ou montrer qu'elles ne sont ni sécantes ni parallèles.
  \item Dans un cube ou un pavé : les arêtes issues d'un sommet sont deux à deux orthogonales ; une arête est orthogonale aux faces qui ne la contiennent pas.
\end{itemize}
\end{bilan}

\begin{aretenir}[Checklist avant le Probatoire]
\begin{itemize}
  \item[$\square$] Je sais citer les trois positions relatives possibles de deux droites de l'espace.
  \item[$\square$] Je sais distinguer \emph{orthogonales} (pas forcément sécantes) et \emph{perpendiculaires} (sécantes).
  \item[$\square$] Je sais que perpendiculaires $\Rightarrow$ orthogonales, mais pas la réciproque.
  \item[$\square$] Je connais le critère fondamental : $(d) \perp (P)$ si $(d)$ est orthogonale à \textbf{deux droites sécantes} de $(P)$.
  \item[$\square$] Je n'oublie jamais de vérifier (et d'écrire) que les deux droites du plan sont bien \textbf{sécantes}.
  \item[$\square$] Je sais que $(d) \perp (P)$ entraîne $(d)$ orthogonale à \textbf{toute} droite de $(P)$.
  \item[$\square$] Je connais les théorèmes de liaison entre orthogonalité et parallélisme (droites $\perp$ même plan ; plans $\perp$ même droite).
  \item[$\square$] Je sais montrer que deux plans sont orthogonaux en exhibant une droite de l'un orthogonale à l'autre.
  \item[$\square$] Je sais exploiter les propriétés du cube et du pavé droit (arêtes, faces, diagonales).
  \item[$\square$] Je sais prouver que deux droites sont non coplanaires (ni sécantes, ni parallèles).
  \item[$\square$] Je rédige mes démonstrations en citant précisément les hypothèses et le théorème utilisé.
  \item[$\square$] Je fais une figure claire, en perspective cavalière, avant de raisonner.
\end{itemize}
\end{aretenir}

\findecours

\end{document}
